\documentclass[ecta,nameyear,draft]{econsocart}
\usepackage{amsmath,amsfonts,amssymb,amsthm,booktabs,array,longtable}
\usepackage[T1]{fontenc}
\usepackage{mathpazo,microtype,needspace,xurl}
\RequirePackage[colorlinks,citecolor=blue,linkcolor=blue,urlcolor=blue]{hyperref}
\startlocaldefs
\newtheorem{theorem}{Theorem}
\newtheorem{proposition}[theorem]{Proposition}
\newtheorem{lemma}[theorem]{Lemma}
\newtheorem{corollary}[theorem]{Corollary}
\newcommand{\E}{\mathbb E}
\newcommand{\R}{\mathbb R}
\endlocaldefs
\setlength{\emergencystretch}{3em}
\makeatletter\def\copyright@text{Prepared for referee review}\makeatother

\usepackage{xr-hyper}
\externaldocument[M-]{build/complete}[ECTA.pdf]
\begin{document}
\begin{frontmatter}
\title{Supplement to ``Neural Bellman Operators''}
\runtitle{Neural Bellman Operators: Supplement}
\begin{aug}
\author[id=au1,addressref={add1}]{\fnms{Qian}~\snm{QI}}
\address[id=add1]{Peking University}
\end{aug}
\begin{abstract}
This supplement supplies proofs and implementation details for the R48 revision, retaining the R44--R47 results: centered economic transport, direct full-policy certification, retained-policy recertification, and coupled policy-cost comparisons. It reports all new comparison intervals and separates mathematical, numerical, and statistical guarantees. The constructive neural backend has its own complete primitive-budget proof below. Complete earlier proofs and economic applications are materialized unchanged in the current source tree.
\end{abstract}
\end{frontmatter}
\section*{Reading notice for the complete development edition}
This volume retains the baseline development verbatim, including historical descriptions, followed by the current R52 argument. Its historical reference to a formerly authoritative revision is not a second current scientific claim. The active article, technical supplement and response govern the new evidence and numerical corrections. Every baseline result and its original assumptions remain available for examination.

\section{Notation and economic domains}
Dates are $0,\ldots,T$; $B_0=1$ and $B_t=\prod_{j<t}\beta_j$ with $\beta_j>0$. The state domain is invariant under all admissible actions and the declared innovation law. The fixed-policy operator is
$\mathcal T_t^\pi v(x)=c_t(x,\pi_t(x))+\beta_t\rho_t(v(F_t(x,\pi_t(x),Z)))$.
The Bellman operator is its infimum over feasible current actions. Every cash-invariance and monotonicity statement is restricted to a common domain containing the values, fitted functions, and constant translates used in the comparison. This requirement is substantive for recursive utility.

For compact continuous expected-cost problems with continuous feasible correspondences, continuous integrals, and measurable minimizers, backward induction produces the Bellman optimum. An adapted policy cannot lower its conditional continuation cost below the Bellman value at any date. Induction therefore proves optimality against history-dependent admissible controls, rather than only Markov grid policies. The same reasoning applies to the stipulated monotone recursive evaluation when its conditional recursion and measurable selection conditions hold. The Gaussian entropic construction in the original article uses separate moment margins rather than compactness.

\section{Centered decision and transport proofs}
Write $H(a)=r_\theta(s,a)+S_\theta(T_\theta(s,a))$ and
$\widehat H(a)=r_\theta(s,a)+\widehat S_0(T_\theta(s,a))$.
Then $\widehat H-H=e\circ T_\theta-\Delta_\theta\circ T_\theta$.
For an arbitrarily near-optimal true action $a^*$,
\[
 H(a^*)-H(\widehat a)
 \leq\delta+(H-\widehat H)(a^*)-(H-\widehat H)(\widehat a).
\]
Bounding the two oscillations and letting the approximation to $a^*$ vanish proves the centered bound, even if a supremum is not attained. For the finite menu let $\bar e=\sum_jp_je_j$. For distinct $i,j$, weighted Cauchy--Schwarz gives
\[
 |e_i-e_j|^2
 =\left|\sum_kp_k(e_k-\bar e)
       \left(\frac{\mathbf1_{k=i}}{p_i}-\frac{\mathbf1_{k=j}}{p_j}\right)\right|^2
 \leq R_p(p_i^{-1}+p_j^{-1}).
\]
Maximizing over $i,j$ and bounding the continuation perturbation by $2b_\theta$ proves the menu statement. It is a statement about a common menu; a method-specific pair of actions cannot substitute for the full set.

For transport, couple the future economies with the same innovations, starting from the same postdecision state. The assumed transition bounds imply pathwise
$\|X_{k+1}^\theta-X_{k+1}^0\|\leq L_k\|X_k^\theta-X_k^0\|+\epsilon_{\theta,k}$.
Thus $\|X_k^\theta-X_k^0\|\leq d_k$ by induction. Each future reward differs by at most $\eta_{\theta,k}+H_kd_k$. Summation with nonnegative weights and expectation prove the transport bound. If the recursion is nonlinear but sup-norm Lipschitz on a common utility domain, add and subtract the anchor map at the changed continuation to obtain $b_k\leq\nu_k+\beta_kb_{k+1}$; unwind this inequality backward.

These proofs depend only on the declared economic perturbation class. They do not assert that all policy changes, altered constraints, or opponent deviations preserve the same domain or constants.

\section{Signed full-policy certificate and retained actors}
Suppose $\ell_t\leq f_t-\mathcal T_tf_{t+1}\leq u_t$ and
$\mathcal T_t^\pi f_{t+1}-\mathcal T_tf_{t+1}\leq\eta_t$ for $t<T$, with terminal enclosure $\ell_T\leq f_T-g\leq u_T$.
For each date let $B_{t,k}=\prod_{j=t}^{k-1}\beta_j$ and set
\[
 a_t=\sum_{k=t}^{T-1}B_{t,k}u_k+B_{t,T}u_T,
 \qquad
 b_t=\sum_{k=t}^{T-1}B_{t,k}(\eta_k-\ell_k)-B_{t,T}\ell_T.
\]
We prove $V_t^*\geq f_t-a_t$ and $J_t^\pi\leq f_t+b_t$ simultaneously. The terminal claim is the terminal enclosure. For the lower induction,
\[
 V_t^*=\mathcal T_tV_{t+1}^*
 \geq\mathcal T_t(f_{t+1}-a_{t+1})
 \geq f_t-u_t-\beta_ta_{t+1}=f_t-a_t.
\]
For the upper induction,
\[
 J_t^\pi=\mathcal T_t^\pi J_{t+1}^\pi
 \leq\mathcal T_t^\pi(f_{t+1}+b_{t+1})
 \leq f_t-\ell_t+\eta_t+\beta_tb_{t+1}=f_t+b_t.
\]
Subtraction yields the theorem. Additive shifts $d_t$ change each residual by $d_t-\beta_td_{t+1}$ and the terminal residual by $d_T$; widths and action gaps are unchanged. This proves the invariance without assuming a contraction factor less than one.

For retained-policy recertification, new residual enclosures are substituted in exactly the same induction. The old actor allowance remains valid because its hypothesis is a uniform statement about the original actor and the unchanged functions. It has no dependence on the mesh on which the new residual happens to be validated. The terminal allowance can also be retained. Intersecting two valid residual intervals preserves validity and cannot increase width. In the executed diagnostics the new residuals are used directly, with the old actor and terminal enclosures; the record contains every component needed to reconstruct the bound.

The zero-shift case does not imply that constant residuals disappear from all value brackets. They disappear from the policy-loss bound through upper-minus-lower cancellation. Keeping signed endpoints is consequently important for direct cost scores as well as for sharp policy bounds.

\section{Verification geometry and numerical enclosures}
\subsection{The scalar primitive economy}
The scalar state is $x\in[0,1]$, action $a\in[0,1/4]$, horizon four, and discount $15/16$. The transition is
\[
 F(x,a,z)=\operatorname{clip}_{[0,1]}
       \{(13/16)x+a+(1/16)x(1-x)+z\},
\]
with $z\in\{-1/16,0,1/16\}$ and probabilities $(1/4,1/2,1/4)$. The cost is
\[
 c(x,a)=(x-11/16)^2+pa^2+4a^4+2(3/8-x)_+^2,
\]
and terminal cost is $g(x)=2(x-11/16)^2+2(3/8-x)_+^2$. The two evaluations are expectation and the entropic certainty equivalent $\rho_\theta(v)=\theta^{-1}\log\E e^{\theta v}$ with $\theta=1$; expectation is $\theta=0$. The discrete law here is an economic primitive, unlike the continuous uniform law of the coupled experiment.

For a scalar ReLU network, exact dyadic coefficients permit rational breakpoints and slopes. The largest absolute slope bounds its Lipschitz constant $L_t$. The Bellman state and action constants used by the verifier are
\[
 L_t^x=2.875+(15/16)(0.875)L_{t+1},\qquad
 L_t^a=p/2+1/4+(15/16)L_{t+1}.
\]
For $N+1$ state nodes and $A+1$ action nodes, the action cover allowance is $L_t^a/(8A)$ and the residual state allowance is $(L_t+L_t^x)/(2N)$. Each residual endpoint includes the state allowance; the nodal Bellman lower bound includes the action cover. The nearest-node actor allowance adds the selected nodal action gap, action cover, and $L_t^x/N$. All operations are outward enclosed.

The terminal residual is piecewise quadratic. Its extrema occur at exact network breakpoints, the cost kink, endpoints, or stationary points in each piece. The implementation enumerates these candidates using rational arithmetic and rounds the extrema outward. It therefore does not replace an exact terminal calculation with an uncharged sampled check.

\subsection{The coupled action and state certificates}
For fixed $x$, the coupled action objective $q(a)$ is $\mu$-strongly convex on the simplex, with $\mu=2p-1/2$. For any feasible proposal $a$ and nonnegative multiplier $\nu$, the global quadratic lower model yields
\begin{align*}
 \inf_{b\geq0,\,\mathbf1'b\leq h}q(b)
 &\geq q(a)-\nabla q(a)'a+\tfrac\mu2\|a\|^2-\nu h\\
 &\quad-\frac1{2\mu}\sum_j
          \bigl(\min\{\nabla_jq(a)-\mu a_j+\nu,0\}\bigr)^2,
 \qquad h=1/8.
\end{align*}
Indeed, strong convexity gives a quadratic lower model; adding $\nu(\mathbf1'b-h)$ gives a lower bound on the constrained infimum, and minimizing the separable quadratic over $b\geq0$ gives the displayed expression. Any nonnegative multiplier is valid. Interval lower evaluation and the feasible proposal's upper value enclose the continuous-action infimum without trusting numerical first-order convergence.

Let $\delta$ bound the nodal primal--dual gap, and let $L_{xa}$ bound the change of the action gradient with state. A node action with gap $\delta$ is within $\sqrt{2\delta/\mu}$ of the node optimizer. At a state distance $r$, the gradient perturbation is at most $L_{xa}r$. The loss of retaining the node action is bounded by
\[
 \eta=\delta+(L_{xa}r)\sqrt{2\delta/\mu}
                 +\frac{(L_{xa}r)^2}{2\mu}.
\]
This follows by combining strong convexity at the node with the cross-state change of objective differences and maximizing the resulting linear-minus-quadratic expression in the distance to the node optimizer. On the square tensor grid, $r=\sqrt2/(2N)$. Rounded acquisition adds its own state-distance allowance before applying the formula. The ideal policy and rounded-state policy bounds are separately named in the original records.

If the gradient of the fitted current continuation has Lipschitz constant $M_t$ and the Bellman envelope has gradient Lipschitz constant $H_{t+1}$, bilinear interpolation on cells of side $1/N$ has residual error at most $(M_t+H_{t+1})/(4N^2)$. To see this, condition successively on a coordinate and use the one-dimensional linear-interpolation error $Lh^2/8$ for a function with Lipschitz derivative. Summing two coordinate errors gives $Lh^2/4$. The source-bound derivative envelopes, rather than empirical finite differences, supply these constants. The terminal allowance uses $(M_T+18)/(4N^2)$.

In the retained-policy diagnostic, these residual terms are refined at $N=256$, but the actor term remains the one certified for $N=128$. The record may also contain the fine-grid proposed actor's allowance as part of the inherited verifier output. That allowance is not used to claim an unchanged-policy result.

\subsection{Arithmetic assumptions and implementation boundary}
The inherited interval implementation applies outward rounding to elementary binary operations and encloses its exponential and logarithm approximations by explicit series remainders. Stored network and actor coefficients are interpreted as exact dyadics. Analytic innovation moments are enclosed using interval arithmetic, including cancellation. Exact rational breakpoint calculations and independently checked square roots are used where specified.

The resulting certificates concern the defined policies, with numerical enclosures around their evaluation. The original rounded-state and execution accounts remain separate. The new retained-policy bound does not silently assert that an arbitrary hardware implementation with different state acquisition, optimizer, or clipping conventions inherits the same target. A deployed implementation must satisfy the execution assumptions or add its own allowance.

\section{Paired policy costs: probability and numerical error}
\subsection{An unbiased score with certificate-derived range}
Under expectation and a fixed policy, the cost score is
\[
 Z_m=f_0^m(x_0)+\sum_{t<T}B_t(Q_t(f_{t+1}^m;X_t^m,\pi_t^m)-f_t^m(X_t^m))
            +B_T(g(X_T^m)-f_T^m(X_T^m)).
\]
Taking expectations cancels the fitted values at every date, including both endpoints. The score has expected value equal to actual policy cost even if the networks are poor approximations. To obtain bounded support, write the centered summand as the nonnegative selected-action gap minus the signed Bellman residual. It lies in $[-u_t,\eta_t-\ell_t]$. The terminal term lies in $[-u_T,-\ell_T]$. Thus the score interval has width precisely the corresponding ideal policy bound. Subtracting two scores has a support width equal to the sum of their policy bounds, but its expectation is the direct difference of policy costs. The distinction is that the scores are evaluated along the policies' trajectories; two regret upper bounds are never used as a surrogate observation of the cost difference.

\subsection{Continuous innovations through uniform bins}
For each comparison, draw independent integer arrays $K_{itj}$ uniformly in $\{0,\ldots,2^{40}-1\}$. A bin represents the continuous interval
\[
 \left[\frac1{32}\left(\frac{2K_{itj}}{2^{40}}-1\right),
       \frac1{32}\left(\frac{2(K_{itj}+1)}{2^{40}}-1\right)\right].
\]
The exact uniform innovation law can be represented by a uniform bin followed by a uniform draw inside it. The implementation encloses every draw in that bin and needs no additional within-bin random number. Both policies use the same innovation bins but propagate their own states and choose their own actions.

At each date and in each coordinate, the actor is the half-up nearest-node rule. If a state interval meets more than one cell, the implementation includes all eligible node actions and propagates their componentwise hull. Such a hull need not itself be a feasible action, but interval evaluation over the hull still encloses every feasible action actually selected; it is used for enclosure, not as a new policy. The bounding arithmetic covers network values, analytic expectations, transition, stage cost, and terminal cost. Intersecting sample score enclosures with the deterministic support is valid for every bin.

For a bin $K_i$, the computed endpoints $A_i,D_i$ are deterministic functions with $A_i\leq Z_i\leq D_i$ for every within-bin realization. Hence $\E A_i\leq\E Z_i\leq\E D_i$. Independence of bin draws makes each endpoint sequence independent and identically distributed, conditional on the frozen policies. Endpoint widening changes numerical precision, not the underlying innovation law.

\subsection{Simultaneous coverage}
Let $q=24$, $\alpha=0.01$, and $n=65536$. Scale each endpoint sequence to its deterministic support of width $B$. Apply the one-sided empirical Bernstein inequality of \citet{maurer2009}, Theorem 4, to the lower sequence with reversed sign and to the upper sequence in its usual direction, each at tail probability $\alpha/(2q)$. Rescaling gives the radius in the main article with logarithm $\log(4q/\alpha)=\log(9600)$. The union bound over all 48 tails yields simultaneous coverage at least $0.99$. Dependence between comparisons is not excluded or required to vanish; only within-comparison trajectory independence enters this argument.

The code computes means by outward pairwise summation. If $[\underline m,\overline m]$ contains the true sample mean, interval squares of $x_i-[\underline m,\overline m]$ bound every squared deviation from that mean. Their outward sum divided by $n-1$ therefore bounds sample variance above. The logarithm uses the inherited interval series. A proposed floating square root is checked by exact rational squaring and enlarged until it is an upper bound. Since the radius is increasing in variance, range, and logarithm, these substitutions preserve coverage.

This is a probability theorem for the stipulated independent simulation model. The fixed PCG64 seed identifies the reproducible realization; it is not a mathematical proof that a deterministic sequence is independent. The comparison states, policies, sample count, and multiplicity are fixed before the new simulations. No optional stopping or policy selection is performed. All 24 intervals contain zero, so no sign or equivalence conclusion is made.

\begin{center}\small
\begin{longtable}{rrrrr}
\caption{All direct cost intervals. Endpoints are multiplied by $10^4$. States are ordered as in the main article.}\label{tab:pairedfull}\\
\toprule Price & Training seed & State & Lower & Upper\\\midrule\endfirsthead
\toprule Price & Training seed & State & Lower & Upper\\\midrule\endhead
1 & 104729 & 1 & -1.302334 & 1.319566\\
1 & 104729 & 2 & -1.343033 & 1.331605\\
1 & 104729 & 3 & -1.272277 & 1.395933\\
1 & 104729 & 4 & -1.318164 & 1.380020\\
1 & 130363 & 1 & -0.913721 & 1.096187\\
1 & 130363 & 2 & -0.930282 & 1.016606\\
1 & 130363 & 3 & -0.877286 & 1.070411\\
1 & 130363 & 4 & -0.930738 & 1.031036\\
1 & 155921 & 1 & -1.119909 & 1.057248\\
1 & 155921 & 2 & -1.155572 & 1.015951\\
1 & 155921 & 3 & -1.119496 & 1.035950\\
1 & 155921 & 4 & -1.182758 & 1.039751\\
4 & 104729 & 1 & -1.110817 & 0.892852\\
4 & 104729 & 2 & -1.089372 & 0.876343\\
4 & 104729 & 3 & -1.101773 & 0.870680\\
4 & 104729 & 4 & -1.301186 & 0.663547\\
4 & 130363 & 1 & -1.217968 & 1.253187\\
4 & 130363 & 2 & -1.218284 & 1.248764\\
4 & 130363 & 3 & -1.219101 & 1.252803\\
4 & 130363 & 4 & -1.227610 & 1.207526\\
4 & 155921 & 1 & -1.192728 & 1.201462\\
4 & 155921 & 2 & -1.190404 & 1.185857\\
4 & 155921 & 3 & -1.176320 & 1.205789\\
4 & 155921 & 4 & -1.171505 & 1.216969\\
\bottomrule\end{longtable}\end{center}


\section{Work accounts, component diagnostics, and reproduction}
\begin{table}[htbp]
\centering\small
\caption{Additional verification work and local diagnostic clocks}
\label{tab:diagwork}
\begin{tabular}{lrrrr}
\toprule
Service & Mesh & Operations & Retained actor & Seconds\\
\midrule
041 & 4096 & 50,393,100 & 8,196 & 11.920\\
089 & 4096 & 50,393,100 & 8,196 & 67.576\\
120 & 4096 & 50,393,100 & 8,196 & 59.052\\
124 & 4096 & 50,393,100 & 8,196 & 14.998\\
128 & 256 & 2,054,388,096 & 133,128 & 140.429\\
145 & 4096 & 50,393,100 & 8,196 & 87.250\\
147 & 4096 & 50,393,100 & 8,196 & 15.473\\
\bottomrule
\end{tabular}
\end{table}
Operations count Bellman transitions for scalar services and neuron expectations for the coupled service; these are different units, not cross-method speed estimates.

The additional verification clocks include computations of fine-grid action proposals even though those proposals are not deployed. They do not repeat training or reconstruct the original service clock. Jobs were scheduled concurrently on a resource-limited local host; these clocks describe those diagnostics and cannot be read as isolated cross-method speed measurements.

The original records expose state coverage, action coverage, nodal gaps, and terminal bounds. A representation/optimizer decomposition is not identifiable from them: different parameter errors and fitting trajectories can produce the same signed Bellman residual. We report the available components as an error account, not a causal identification of each source of approximation error. Training residuals are available as separate diagnostics and are not added twice to a certificate that already evaluates the trained functions.

A fresh reproduction obtains the two pinned input artifacts, verifies their archive digests, and uses a clean result directory. The study commands deliberately refuse to overwrite completed records. The deterministic audit reconstructs target counts and signed comparisons from all original records, checks retained network and actor hashes, and reconstructs every new confidence endpoint from stored moments and support. Regression tests exercise interval arithmetic, exact finite-state policy inequalities, offset invariance, and source identities. The release manifest binds every active source, table, result, and compiled document.

\section{Relationship to the retained theory}
The full controlled-diffusion and viscosity analysis, economic boundary and transversality conditions, recursive utility and endogenous-preference models, temporal selves, games, quadratic hidden-factor construction, numerical acceptance theorems, and historical experiments remain in their unchanged source files and compiled companions. The new main article does not purport to re-prove them by citing the nonlinear experiment. The preservation map records file hashes and LaTeX labels, and the response gives the location of each current answer to the referee.

The current revision supplies an integrated paper and new proofs and evidence. The new constructive catalogue supplies a separate horizon/price study and repeated isolated timings for its declared backend. An adaptive-grid comparator, a nonlinear dimension frontier, and repeated isolated comparisons for the original optimizers remain distinct evidence requirements; retained-policy and paired-cost results are not substitutes for them.
\section{Constructive neural policies from primitive moduli}\label{app:constructive45}
This section supplies the complete implementation-independent proof of the constructive theorem in the main article. Its assumptions are a compact state set $X\subset\mathbb R^d$, a common compact action set $A\subset\mathbb R^a$, a finite horizon, nonnegative discount factors, an invariant transition, a common innovation law, and state/action-independent monotone cash-invariant conditional evaluation on a common bounded-function domain. The primitive moduli are those in the main article's constructive section. Effective compactness and finite certified numerical evaluation are used only for finite termination, not for the analytic inequalities.

\subsection{Bellman moduli and labels}
If $\|u-v\|_\infty\le r$, then $v-r\le u\le v+r$. Monotonicity and cash invariance give $|\rho_t(u)-\rho_t(v)|\le r$. Consequently, when $f_{t+1}$ is $L_{t+1}$-Lipschitz,
\[
 |Q_t(f_{t+1};x,a)-Q_t(f_{t+1};y,b)|
 \le L_t\|x-y\|_1+L^a_t\|a-b\|_1,
\]
where $L_t=C_{x,t}+\beta_tF_{x,t}L_{t+1}$ and $L^a_t=C_{a,t}+\beta_tF_{a,t}L_{t+1}$. Infima over the same action set preserve the state modulus. Compactness and continuity give minima. Alternatively, the same argument uses arbitrarily close minimizers and then takes their error to zero.

Let $H_t=T_tf_{t+1}$, $q_{ij}=Q_t(f_{t+1};x_i,a_j)$, and $|\widehat q_{ij}-q_{ij}|\le e$. The state and action grids have radii $h,k$. Since every grid action is feasible,
\[
 y_i=\min_j\widehat q_{ij}\ge \min_jq_{ij}-e\ge H_t(x_i)-e.
\]
Choose an action minimizer $a^*$ and a grid action within $k$ of it. Then
\[
 y_i\le q_{ij^*}+e\le H_t(x_i)+L^ak+e.
\]
These are bounds on the label array actually produced by the own-future recursion. They do not use $V_{t+1}^*$ and do not assume that a regression loss is small.

\subsection{Envelope, residual, and selected action}
For arbitrary finite labels $y_i$, define $E_y(x)=\min_i\{y_i+L\|x-x_i\|_1\}$. Each cone is $L$-Lipschitz. For a minimizing cone at $x'$, $E_y(x)-E_y(x')\le L\|x-x'\|_1$, and reversing the points proves the Lipschitz assertion. If $\max_i|y_i-y'_i|\le\delta$, then every cone shifts by at most $\delta$, so $\|E_y-E_{y'}\|_\infty\le\delta$. In particular, numerical label errors cannot inflate the certified slope.

For every node, $y_i+L\|x-x_i\|_1\ge H_t(x_i)-e+L\|x-x_i\|_1\ge H_t(x)-e$. Taking the minimum proves the lower residual bound. Choosing a nearest node $i(x)$ gives
\[
 E_y(x)\le H_t(x_{i(x)})+L^ak+e+Lh
        \le H_t(x)+L^ak+e+2Lh.
\]
If $j(i)$ minimizes the numerical queries, then
\[
 q_{i,j(i)}\le\widehat q_{i,j(i)}+e=y_i+e
       \le H_t(x_i)+L^ak+2e.
\]
The difference of the chosen-action value and $H_t$ has state modulus at most $2L$. Thus the nearest-node actor has gap at most $L^ak+2e+2Lh$. A deterministic rule for ties among finitely many closed Voronoi cells gives a Borel actor. Any tied nearest node satisfies the same inequality. Because the action set is common, this actor is feasible at every state. This proof would not establish feasibility for an unrestricted state-dependent correspondence.

At the terminal date, $H_T=g$ and there is no action minimization. Labels within $e_T$ of $g$ yield $-e_T\le f_T-g\le e_T+2L_gh_T$. Starting from the terminal network and proceeding backward establishes the assumed continuation moduli at every date. The policy theorem, proved earlier in this supplement, then yields
\[
 J_0^\pi(x)-V_0(x)\le
 \sum_{t<T}B_t(2L^a_tk_t+4e_t+4L_th_t)
   +B_T(2e_T+2L_gh_T).
\]
The comparison is with the optimum over all admissible adapted policies, not just policies constant on the grid. The action covering allowance is precisely what makes the continuous-action comparison possible.

\subsection{Realization and finite termination}
An absolute value uses two ReLU gates. A sum of $d$ absolute values with affine coefficients computes one cone. The identity $\min(u,v)=u-\max(u-v,0)$ computes pairwise minima. A balanced finite tree consumes $K$ cones in $O(\log K)$ stages, with $O(K)$ additional gates. Positive and negative channels represent signed intermediate scalars in an ordinary feed-forward realization. The total number of coefficients is $O(dK)$; a bounded-fan-in realization can replace each affine sum by a logarithmic-depth addition tree. This is an explicit network with weights determined by labels, state nodes, and the primitive upper modulus.

All constants may be replaced by computable rational upper bounds. The labels can be rational rounded midpoints of certified numerical intervals. State/action covers must be supplied with verifiable covering and feasibility radii, rather than presumed because a set is topologically compact. For the budget $s=\varepsilon/\sum B_t$, the nonterminal terms satisfy
\[
 2(s/8)+4(s/8)+4(s/16)=s,
\]
and terminal terms satisfy $2(s/4)+2(s/4)=s$. Positive-radius covers have finitely many nodes and numerical queries reach the specified tolerances by hypothesis. Backward construction therefore takes finite work and the policy loss is at most $\varepsilon$. This is a deterministic algorithm theorem. It does not turn a finite implementation cap into a termination guarantee, and it does not provide a convergence theorem for Adam, stochastic gradient descent, or an arbitrary nonconvex loss.

The work expression in the article counts all state--action--innovation evaluations of the already constructed future network. If numerical precision increases with the requested tolerance, the bit cost of each operation and of each weight must also increase in the account. A tensor cover retains its exponential state/action dimension factors. Analytic or sparse integration can reduce the innovation work but does not remove the covering factors merely by being called neural computation.

\subsection{Exact scalar compilation}
Let $0=x_0<\cdots<x_N=1$ and let the labels be arbitrary rational numbers. Define
\[
 z_i=\min_j\{y_j+L|x_i-x_j|\}.
\]
A left-to-right pass computes the minimum over $j\le i$ and a right-to-left pass incorporates $j\ge i$. The result is the greatest $L$-Lipschitz vector lying below the labels: any such vector $w$ satisfies $w_i\le w_j+L|x_i-x_j|\le y_j+L|x_i-x_j|$ for every $j$. In particular $|z_{i+1}-z_i|\le L(x_{i+1}-x_i)$.

For $x\in[x_i,x_{i+1}]$, every cone centered to the left is represented by the best left intercept and every cone centered to the right by the best right intercept. Therefore
\[
 E_y(x)=\min\{z_i+L(x-x_i),\ z_{i+1}+L(x_{i+1}-x)\}.
\]
For $L>0$, the intersection is
\[
 r_i=\frac{x_i+x_{i+1}}2+\frac{z_{i+1}-z_i}{2L}\in[x_i,x_{i+1}].
\]
When the intersection is interior it is the only additional knot; endpoint intersections cause no duplicate knot. When $L=0$ the envelope is constant. This proves equality of the compiled piecewise-linear function and the defining network everywhere, including non-dyadic intersections and ties.

For rational knots $r_j$, values $v_j$, and slopes $s_j$, the antiderivative on one cell is
\[
 A(x)=A(r_j)+v_j(x-r_j)+\tfrac12s_j(x-r_j)^2.
\]
Rational trapezoidal accumulation computes $A(r_j)$ exactly. Under a continuous uniform innovation on $[-\sigma,\sigma]$, the own-future expectation is
$[A(b+\sigma)-A(b-\sigma)]/(2\sigma)$.
This is an identity under the original law, not a finite-bin change of the economy.

The numerical evaluator encloses each rational coefficient and each native arithmetic operation outward. At a query near a rounded non-dyadic knot, it resolves the cell using the exact rational knot and the exact rational value of the binary64 query. It does not use a rounded knot as a silent change of the partition. An uncertain antiderivative argument is evaluated at its midpoint and expanded by $\sup|f|$ times its radius; an uncertain innovation center is expanded by $\operatorname{Lip}(f)$ times its radius. Intersections with $[0,1]$ discard only rounding overhang when the exact primitive argument is known to lie in that invariant interval. These steps separate analytic integration from arithmetic rounding.

\subsection{Matched spline account and measurement scope}
For the spline comparator, linear interpolation of labels with node errors between $-e$ and $L^ak+e$ gives a residual interval
$[-e-Lh,\ L^ak+e+Lh]$.
Indeed the interpolating weights are nonnegative and sum to one, and their weighted distance to the query point is at most $h$ on the uniform one-dimensional mesh. Its width is $L^ak+2e+2Lh$, the same form as the envelope width. The nearest-node actor allowance is also the same form. The spline's actual exact slope, not the neural primitive upper bound, is propagated into its next date's query modulus. Both terminal certificates conservatively use the primitive terminal modulus. Neither comparator borrows the other's targets.

The catalogue records exact rational gap formulas, upper binary64 endpoints, all interval-error allowances, defining labels, compiled nodes, and selected actions for every rung. It also records evaluated Bellman queries, continuous-law integrals, antiderivative queries, binary-search upper counts, deployed actor scalars, rational storage quantities, maximum stored numerator/denominator bit lengths, checkpoint bytes, process peak resident memory, internal clocks, and process clocks. Primitive-object counts are not described as a complete floating-point or bit-operation count. Peak memory includes the interpreter and warm-up. The internal target clock includes failed rungs and durable checkpoint output; independent economic comparison is absent rather than costless. These scope qualifications are part of the empirical estimand.

A first-rung timing repetition is not a new random economic environment. For each fixed method/horizon/price combination, the three exact certificate arrays and deployed objects should coincide; their clocks need not. Equality of object hashes is audited separately from timing dispersion. This supplies reproducibility for the declared deterministic generator without inferring a probability of success for another training algorithm.

\section{Complete constructive verification frontiers}
\begin{longtable}{lrrrrrr}
\caption{All 48 distinct method--economy--resolution rows}\label{tab:frontier45}\\
\toprule
Generator & $T$ & $p$ & $N$ & Gap & Prefix queries & Median s\\
\midrule\endfirsthead
\toprule
Generator & $T$ & $p$ & $N$ & Gap & Prefix queries & Median s\\
\midrule\endhead
ReLU envelope & 2 & 1 & 16 & 1.96222 & 578 & 0.0104 \\
ReLU envelope & 2 & 1 & 32 & 0.98111 & 2,756 & 0.0252 \\
ReLU envelope & 2 & 1 & 64 & 0.49056 & 11,206 & 0.0534 \\
ReLU envelope & 2 & 1 & 128 & 0.24528 & 44,488 & 0.1200 \\
ReLU envelope & 2 & 1 & 256 & 0.12264 & 176,586 & 0.3009 \\
ReLU envelope & 2 & 1 & 512 & 0.06132 & 702,924 & 0.8474 \\
Spline & 2 & 1 & 16 & 1.74262 & 578 & 0.0072 \\
Spline & 2 & 1 & 32 & 0.88564 & 2,756 & 0.0167 \\
Spline & 2 & 1 & 64 & 0.44641 & 11,206 & 0.0345 \\
Spline & 2 & 1 & 128 & 0.22410 & 44,488 & 0.0800 \\
Spline & 2 & 1 & 256 & 0.11228 & 176,586 & 0.2144 \\
Spline & 2 & 1 & 512 & 0.05619 & 702,924 & 0.6777 \\
ReLU envelope & 2 & 4 & 16 & 2.00763 & 578 & 0.0101 \\
ReLU envelope & 2 & 4 & 32 & 1.00382 & 2,756 & 0.0248 \\
ReLU envelope & 2 & 4 & 64 & 0.50191 & 11,206 & 0.0529 \\
ReLU envelope & 2 & 4 & 128 & 0.25095 & 44,488 & 0.1192 \\
ReLU envelope & 2 & 4 & 256 & 0.12548 & 176,586 & 0.2944 \\
ReLU envelope & 2 & 4 & 512 & 0.06274 & 702,924 & 0.8450 \\
Spline & 2 & 4 & 16 & 1.80513 & 578 & 0.0075 \\
Spline & 2 & 4 & 32 & 0.91577 & 2,756 & 0.0168 \\
Spline & 2 & 4 & 64 & 0.46103 & 11,206 & 0.0345 \\
Spline & 2 & 4 & 128 & 0.23136 & 44,488 & 0.0799 \\
Spline & 2 & 4 & 256 & 0.11587 & 176,586 & 0.2151 \\
Spline & 2 & 4 & 512 & 0.05799 & 702,924 & 0.6740 \\
ReLU envelope & 4 & 1 & 16 & 3.91560 & 1,156 & 0.0174 \\
ReLU envelope & 4 & 1 & 32 & 1.95780 & 5,512 & 0.0430 \\
ReLU envelope & 4 & 1 & 64 & 0.97890 & 22,412 & 0.0927 \\
ReLU envelope & 4 & 1 & 128 & 0.48945 & 88,976 & 0.2132 \\
ReLU envelope & 4 & 1 & 256 & 0.24473 & 353,172 & 0.5335 \\
ReLU envelope & 4 & 1 & 512 & 0.12236 & 1,405,848 & 1.5540 \\
Spline & 4 & 1 & 16 & 2.99898 & 1,156 & 0.0127 \\
Spline & 4 & 1 & 32 & 1.52738 & 5,512 & 0.0293 \\
Spline & 4 & 1 & 64 & 0.77070 & 22,412 & 0.0618 \\
Spline & 4 & 1 & 128 & 0.38711 & 88,976 & 0.1455 \\
Spline & 4 & 1 & 256 & 0.19399 & 353,172 & 0.4028 \\
Spline & 4 & 1 & 512 & 0.09711 & 1,405,848 & 1.2654 \\
ReLU envelope & 4 & 4 & 16 & 4.00092 & 1,156 & 0.0171 \\
ReLU envelope & 4 & 4 & 32 & 2.00046 & 5,512 & 0.0424 \\
ReLU envelope & 4 & 4 & 64 & 1.00023 & 22,412 & 0.0931 \\
ReLU envelope & 4 & 4 & 128 & 0.50012 & 88,976 & 0.2125 \\
ReLU envelope & 4 & 4 & 256 & 0.25006 & 353,172 & 0.5349 \\
ReLU envelope & 4 & 4 & 512 & 0.12503 & 1,405,848 & 1.5495 \\
Spline & 4 & 4 & 16 & 3.17460 & 1,156 & 0.0128 \\
Spline & 4 & 4 & 32 & 1.61224 & 5,512 & 0.0293 \\
Spline & 4 & 4 & 64 & 0.81227 & 22,412 & 0.0612 \\
Spline & 4 & 4 & 128 & 0.40770 & 88,976 & 0.1436 \\
Spline & 4 & 4 & 256 & 0.20423 & 353,172 & 0.4038 \\
Spline & 4 & 4 & 512 & 0.10221 & 1,405,848 & 1.2705 \\
\bottomrule\end{longtable}
Every row represents three byte-identical policy and certificate checkpoints with distinct clocks. Complete counts, storage, memory, and the individual clocks are deposited in the machine-readable records. The full frontier includes unsuccessful rungs; no refined neural row is substituted for an original failed service.

\section{Witness-preserving policy construction}\label{app:witness46}
This section establishes the additional policy sandwich and the witness-preserving implementation. It uses the same primitive economy, admissible-policy comparison and common evaluation domain as the corresponding main-text results. All inequalities concern the continuation actually constructed from its own fitted future.

\subsection{Order comparison and signed accounts}
Suppose $f_t-\mathcal T_tf_{t+1}\le u_t$ and $\mathcal T_t^\pi f_{t+1}-f_t\le d_t$. Define
\[
 U_t=\sum_{j=t}^{T-1}\left(\prod_{k=t}^{j-1}\beta_k\right)u_j
       +\left(\prod_{k=t}^{T-1}\beta_k\right)u_T,
\]
with empty products equal to one. The terminal inequality gives $V_T^*=g\ge f_T-U_T$. If $V_{t+1}^*\ge f_{t+1}-U_{t+1}$, monotonicity and cash invariance yield
\[
 V_t^*=\mathcal T_tV_{t+1}^*
 \ge\mathcal T_tf_{t+1}-\beta_tU_{t+1}
 \ge f_t-u_t-\beta_tU_{t+1}=f_t-U_t.
\]
The same induction for the fixed-policy map, beginning at $g\le f_T-\ell_T$, gives
\[
 J_0^\pi\le f_0+\sum_{t<T}B_td_t-B_T\ell_T.
\]
Subtracting proves the one-sided sandwich. Neither $u_t$ nor $d_t$ is replaced by an absolute value; their signs cancel the effect of arbitrary date-specific changes of origin. The lower loss bound follows from Bellman optimality against all admissible adapted policies. This argument uses no independence of simulations and no linearity of the recursive functional. It is distinct from the expectation-specific paired telescoping theorem retained earlier in this supplement.

The old centered certificate is the special bound $d_t\le\eta_t-\ell_t$. Its possible slack is due to taking two separate extrema of quantities that share the same fitted value. The new witness result controls their sum directly.

\subsection{Action transport through the minimum circuit}
At a node $x_i$, the numerical action query attaining $y_i$ corresponds to an actual feasible action $a_i$. The query enclosure gives $Q(x_i,a_i)\le y_i+e$. Uniform state Lipschitz continuity of $Q(\cdot,a)$ therefore implies
\[
 Q(x,a_i)-f(x)\le e+y_i+L\|x-x_i\|_1-f(x)
\]
for every node and every state. Selecting an index attaining the minimum cone makes the last three terms sum to zero. Finite minima attain their values, and a fixed ordering of the finite indices defines a Borel selector. The common action set is essential: the proof does not transport an action into a state at which that action is infeasible.

Combining $d_t=e_t$ with $u_t=L_t^ak_t+e_t+2L_th_t$ proves the displayed witness gap. It also implies the separated actor allowance
\[
 Q_t(f_{t+1};x,\pi_t^{\rm w}(x))-\mathcal T_tf_{t+1}(x)
 \le e_t+u_t=L_t^ak_t+2e_t+2L_th_t.
\]
Thus the earlier numerical separated budget is also a valid, though weaker, bound for the new witness actor. Using it as an alternative does not substitute the old nearest-node policy. At each nonterminal date, the new sufficient resource choices obey $s/4+2(s/4)+2(s/8)=s$. The terminal bracket is also at most $s$. Effective finite covers and certified numerical evaluation then give finite termination.

Only the critic is an affine--ReLU circuit. Returning a comparison index and an attached action is a finite measurable algorithm with a possibly discontinuous policy. Classical Lipschitz extension and elementary minimum gates remain credited in the main article. The result does not assert their novelty or claim that an equivalent conventional envelope implementation is impossible.

\subsection{Approximate indices and numerical deployment}
Let $x'$ be an acquired state with $\|x-x'\|_1\le r$. Suppose the selected index satisfies
\[
 y_i+L\|x'-x_i\|_1\le f(x')+\nu.
\]
The cone and the envelope are both $L$-Lipschitz. Hence
\[
 y_i+L\|x-x_i\|_1-f(x)
 \le \nu+2Lr.
\]
A feasible executed action within distance $\alpha$ of $a_i$ adds at most $L^a\alpha$ to $Q$. This proves the stated execution allowance. Rounding a rational switching boundary without a checked comparison can change the selected index; it must be covered by $\nu$ or by a certified interval containing all eligible indices. An exact tie presents no difficulty, since both attaining cones give the same certified upper bound. The current records store exact rational boundaries and define the ideal selector. They do not claim a measured deployment performance for arbitrary hardware.

\subsection{Scalar compilation with original witnesses}
Let $0=x_0<\cdots<x_N=1$ and let the original numerical labels $y_i$ be arbitrary rational numbers. Define pairs
\[
 (p_i,l_i)=\min_{j\le i}(y_j+L(x_i-x_j),j),\qquad
 (q_i,r_i)=\min_{j\ge i}(y_j+L(x_j-x_i),j),
\]
where pairs are ordered first by value and then by index. A forward pass computes the first array by comparing $(y_i,i)$ with $(p_{i-1}+L(x_i-x_{i-1}),l_{i-1})$; a backward pass computes the second. The index always refers to an original label and its action, not to a newly regularized value with no action witness.

On $[x_i,x_{i+1}]$, the envelope is the minimum of
\[
 p_i+L(x-x_i)\quad\hbox{and}\quad
 q_{i+1}+L(x_{i+1}-x),
\]
with witnesses $l_i$ and $r_{i+1}$. When $L>0$ their intersection is
\[
 s_i=\frac{x_i+x_{i+1}}2+\frac{q_{i+1}-p_i}{2L}.
\]
Unlike an intersection computed after closing both labels, this intersection may lie outside the cell. In that case one affine candidate attains the minimum throughout the cell. An interior intersection divides the cell into exactly two witness regions. If $L=0$, compare the two constants directly. The implementation preserves every original node boundary and every interior switch, and chooses the right cell at an exact boundary. Continuity ensures that an adjoining witness remains attaining at the boundary. Thus there are at most $2N$ witness intervals, with rational endpoints and original action indices.

Once rational endpoints are loaded, an ordered search locates a region using a logarithmic number of comparisons. Reading and parsing all stored rational cuts anew costs linear work and is not removed by that comparison bound. The reference accessor used by the tests performs that parsing; the catalogue measures construction and serialization, not repeated policy-query throughput. The record's binary-search count is only a count of ordered comparisons, not the complete cost of that accessor. Rational numerator/denominator lengths, cuts, witness indices and original actions are stored separately.

\subsection{A common strengthened nearest-node certificate}
For a nearest-node policy and any fitted piecewise-affine continuation, let
\[
 D=\max_i\sup_{x\in C_i}\{y_i+L|x-x_i|-f(x)\}.
\]
The selected-node query and the state modulus give $Q(x,\pi(x))-f(x)\le e+D$. On each intersection of a nearest-node cell with a linearity interval of $f$, divided also at $x_i$, the expression is affine. Its maximum is attained at an endpoint. Consequently it suffices to enumerate the union of the fitted knots, original state nodes and nearest-cell boundaries. At a shared nearest-cell boundary both possible node indices are included. This is an exact all-state calculation, rather than a sampled maximum.

Both the cone-nearest ablation and spline-nearest comparator use this calculation and their own valid optimal-residual upper bounds. The former has $D\ge0$ because the cone envelope lies below every candidate cone. Its witness counterpart achieves zero for this particular label-transport upper account. This statement is about the account, not an optimality theorem for the actual policy among all feasible controls. The spline's continuation is not replaced by the cone envelope and continues to propagate its own exact Lipschitz constant.

\subsection{Evidence identities and scope of reproduction}
The protocol predates the complete 36-service execution. Fourteen new checks examine exact compilation, noisy and dominated labels, ties, finite-policy comparison, shift cancellation, common nearest-cell bounds, exact continuous-uniform integration in small examples, and immutable numerical dependencies. They supplement the analytic proof rather than replace it. The source hashes, all 216 checkpoints, and complete process clocks are retained. Formula reconstruction verifies each one-sided component, every first crossing, and the identity of the two neural critics in the actor ablation. Three repetitions of a deterministic policy are not three economic draws.

The latest R43 supplemental report concerns a different, earlier incomplete snapshot. The current reply does not claim that its missing source parts were present. It identifies the complete R46 source tree as the new review object and retains the earlier review and incomplete historical branch unchanged. The predeclared R46 execution, R45 construction catalogue, R44 postselected diagnostics, original R42 capped services and original paired cost comparisons remain distinct evidence objects.

\begin{longtable}{rrrrrr}
\caption{Complete matched witness frontiers}\label{tab:witnessfront46}\\
\toprule $T$ & $p$ & $N$ & Witness & Nearest neural & Spline \\ \midrule\endfirsthead
\toprule $T$ & $p$ & $N$ & Witness & Nearest neural & Spline \\ \midrule\endhead
2 & 1 & 16 & 1.097841 & 1.558158 & 1.222038 \\
2 & 1 & 32 & 0.548921 & 0.790405 & 0.622981 \\
2 & 1 & 64 & 0.274460 & 0.398115 & 0.314494 \\
2 & 1 & 128 & 0.137230 & 0.199802 & 0.158000 \\
2 & 1 & 256 & 0.068615 & 0.100089 & 0.079188 \\
2 & 1 & 512 & 0.034308 & 0.050092 & 0.039641 \\
2 & 4 & 16 & 1.120546 & 1.615348 & 1.269376 \\
2 & 4 & 32 & 0.560273 & 0.815732 & 0.645500 \\
2 & 4 & 64 & 0.280137 & 0.411198 & 0.325309 \\
2 & 4 & 128 & 0.140068 & 0.206040 & 0.163341 \\
2 & 4 & 256 & 0.070034 & 0.103215 & 0.081831 \\
2 & 4 & 512 & 0.035017 & 0.051637 & 0.040957 \\
4 & 1 & 16 & 2.060395 & 2.916323 & 2.035153 \\
4 & 1 & 32 & 1.030197 & 1.478057 & 1.040370 \\
4 & 1 & 64 & 0.515099 & 0.744163 & 0.525918 \\
4 & 1 & 128 & 0.257549 & 0.373386 & 0.264396 \\
4 & 1 & 256 & 0.128775 & 0.187021 & 0.132558 \\
4 & 1 & 512 & 0.064387 & 0.093593 & 0.066369 \\
4 & 4 & 16 & 2.103056 & 3.034798 & 2.161469 \\
4 & 4 & 32 & 1.051528 & 1.533865 & 1.100879 \\
4 & 4 & 64 & 0.525764 & 0.772861 & 0.555422 \\
4 & 4 & 128 & 0.262882 & 0.387323 & 0.278981 \\
4 & 4 & 256 & 0.131441 & 0.193954 & 0.139801 \\
4 & 4 & 512 & 0.065720 & 0.097039 & 0.069979 \\
\bottomrule\end{longtable}

\begin{longtable}{llrrrr}
\caption{Full-frontier isolated construction clocks in seconds}\label{tab:witnesstime46}\\
\toprule Method & $(T,p)$ & Median & Minimum & Maximum & Peak MiB \\ \midrule\endfirsthead
\toprule Method & $(T,p)$ & Median & Minimum & Maximum & Peak MiB \\ \midrule\endhead
cone-witness & $(2,1)$ & 0.9631 & 0.9595 & 0.9817 & 113.8 \\
cone-witness & $(2,4)$ & 0.9596 & 0.9583 & 0.9709 & 113.9 \\
cone-witness & $(4,1)$ & 1.8000 & 1.7928 & 1.8025 & 116.4 \\
cone-witness & $(4,4)$ & 1.7456 & 1.7104 & 1.7899 & 116.4 \\
cone-nearest & $(2,1)$ & 1.1797 & 1.1739 & 1.2017 & 113.7 \\
cone-nearest & $(2,4)$ & 1.1798 & 1.1691 & 2.1261 & 113.7 \\
cone-nearest & $(4,1)$ & 2.2305 & 2.1667 & 2.2802 & 115.4 \\
cone-nearest & $(4,4)$ & 2.2163 & 2.2158 & 2.2221 & 115.4 \\
spline-nearest & $(2,1)$ & 0.8632 & 0.8424 & 0.8706 & 113.1 \\
spline-nearest & $(2,4)$ & 0.8755 & 0.8725 & 0.8964 & 113.0 \\
spline-nearest & $(4,1)$ & 1.6690 & 1.6628 & 1.6888 & 115.1 \\
spline-nearest & $(4,4)$ & 1.6707 & 1.5893 & 1.6783 & 115.1 \\
\bottomrule\end{longtable}
Clocks include all six rungs through checkpoint fsync. Peak resident memory includes the interpreter and warm-up. These are three local repetitions, not independent economic draws.


\section{Feasible transport under endogenous constraints}\label{app:feasible46}
We give the details underlying the feasible-witness transport theorem in the main article. All norms below are the specified $\ell^1$ norms. The evaluation functional is the same at the compared state--action pairs. Its monotonicity and cash invariance imply that changing its bounded random input by at most $b$ changes its value by at most $b$. Thus the primitive state and action inequalities on the feasible graph imply
\[
 |Q_t(f_{t+1};x,a)-Q_t(f_{t+1};z,b)|
 \le A_t^Q\|x-z\|_1+D_t^Q\|a-b\|_1
\]
whenever $a\in A_t(x)$ and $b\in A_t(z)$. The proof never evaluates $Q_t$ outside that graph.

\subsection{The Bellman modulus and noisy feasible labels}
Fix $a\in A_t(x)$. Compactness and the Hausdorff condition give $b\in A_t(z)$ with $\|a-b\|_1\le K_t^A\|x-z\|_1$. It follows that
\[
 \mathcal T_tf_{t+1}(z)
 \le Q_t(f_{t+1};x,a)+L_t\|x-z\|_1.
\]
Taking the infimum over $a$ and then reversing $x,z$ establishes the claimed Bellman modulus, without assuming that the minimizer is known. At node $i$, the minimum of exact objective values over its feasible action net is at least the continuous feasible infimum and at most that infimum plus $D_t^Qk_t$. Taking the minimum of queries with absolute error at most $e_t$ adds at most $e_t$ in each direction.

For every cone, the lower label inequality and the Bellman modulus give
\[
 y_{t,i}+L_t\|x-x_{t,i}\|_1\ge\mathcal T_tf_{t+1}(x)-e_t.
\]
For a nearest state node, the upper label inequality gives an upper bound $\mathcal T_tf_{t+1}(x)+D_t^Qk_t+e_t+2L_th_t$. The finite minimum has the same bounds and is $L_t$-Lipschitz. The next backward step therefore has precisely the modulus required by the construction, independently of the numerical label errors.

\subsection{Repair, selection and the dynamic account}
The stored action $a_{t,i}$ attains an original numerical query; hence its exact node objective is at most $y_{t,i}+e_t$. For the minimizing cone, the feasible-pair inequality and repair displacement give
\[
 Q_t(f_{t+1};x,\Gamma_{t,i}(x))-f_t(x)
 \le e_t+D_t^Q\zeta_t.
\]
A fixed ordering of finitely many continuous cone values gives a Borel index selector. Composing the finitely many Borel repairs with that selector gives a feasible Borel policy. The one-sided sandwich now supplies the discounted loss bound, with the original terminal interval unchanged. The full-policy comparison is exactly the Bellman comparison assumed in the controlled economy; it is not a comparison restricted to the node action arrays.

The four nonterminal resource allocations make the bracket at most $s/4+s/4+s/4+s/4=s$. The terminal bracket is at most $s$. Effective finite covers, feasible action nets, query enclosures and repair procedures consequently give finite termination. Where only an approximate repair is computable, its certified displacement is charged through $D_t^Q\zeta_t$ and its actual action must still be feasible. A numerical tolerance for an infeasible constraint violation is not included in this theorem.

\subsection{A transparent exact check}
Let $X=[0,1]^2$, $b(x)=(x_1+x_2)/2$, $A(x)=[0,b(x)]$, and consider the one-date objective
\[
 Q(x,a)=(a-3/4)^2+x_1/8+x_2/16.
\]
Here $A^Q=1/8$, $D^Q=3/2$, $K^A=1/2$, and $L=7/8$ are valid moduli. The exact optimum is attained at $\min\{3/4,b(x)\}$. For a feasible anchor action, clipping at $b(x)$ changes the action by at most $|b(x)-b(x_i)|\le\|x-x_i\|_1/2$. The transport theorem therefore applies directly, and a reference implementation can compare the exact repaired-policy loss with the analytic bound using rational arithmetic.

The additional tests independently form feasible node nets and noisy labels, compute the cone minimum by direct enumeration, and check the Bellman and selected-policy enclosures on a fixed rational set. They also check the Hausdorff modulus, exact feasibility, approximate-repair accounting, nonminimizing-index allowances and the zero-constraint-variation reduction. The analytic proof above, rather than a finite grid test, establishes the all-state claim. The tests are not training data, timing observations, or evidence of a high-dimensional speed advantage. No R46 catalogue record is altered by these checks.

\section{Acquired states and implementable witnesses}\label{supp:acquisition47}
This section supplies the implementation and statistical details for the new acquisition-aware results. The earlier proofs remain in force on their stated domains. To avoid conflating economic objects, the constrained study and the direct scalar comparison are separate experiments: the former constructs new two-state policies, while the latter evaluates the actual scalar policies previously constructed.

\subsection{Causal information and the one-sided comparison}
Let $\mathcal H_t$ contain the primitive history and the observations available before date-$t$ action. The observation process is causal and preserves the stipulated conditional law, or conditional recursive evaluator, for the next primitive innovation. This condition permits state errors to depend arbitrarily on the available past. It does not permit a sensor to reveal a future innovation without changing the economic information structure. The full-information optimum is over feasible policies adapted to this permitted history. Uniform one-step inequalities therefore apply to observation-dependent actions as well as exact-state Markov selectors.

Suppose $f_t-\mathcal T_tf_{t+1}\le u_t$, and suppose every permitted observation and true state obey $\mathcal T_t^{\widehat\pi}f_{t+1}-f_t\le d_t$. If $\ell_T\le f_T-g\le u_T$, induction gives
\[
 V_t^*\ge f_t-\sum_{j=t}^{T-1}B_{t,j}u_j-B_{t,T}u_T,
 \qquad
 J_t^{\widehat\pi}\le f_t+\sum_{j=t}^{T-1}B_{t,j}d_j-B_{t,T}\ell_T.
\]
The second inequality is interpreted on $\mathcal H_t$ when the policy uses the observation history. Monotonicity and cash invariance justify every induction step on the common evaluation domain. Subtraction yields the sum of $u_j+d_j$ and terminal width; no actor continuity is used. A date shift $f_t\mapsto f_t+s_t$ adds $s_t-\beta_ts_{t+1}$ to $u_t$ and subtracts the same quantity from $d_t$. The policy-loss account is invariant.

At a selected node, retain the original action $a_{t,i}$ and a label whose numerical error is at most $e_t$. For a feasible deployed action $\widehat a$, state and action transport give
\[
 Q_t(f_{t+1};x,\widehat a)
 \le y_{t,i}+e_t+A_t^Q\|x-x_{t,i}\|_1
                  +D_t^Q\|\widehat a-a_{t,i}\|_1.
\]
Substitute the robust-repair inequality and $L_t=A_t^Q+D_t^QK_t^A$. The cone at $x$ is within $L_tr_t$ of its value at the acquired state $\widehat x$. An index with numerical excess $\nu_t$ is within $\nu_t$ of the envelope there. Finally the envelope, being a minimum of $L_t$-Lipschitz functions, is $L_t$-Lipschitz, adding another $L_tr_t$. Hence
\[
 Q_t(f_{t+1};x,\widehat a)-f_t(x)
 \le e_t+\nu_t+2L_tr_t+D_t^Q\zeta_t^{\rm acq}.
\]
Combining this with the optimal-residual upper bound $D_t^Qk_t+e_t+2L_th_t$ proves the main theorem. This proof charges an objective excess even when two selected actions differ substantially at a boundary.

For a capacity interval, exact repair at the true state has displacement $(a_i-\kappa(x))_+\le K^A\|x-x_i\|_1$. Replacing capacity by a nonnegative certified lower value changes the repaired action by at most $\kappa(x)-\underline\kappa(\widehat x)$. Downward quantization adds less than $\Delta$. These inequalities prove both feasibility and the stated displacement bound. In the executed affine-capacity economy the acquisition box has coordinate endpoints $\max(0,\widehat x_j-\eta)$ and $\min(1,\widehat x_j+\eta)$. Its minimum capacity is computed at the lower corner. All these quantities are dyadic, so that the capacity and downward action rounding are exact at the specified resolutions.

\subsection{Finite-precision indices and the identical envelope}
The scalar and two-state experiments store labels on a dyadic lattice. In the acquired-state deployment, nodes and measured coordinates have the declared dyadic denominators; their coordinate differences and the sum of two absolute differences are exactly representable in binary64. No overflow or underflow occurs in the evaluated range. For a cone $v_i=y_i+L\|\widehat x-x_i\|_1$, let $\widetilde L$ denote the binary64 slope and $\widetilde v_i$ its computed score. The standard product-and-sum bound, with unit roundoff $u=2^{-53}$, gives
\[
 |\widetilde v_i-v_i|
 \le d|\widetilde L-L|+\gamma_2(Y+d|\widetilde L|)=\tau,
 \qquad \gamma_2=\frac{2u}{1-2u},\quad Y=\max_i|y_i|.
\]
If $i$ minimizes the computed scores and $j$ minimizes the exact scores, then $v_i\le\widetilde v_i+\tau\le\widetilde v_j+\tau\le v_j+2\tau$. Thus $\nu=2\tau$ works regardless of a tie. More general state formats must charge the formation of distances too; that cost is not suppressed by this dyadic example.

An affine--ReLU minimum gate is $a-(a-b)_+$. For exact inputs it equals the native minimum. Attaching the same index and the same tie rule proves equality of the witness actor, including feasible repair. The proof is algebraic and covers all states, not only the measurement grid. The scalar implementation's compiled right-cell tie convention is held fixed in this control.

Ordinary rounded ReLU gates need not agree bit for bit with native minima. If the exact intermediate minima have absolute value at most $M$, input errors are at most $E$, and two subtractions implement one gate, a conservative next-level bound is
\[
 E_{\rm next}\le E+6u(M+E)
\]
under the preceding finite-range conditions. Indeed the first subtraction contributes at most $2u(M+E)$, the positive-part map is nonexpansive, and the last subtraction contributes at most $u(3+2u)(M+E)$. The exact minimum is nonexpansive in its two arguments. Iteration from the leaf conversion bound controls a balanced circuit. The retained diagnostic checks all 24 scalar witness checkpoints on 1,025 dyadic points against this budget. Its gate-evaluation clock excludes construction and is not a method-level speed comparison. The native representation and the ReLU representation share parameters; they are not counted as independent trained policies.

\section{Explicit constrained primitives and resource allocation}\label{supp:constrained47}
For the new two-state economy, the action lies in $[0,1/4]$. Bounding the nonnegative state terms separately gives $F_1\in[1/32,29/32]$ and $F_2\in[1/32,27/32]$. These bounds already establish invariance before any numerical enclosure is intersected with the domain. The two state-column sums of the transition Jacobian are $11/16-x_2/8$ and $11/16-x_1/8$; every term used in these sums is nonnegative. The induced $\ell^1$ state modulus is therefore $11/16$. The action column has norm $1/2+1/4=3/4$. Capacity has gradient $(1/16,1/16)$ and hence $\ell^1$ Lipschitz constant $1/16$.

Let $s=(1/2-x_1-x_2)_+$. The first state derivative of stage cost is $2(x_1-5/8)+(x_1-x_2)/2-4s$; the second is the symmetric expression. On each of the two shortage regions these derivatives are affine. At the vertices of the resulting polygons their absolute values are at most $13/4$. Doubling the squared-target part changes this bound to $9/2$ for terminal cost. The action derivative $2pa+16a^3$ is nonnegative and at most $p/2+1/4$. These are constants for the primitive equations, not estimates from the evaluation grid.

For a uniform tensor grid, each coordinate is within $1/(2N)$ of a node, giving $h=1/N$. The maximum feasible interval length is $1/4$, so a uniform $N$-subdivision feasible net has radius at most $1/(8N)$. At different states the capacity intervals have Hausdorff distance at most $(1/16)\|x-x'\|_1$. The map $a\mapsto\min\{a,\kappa(x)\}$ is measurable and feasible, with the displacement bound proved above. Thus neither an unspecified covering oracle nor an unspecified feasible projection is used here.

A shock cell has length $1/(16M)$. Its conditional mean distance to the midpoint is $1/(64M)$. The transition's shock derivative is $(1,-1)$, of $\ell^1$ norm two. Consequently a future continuation of modulus $L$ has midpoint expectation error at most $L/(32M)$. The program encloses numerical evaluation at all midpoints, averages outward, adds this continuous-law remainder, and discounts it. Labels are dyadic rounded centers with an error enclosing the full numerical query. The economic law remains continuous; this is not a proof for a substituted finite-state shock process.

\subsection{The strengthened bilinear comparator}
Write $\lambda_i(x)$ for the four bilinear interpolation weights in a tensor cell. They are nonnegative, sum to one and reproduce the two coordinates. In one coordinate, at relative position $q\in[0,1]$, the weighted distance to the two endpoints is $2q(1-q)/N\le1/(2N)$. Adding coordinates gives $\sum_i\lambda_i(x)\|x-x_i\|_1\le1/N$. If the nodal Bellman labels have lower error $-e_t$ and upper error $D_t^Qk_t+e_t$, the graph-Lipschitz Bellman target therefore obeys
\[
 -e_t-L_t/N\le f_t(x)-\mathcal T_tf_{t+1}(x)
 \le D_t^Qk_t+e_t+L_t/N.
\]
The interpolant's own Lipschitz modulus is $N$ times the largest adjacent nodal difference in absolute value. Each partial derivative is a convex combination of two adjacent differences. In this execution the differences are dyadic with sufficiently few significant bits to be represented exactly; the computed modulus is consequently an upper bound without a hidden downward rounding.

For the nearest-node actor, each closed Voronoi cell is split at its node coordinates. On each of the at most four resulting rectangles the absolute-distance term is affine, and the continuation is a single bilinear polynomial. Fixing one coordinate makes the difference affine in the other. Successively maximizing both coordinates therefore attains a maximum at a vertex. The union contains at most nine points. Evaluating every point outward bounds the complete supremum, including boundary ties. This proves the selected-policy allowance used by the comparator rather than comparing a sharpened witness bound to an unnecessarily separated baseline.

\subsection{Sufficient caps versus executed caps}
For the cone generator,
\[
 L_t=C_x+K^AC_a+\beta(F_x+K^AF_a)L_{t+1}.
\]
Here $b=\beta(F_x+K^AF_a)=705/1024<1$, giving the horizon-uniform bound $\overline L=\max\{L_g,(C_x+K^AC_a)/(1-b)\}$. The action modulus is at most $\overline D=C_a+\beta F_a\overline L$. With $M=N$, nonterminal label error is at most $\beta\overline L/(32N)+\xi$. Substituting into the one-sided loss bound yields exactly
\[
 \mathcal G_N\le\frac{S_T(\overline D/8+2\overline L+\beta\overline L/16)+2\beta^TL_g}{N}
                      +2S_T\xi+2\beta^T\xi_T.
\]
Choosing the spatial and numerical allocations as in the article proves the sufficient cap. It is not a statement that the fixed ladder $4,8,16$ contains that cap, nor that the labels rounded to a fixed lattice support arbitrary tolerances. Acquisition, quantization and selector excess are additional summands in the same budget.

There are $(N+1)^2$ nodes, $N+1$ feasible actions per node, $N$ shock midpoints per query and $T$ backward dates. A direct cone query evaluates $(N+1)^2$ pieces, while a bilinear query evaluates one local patch. The resulting operation bounds are $O(TN^6)$ and $O(TN^4)$, respectively. These formulas include label formation, not only verification. Actor storage and critic-label storage are $O(TN^2)$. Circuit compilation, cell verification, selector queries, feasible repair, numerical enclosure, coefficient conversion and serialization are separately counted or bounded in the stated implementation. Arbitrary precision multiplies arithmetic counts by the cost of the chosen bit operations and requires enough guard precision to meet the query error. These are upper bounds for the specified routines, not lower bounds against alternative algorithms. In general dimension the original state and action covering factors remain visible. No computational saving follows from relabeling the identical envelope as a neural circuit.

\section{Direct policy values with numerical and sampling error}\label{supp:direct47}
\subsection{Identity and deterministic support}
For a fixed policy $m$ and its own fitted continuation, let
\begin{align*}
 Z_m={}&f_0^m(X_0)+B_T(g(X_T)-f_T^m(X_T))\\
 &+\sum_{t<T}B_t\{c_t(X_t,A_t^m)+\beta_t\E[f_{t+1}^m(X_{t+1})\mid\mathcal H_t]-f_t^m(X_t)\}.
\end{align*}
Conditional expectation and telescoping show $\E Z_m=J_m$ for either a fixed or random initial state. Policies in a pair may be driven by common shocks; preservation of each marginal law is sufficient. The innovation dependence between the two policies changes variance, not this identity.

For $f_t^m-\mathcal T_tf_{t+1}^m\le u_t^m$ and $\mathcal T_t^{\pi_m}f_{t+1}^m-f_t^m\le d_t^m$, the realized conditional residual is in $[-u_t^m,d_t^m]$. The lower bound follows because a feasible action cannot be below the Bellman infimum. Adding the ranges of these residuals, the terminal differences, and the initial difference gives the deterministic support in the article. For a common random initial state the initial difference is piecewise affine on the union of the two knot sets; extrema occur at those knots. For fixed initial states it is evaluated exactly as a rational number. Initial randomness is therefore not replaced by an unsupported average over four hand-selected points.

The frozen scalar terminal allowance is reconstructed from its stored width. If terminal node error is $e_T$, a cone continuation gives $g-f_T\in[-e_T-2L_gh_T,e_T]$, and a spline gives $g-f_T\in[-e_T-L_gh_T,e_T+L_gh_T]$. All signs are retained when forming pair support. Subtracting separate regret bounds would not give this estimand.

\subsection{Bins, discontinuities, and endpoint statistics}
Each independent uniform integer specifies one of $2^{40}$ equal innovation bins, and, for the uniform initial-state law, one such initial bin. The lower and upper path computation encloses every continuous point of that bin. Conditional uniform randomization within each independent bin recovers the original continuous uniform law. It need not be explicitly generated because every such value is already enclosed.

At a selector boundary, stored rational cuts are enclosed outward. Binary search obtains the first and last cells intersecting the state interval. The action interval is the hull of all actions in those cells, not merely its two endpoint actions. Propagating the hull may be conservative but is valid even for a nonmonotone actor. No path is removed. Intersections with the proved invariant state domain and deterministic score support preserve enclosure. The stored ambiguity counts quantify this part of the computation. These endpoint observations are independent across paths under the stated sampling model, even though endpoints and contrasts within a path are dependent.

Here are the numerical details of the empirical Bernstein calculation. For a binary64 endpoint sample $x_1,\ldots,x_n\in[a,b]$, let $M=\max(|a|,|b|)$ and $\gamma=(n+4)u/(1-(n+4)u)$. Denote the computed mean and mean square by $\widetilde\mu,\widetilde q$. The exact moments obey
\[
 \mu\in[\widetilde\mu-\gamma M,\widetilde\mu+\gamma M],\qquad
 q\in[\widetilde q-\gamma M^2,\widetilde q+\gamma M^2].
\]
The bound allows sequential summation, multiplication and division and is conservative for the implementation's pairwise sums. Rational arithmetic constructs these intervals from the exact binary64 outputs. There is no uncharged empirical-variance rounding. If $[\mu_-,\mu_+]$ is the mean enclosure, then
\[
 V^+=\frac{n}{n-1}\left[q_+-\min_{v\in[\mu_-,\mu_+]}v^2\right]_+
\]
is an upper bound for the unbiased sample variance. The minimum is zero if the mean interval contains zero and otherwise the smaller squared endpoint.

There are $K=120$ contrasts, each with two one-sided tails, and total error $\alpha=0.01$. The bound of \citet{maurer2009} uses $\log(4K/\alpha)=\log(48000)<11$. The inequality is verified by the exact positive sum $\sum_{j=0}^{40}11^j/j!>48000$. A valid endpoint radius is
\[
 R(V^+)=\sqrt{2V^+\,11/n}+\frac{7(b-a)\,11}{3(n-1)}.
\]
The code rounds the square-root result upward until its exact rational square is at least the radicand. All remaining additions are rational. Subtract this radius from the lower endpoint's mean lower bound and add it to the upper endpoint's mean upper bound. Intersect with known support, then round the final reported display outward. This proves simultaneous coverage with probability at least $0.99$ conditional on the fixed-policy, independent-bin model. A seed supplies reproducibility, not evidence that a pseudorandom generator satisfies a mathematical independence axiom.

\subsection{Decision margins and selection}
The protocol fixed the superiority threshold at zero and the separate decision tolerance at $1/1024$ before the new comparison. Thus a full interval inside $[-1/1024,1/1024]$ establishes two-sided cost proximity in the normalized units on the joint coverage event. It does not establish equality, a signed ranking, or a calibrated welfare equivalence. The theorem for choosing among a finite catalogue is a direct inequality on this same event: $J_{\widehat m}-J_b\le U_{\widehat m}$ and $U_m\le J_m-J_b+U_m-L_m$. Minimizing $U_m+k_m$ gives the stated oracle inequality. Charges are known inputs to that decision problem; converting observed seconds to an economic charge requires an additional model and is not performed here.

The complete cost intervals follow. The same fixed policy may appear under distinct target labels; its identity is retained in every row. A timing repetition is not a new trained policy. All lower and upper endpoint moments, supports, seeds, stream hashes, ambiguity counts and input checkpoint hashes are deposited with the result records.
\begin{longtable}{rrllcrr}\caption{All 120 direct expected-cost intervals}\label{tab:directfull47}\\\toprule
$T$ & $p$ & Target & Initial law & Contrast & Lower & Upper\\\midrule\endfirsthead
\toprule $T$ & $p$ & Target & Initial law & Contrast & Lower & Upper\\\midrule\endhead
\bottomrule\endfoot
2 & 1 & 1/4 & 1/2 & C--S & -0.0002078 & 0.0001255 \\
2 & 1 & 1/4 & 1/2 & W--C & -0.0000865 & 0.0000825 \\
2 & 1 & 1/4 & 1/2 & W--S & -0.0001976 & 0.0001113 \\
2 & 1 & 1/4 & 1/8 & C--S & -0.0001577 & 0.0001711 \\
2 & 1 & 1/4 & 1/8 & W--C & -0.0000918 & 0.0000883 \\
2 & 1 & 1/4 & 1/8 & W--S & -0.0001472 & 0.0001572 \\
2 & 1 & 1/4 & 7/8 & C--S & -0.0001637 & 0.0001701 \\
2 & 1 & 1/4 & 7/8 & W--C & -0.0000785 & 0.0000718 \\
2 & 1 & 1/4 & 7/8 & W--S & -0.0001548 & 0.0001546 \\
2 & 1 & 1/4 & uniform & C--S & -0.0001858 & 0.0001574 \\
2 & 1 & 1/4 & uniform & W--C & -0.0000891 & 0.0001048 \\
2 & 1 & 1/4 & uniform & W--S & -0.0001656 & 0.0001529 \\
2 & 1 & 1/8 & 1/2 & C--S & -0.0000776 & 0.0000893 \\
2 & 1 & 1/8 & 1/2 & W--C & -0.0000443 & 0.0000393 \\
2 & 1 & 1/8 & 1/2 & W--S & -0.0000740 & 0.0000806 \\
2 & 1 & 1/8 & 1/8 & C--S & -0.0000908 & 0.0000736 \\
2 & 1 & 1/8 & 1/8 & W--C & -0.0000444 & 0.0000476 \\
2 & 1 & 1/8 & 1/8 & W--S & -0.0000830 & 0.0000690 \\
2 & 1 & 1/8 & 7/8 & C--S & -0.0000840 & 0.0000830 \\
2 & 1 & 1/8 & 7/8 & W--C & -0.0000387 & 0.0000365 \\
2 & 1 & 1/8 & 7/8 & W--S & -0.0000790 & 0.0000757 \\
2 & 1 & 1/8 & uniform & C--S & -0.0000850 & 0.0000870 \\
2 & 1 & 1/8 & uniform & W--C & -0.0000518 & 0.0000452 \\
2 & 1 & 1/8 & uniform & W--S & -0.0000821 & 0.0000776 \\
2 & 1 & 1/16 & 1/2 & C--S & -0.0000391 & 0.0000443 \\
2 & 1 & 1/16 & 1/2 & W--C & -0.0000205 & 0.0000214 \\
2 & 1 & 1/16 & 1/2 & W--S & -0.0000356 & 0.0000417 \\
2 & 1 & 1/16 & 1/8 & C--S & -0.0000409 & 0.0000414 \\
2 & 1 & 1/16 & 1/8 & W--C & -0.0000230 & 0.0000226 \\
2 & 1 & 1/16 & 1/8 & W--S & -0.0000380 & 0.0000381 \\
2 & 1 & 1/16 & 7/8 & C--S & -0.0000459 & 0.0000377 \\
2 & 1 & 1/16 & 7/8 & W--C & -0.0000189 & 0.0000188 \\
2 & 1 & 1/16 & 7/8 & W--S & -0.0000429 & 0.0000346 \\
2 & 1 & 1/16 & uniform & C--S & -0.0000431 & 0.0000428 \\
2 & 1 & 1/16 & uniform & W--C & -0.0000280 & 0.0000207 \\
2 & 1 & 1/16 & uniform & W--S & -0.0000437 & 0.0000361 \\
2 & 4 & 1/4 & 1/2 & C--S & -0.0001589 & 0.0001743 \\
2 & 4 & 1/4 & 1/2 & W--C & -0.0000783 & 0.0000966 \\
2 & 4 & 1/4 & 1/2 & W--S & -0.0001367 & 0.0001704 \\
2 & 4 & 1/4 & 1/8 & C--S & -0.0001919 & 0.0001304 \\
2 & 4 & 1/4 & 1/8 & W--C & -0.0000923 & 0.0000953 \\
2 & 4 & 1/4 & 1/8 & W--S & -0.0001774 & 0.0001189 \\
2 & 4 & 1/4 & 7/8 & C--S & -0.0001889 & 0.0001489 \\
2 & 4 & 1/4 & 7/8 & W--C & -0.0000785 & 0.0000806 \\
2 & 4 & 1/4 & 7/8 & W--S & -0.0001750 & 0.0001370 \\
2 & 4 & 1/4 & uniform & C--S & -0.0001569 & 0.0001865 \\
2 & 4 & 1/4 & uniform & W--C & -0.0001162 & 0.0001049 \\
2 & 4 & 1/4 & uniform & W--S & -0.0001496 & 0.0001680 \\
2 & 4 & 1/8 & 1/2 & C--S & -0.0000783 & 0.0000884 \\
2 & 4 & 1/8 & 1/2 & W--C & -0.0000449 & 0.0000417 \\
2 & 4 & 1/8 & 1/2 & W--S & -0.0000735 & 0.0000804 \\
2 & 4 & 1/8 & 1/8 & C--S & -0.0000814 & 0.0000800 \\
2 & 4 & 1/8 & 1/8 & W--C & -0.0000458 & 0.0000491 \\
2 & 4 & 1/8 & 1/8 & W--S & -0.0000732 & 0.0000751 \\
2 & 4 & 1/8 & 7/8 & C--S & -0.0000899 & 0.0000793 \\
2 & 4 & 1/8 & 7/8 & W--C & -0.0000385 & 0.0000408 \\
2 & 4 & 1/8 & 7/8 & W--S & -0.0000822 & 0.0000739 \\
2 & 4 & 1/8 & uniform & C--S & -0.0000792 & 0.0000925 \\
2 & 4 & 1/8 & uniform & W--C & -0.0000657 & 0.0000453 \\
2 & 4 & 1/8 & uniform & W--S & -0.0000829 & 0.0000758 \\
2 & 4 & 1/16 & 1/2 & C--S & -0.0000386 & 0.0000450 \\
2 & 4 & 1/16 & 1/2 & W--C & -0.0000215 & 0.0000222 \\
2 & 4 & 1/16 & 1/2 & W--S & -0.0000350 & 0.0000421 \\
2 & 4 & 1/16 & 1/8 & C--S & -0.0000380 & 0.0000427 \\
2 & 4 & 1/16 & 1/8 & W--C & -0.0000233 & 0.0000241 \\
2 & 4 & 1/16 & 1/8 & W--S & -0.0000343 & 0.0000398 \\
2 & 4 & 1/16 & 7/8 & C--S & -0.0000451 & 0.0000395 \\
2 & 4 & 1/16 & 7/8 & W--C & -0.0000200 & 0.0000196 \\
2 & 4 & 1/16 & 7/8 & W--S & -0.0000421 & 0.0000360 \\
2 & 4 & 1/16 & uniform & C--S & -0.0000404 & 0.0000457 \\
2 & 4 & 1/16 & uniform & W--C & -0.0000319 & 0.0000236 \\
2 & 4 & 1/16 & uniform & W--S & -0.0000413 & 0.0000384 \\
4 & 1 & 1/4 & 1/2 & C--S & -0.0001610 & 0.0001274 \\
4 & 1 & 1/4 & 1/2 & W--C & -0.0000776 & 0.0000826 \\
4 & 1 & 1/4 & 1/2 & W--S & -0.0001471 & 0.0001183 \\
4 & 1 & 1/4 & 1/8 & C--S & -0.0001324 & 0.0001547 \\
4 & 1 & 1/4 & 1/8 & W--C & -0.0000808 & 0.0000937 \\
4 & 1 & 1/4 & 1/8 & W--S & -0.0001144 & 0.0001497 \\
4 & 1 & 1/4 & 7/8 & C--S & -0.0001620 & 0.0001263 \\
4 & 1 & 1/4 & 7/8 & W--C & -0.0000784 & 0.0000726 \\
4 & 1 & 1/4 & 7/8 & W--S & -0.0001535 & 0.0001120 \\
4 & 1 & 1/4 & uniform & C--S & -0.0001608 & 0.0001330 \\
4 & 1 & 1/4 & uniform & W--C & -0.0000918 & 0.0000835 \\
4 & 1 & 1/4 & uniform & W--S & -0.0001535 & 0.0001175 \\
4 & 1 & 1/8 & 1/2 & C--S & -0.0000640 & 0.0000806 \\
4 & 1 & 1/8 & 1/2 & W--C & -0.0000382 & 0.0000434 \\
4 & 1 & 1/8 & 1/2 & W--S & -0.0000556 & 0.0000774 \\
4 & 1 & 1/8 & 1/8 & C--S & -0.0000566 & 0.0000869 \\
4 & 1 & 1/8 & 1/8 & W--C & -0.0000395 & 0.0000476 \\
4 & 1 & 1/8 & 1/8 & W--S & -0.0000468 & 0.0000851 \\
4 & 1 & 1/8 & 7/8 & C--S & -0.0000724 & 0.0000719 \\
4 & 1 & 1/8 & 7/8 & W--C & -0.0000370 & 0.0000385 \\
4 & 1 & 1/8 & 7/8 & W--S & -0.0000660 & 0.0000669 \\
4 & 1 & 1/8 & uniform & C--S & -0.0000743 & 0.0000732 \\
4 & 1 & 1/8 & uniform & W--C & -0.0000441 & 0.0000435 \\
4 & 1 & 1/8 & uniform & W--S & -0.0000689 & 0.0000672 \\
4 & 4 & 1/4 & 1/2 & C--S & -0.0001615 & 0.0001292 \\
4 & 4 & 1/4 & 1/2 & W--C & -0.0000777 & 0.0000923 \\
4 & 4 & 1/4 & 1/2 & W--S & -0.0001419 & 0.0001241 \\
4 & 4 & 1/4 & 1/8 & C--S & -0.0001735 & 0.0001139 \\
4 & 4 & 1/4 & 1/8 & W--C & -0.0000906 & 0.0000909 \\
4 & 4 & 1/4 & 1/8 & W--S & -0.0001610 & 0.0001018 \\
4 & 4 & 1/4 & 7/8 & C--S & -0.0001368 & 0.0001565 \\
4 & 4 & 1/4 & 7/8 & W--C & -0.0000782 & 0.0000822 \\
4 & 4 & 1/4 & 7/8 & W--S & -0.0001225 & 0.0001462 \\
4 & 4 & 1/4 & uniform & C--S & -0.0001432 & 0.0001535 \\
4 & 4 & 1/4 & uniform & W--C & -0.0000963 & 0.0000992 \\
4 & 4 & 1/4 & uniform & W--S & -0.0001296 & 0.0001428 \\
4 & 4 & 1/8 & 1/2 & C--S & -0.0000647 & 0.0000808 \\
4 & 4 & 1/8 & 1/2 & W--C & -0.0000451 & 0.0000405 \\
4 & 4 & 1/8 & 1/2 & W--S & -0.0000609 & 0.0000724 \\
4 & 4 & 1/8 & 1/8 & C--S & -0.0000727 & 0.0000710 \\
4 & 4 & 1/8 & 1/8 & W--C & -0.0000444 & 0.0000464 \\
4 & 4 & 1/8 & 1/8 & W--S & -0.0000656 & 0.0000658 \\
4 & 4 & 1/8 & 7/8 & C--S & -0.0000731 & 0.0000734 \\
4 & 4 & 1/8 & 7/8 & W--C & -0.0000405 & 0.0000394 \\
4 & 4 & 1/8 & 7/8 & W--S & -0.0000675 & 0.0000667 \\
4 & 4 & 1/8 & uniform & C--S & -0.0000700 & 0.0000788 \\
4 & 4 & 1/8 & uniform & W--C & -0.0000483 & 0.0000494 \\
4 & 4 & 1/8 & uniform & W--S & -0.0000634 & 0.0000732 \\
\end{longtable}
W is the witness, C the same-critic nearest actor, and S the spline. Every displayed endpoint is rounded outward to seven decimal places; the exact rational intervals control all decisions and are deposited. All 120 intervals lie inside the predeclared band. Repeated target labels and policy hashes remain explicit in the underlying records.


\section{Complete frontiers and experiment records}\label{supp:frontiers47}
For each method, the first-crossing function is constant between adjacent values of its recorded certified bounds. Partitioning the positive tolerance axis by the union of both methods' bounds therefore gives the exact comparison everywhere, with left endpoints included because attainment is weak inequality. The full-prefix Bellman count charges every earlier failed rung. Above the largest breakpoint both methods use their first rung. An interval in which only one method attains is distinguished from an interval in which both attain with different query counts. This descriptive reconstruction is not a new experimental target design.
\begin{longtable}{rrlllr}\caption{Every interval of the scalar first-crossing curves}\label{tab:frontierfull47}\\\toprule
$T$ & $p$ & Tolerance lower & Tolerance upper & Relation & $N_{\rm W}/N_{\rm S}$\\\midrule\endfirsthead
\toprule $T$ & $p$ & Lower & Upper & Relation & $N_{\rm W}/N_{\rm S}$\\\midrule\endhead\bottomrule\endfoot
2 & 1 & 0.0000000 & 0.0343075 & Neither & --/-- \\
2 & 1 & 0.0343075 & 0.0396413 & W only & 512/-- \\
2 & 1 & 0.0396413 & 0.0686151 & Same & 512/512 \\
2 & 1 & 0.0686151 & 0.0791884 & W fewer & 256/512 \\
2 & 1 & 0.0791884 & 0.1372302 & Same & 256/256 \\
2 & 1 & 0.1372302 & 0.1579999 & W fewer & 128/256 \\
2 & 1 & 0.1579999 & 0.2744603 & Same & 128/128 \\
2 & 1 & 0.2744603 & 0.3144939 & W fewer & 64/128 \\
2 & 1 & 0.3144939 & 0.5489206 & Same & 64/64 \\
2 & 1 & 0.5489206 & 0.6229815 & W fewer & 32/64 \\
2 & 1 & 0.6229815 & 1.0978413 & Same & 32/32 \\
2 & 1 & 1.0978413 & 1.2220383 & W fewer & 16/32 \\
2 & 4 & 0.0000000 & 0.0350171 & Neither & --/-- \\
2 & 4 & 0.0350171 & 0.0409569 & W only & 512/-- \\
2 & 4 & 0.0409569 & 0.0700341 & Same & 512/512 \\
2 & 4 & 0.0700341 & 0.0818309 & W fewer & 256/512 \\
2 & 4 & 0.0818309 & 0.1400683 & Same & 256/256 \\
2 & 4 & 0.1400683 & 0.1633408 & W fewer & 128/256 \\
2 & 4 & 0.1633408 & 0.2801366 & Same & 128/128 \\
2 & 4 & 0.2801366 & 0.3253087 & W fewer & 64/128 \\
2 & 4 & 0.3253087 & 0.5602732 & Same & 64/64 \\
2 & 4 & 0.5602732 & 0.6454998 & W fewer & 32/64 \\
2 & 4 & 0.6454998 & 1.1205463 & Same & 32/32 \\
2 & 4 & 1.1205463 & 1.2693757 & W fewer & 16/32 \\
4 & 1 & 0.0000000 & 0.0643873 & Neither & --/-- \\
4 & 1 & 0.0643873 & 0.0663689 & W only & 512/-- \\
4 & 1 & 0.0663689 & 0.1287747 & Same & 512/512 \\
4 & 1 & 0.1287747 & 0.1325577 & W fewer & 256/512 \\
4 & 1 & 0.1325577 & 0.2575494 & Same & 256/256 \\
4 & 1 & 0.2575494 & 0.2643956 & W fewer & 128/256 \\
4 & 1 & 0.2643956 & 0.5150987 & Same & 128/128 \\
4 & 1 & 0.5150987 & 0.5259184 & W fewer & 64/128 \\
4 & 1 & 0.5259184 & 1.0301974 & Same & 64/64 \\
4 & 1 & 1.0301974 & 1.0403702 & W fewer & 32/64 \\
4 & 1 & 1.0403702 & 2.0351529 & Same & 32/32 \\
4 & 1 & 2.0351529 & 2.0603948 & S fewer & 32/16 \\
4 & 4 & 0.0000000 & 0.0657205 & Neither & --/-- \\
4 & 4 & 0.0657205 & 0.0699794 & W only & 512/-- \\
4 & 4 & 0.0699794 & 0.1314410 & Same & 512/512 \\
4 & 4 & 0.1314410 & 0.1398014 & W fewer & 256/512 \\
4 & 4 & 0.1398014 & 0.2628819 & Same & 256/256 \\
4 & 4 & 0.2628819 & 0.2789813 & W fewer & 128/256 \\
4 & 4 & 0.2789813 & 0.5257639 & Same & 128/128 \\
4 & 4 & 0.5257639 & 0.5554218 & W fewer & 64/128 \\
4 & 4 & 0.5554218 & 1.0515278 & Same & 64/64 \\
4 & 4 & 1.0515278 & 1.1008786 & W fewer & 32/64 \\
4 & 4 & 1.1008786 & 2.1030555 & Same & 32/32 \\
4 & 4 & 2.1030555 & 2.1614686 & W fewer & 16/32 \\
\end{longtable}
Lower endpoints are included and upper endpoints excluded in the exact records. Decimal breakpoint displays are descriptive rounded values, not substitute decision thresholds. The deposited rational endpoints determine the step functions without rounding ambiguity. Query counts equal the full prefix through the indicated rung; above each cell's last displayed upper endpoint both methods use their first rung.

\begin{longtable}{rrlrrrr}\caption{All distinct constrained construction rungs}\label{tab:constrainedfull47}\\\toprule
$T$ & $p$ & Method & $N$ & Bound & Prefix queries & Midpoints\\\midrule\endfirsthead
\toprule $T$ & $p$ & Method & $N$ & Bound & Prefix queries & Midpoints\\\midrule\endhead\bottomrule\endfoot
2 & 1 & Bilinear FVI & 4 & 10.12531 & 250 & 1,000 \\
2 & 1 & Bilinear FVI & 8 & 5.37034 & 1,708 & 12,664 \\
2 & 1 & Bilinear FVI & 16 & 2.79602 & 11,534 & 169,880 \\
2 & 1 & Witness & 4 & 9.25915 & 250 & 1,000 \\
2 & 1 & Witness & 8 & 4.62958 & 1,708 & 12,664 \\
2 & 1 & Witness & 16 & 2.31479 & 11,534 & 169,880 \\
2 & 4 & Bilinear FVI & 4 & 10.26707 & 250 & 1,000 \\
2 & 4 & Bilinear FVI & 8 & 5.43656 & 1,708 & 12,664 \\
2 & 4 & Bilinear FVI & 16 & 2.82885 & 11,534 & 169,880 \\
2 & 4 & Witness & 4 & 9.47650 & 250 & 1,000 \\
2 & 4 & Witness & 8 & 4.73825 & 1,708 & 12,664 \\
2 & 4 & Witness & 16 & 2.36912 & 11,534 & 169,880 \\
3 & 1 & Bilinear FVI & 4 & 14.09329 & 375 & 1,500 \\
3 & 1 & Bilinear FVI & 8 & 7.45424 & 2,562 & 18,996 \\
3 & 1 & Bilinear FVI & 16 & 3.89149 & 17,301 & 254,820 \\
3 & 1 & Witness & 4 & 13.28483 & 375 & 1,500 \\
3 & 1 & Witness & 8 & 6.64242 & 2,562 & 18,996 \\
3 & 1 & Witness & 16 & 3.32121 & 17,301 & 254,820 \\
3 & 4 & Bilinear FVI & 4 & 14.28140 & 375 & 1,500 \\
3 & 4 & Bilinear FVI & 8 & 7.53684 & 2,562 & 18,996 \\
3 & 4 & Bilinear FVI & 16 & 3.93210 & 17,301 & 254,820 \\
3 & 4 & Witness & 4 & 13.64263 & 375 & 1,500 \\
3 & 4 & Witness & 8 & 6.82132 & 2,562 & 18,996 \\
3 & 4 & Witness & 16 & 3.41066 & 17,301 & 254,820 \\
\end{longtable}
Each row has three isolated timing repetitions with identical policy and certificate checkpoint hashes. All 72 raw rungs remain deposited. Every state-dependent action net, repaired actor, continuation label and numerical allowance can be reconstructed from its checkpoint.


\section{Continuous tensor compilation and feasible witness identity}\label{supp:compiler48}
This section proves Theorem~\ref{M-thm:compiler48} of the article. We state its identities explicitly so that the numerical compiler can be checked without relying on an unproved interpolation rule. The original labels and original witness indices, not transformed-node actions, are the inputs.

\subsection{The line transform and its tie invariant}
For ordered line locations $s_0<\cdots<s_m$, attach to each location a pair $(q_j,a_j)$, where $q_j$ is a real cost and $a_j$ an original index. Order such pairs lexicographically. Addition of a scalar to a pair adds it only to the cost. Start with $u_j=(q_j,a_j)$ and perform
\begin{align}\label{eq:line48}
 u_j&\leftarrow\min_{\rm lex}\{u_j,u_{j-1}+L(s_j-s_{j-1})\},
       &&j=1,\ldots,m,\\
 u_j&\leftarrow\min_{\rm lex}\{u_j,u_{j+1}+L(s_{j+1}-s_j)\},
       &&j=m-1,\ldots,0.\notag
\end{align}
After the forward pass, induction gives $u_j=\min_{k\le j,\rm lex}(q_k+L(s_j-s_k),a_k)$. The backward pass makes $u_j$ the minimum over all original line inputs of $(q_k+L|s_j-s_k|,a_k)$. To see this directly, a site to the right can be reached by the monotone backward path; a site to the left has already been considered in the forward pass. Any path that reverses direction has weakly larger distance because $L\ge0$, so it cannot improve the cost. With $L=0$ the original smallest index propagates in both directions. Lexicographic addition and minimum preserve the tie information throughout.

For a tensor grid, define the coordinate operator $T_j$ by taking this minimum along coordinate $j$ with all other coordinates fixed. Expanding the finite minima shows
\[
 T_d\cdots T_1(y_i,i)(v)
 =\min_{i,\rm lex}(y_i+L\|v-x_i\|_1,i).
\]
Thus the final result is $(z_v,o_v)$ in~\eqref{M-eq:transform48}. Each grid edge is relaxed at most twice in each coordinate. The edge count is at most $dS$, and all updates can be made in place with one cost and one index per node. The separable distance-transform part is the classical algorithm; see \citet{felzenszwalb2012}. Original action owners are auxiliary information propagated by the same minimum operations, not independently recomputed choices.

\subsection{Off-grid extension and the original tie rule}
For any corner $v$ and any original site $i$, the triangle inequality gives
\[
 y_i+L\|x-x_i\|_1\le y_i+L\|v-x_i\|_1+L\|x-v\|_1.
\]
Minimizing first over $i$ proves the lower inequality in~\eqref{M-eq:corner48}. For the converse, fix a site attaining $f(x)$. If its $j$th coordinate is at or to the left of the cell, choose the left corner coordinate; if it is at or to the right, choose the right corner coordinate. A grid site cannot lie strictly between adjacent grid coordinates. When $x$ lies on the cell boundary, either admissible corner choice gives the same needed identity. Summing the coordinate identities gives
\[
 \|x-x_i\|_1=\|x-v\|_1+\|v-x_i\|_1.
\]
Using this site in $z_v$ proves the upper inequality and hence equality.

Now take the smallest attaining original index $i_*(x)$. For its associated corner $v$, equality in the preceding chain implies
$z_v=y_{i_*}+L\|v-x_{i_*}\|_1$.
If the original owner $o_v$ were smaller than $i_*$, then
\[
 y_{o_v}+L\|x-x_{o_v}\|_1
 \le z_v+L\|x-v\|_1=f(x).
\]
It would be an attaining original site with a smaller index, a contradiction. Therefore $o_v=i_*(x)$. Evaluating the original cones of all corner owners and taking their lexicographic minimum proves~\eqref{M-eq:owner48}. This argument covers inconsistent or dominated labels: it never uses $z_v=y_v$.

Attaching the original action $a_{o_v}$ to a transformed label is essential. Attaching $a_v$ need not preserve the policy. Under the established identity, both implementations select $i_*(x)$ and apply the same feasible repair $\Gamma_{i_*}(x)$. They consequently induce the same action at every state, the same controlled transition law and, by backward evaluation, the same policy value. The argument is not restricted to expected costs; identity of the actual policy preserves any well-defined common recursive evaluation. It does not equate two policies generated from different rounded label arrays.

\subsection{Finite arithmetic and uncertain queries}
All preprocessing inputs in the executable study are exact rational representations of the stored numerical labels, slopes and coordinate knots. Comparisons in~\eqref{eq:line48} are exact. The compiler stores the resulting rational $z_v$ and integer $o_v$. Its arithmetic-operation count does not suppress numerator/denominator growth: coefficient bit lengths and actual preprocessing time are recorded.

At an exact represented point, outward arithmetic encloses each term of~\eqref{M-eq:corner48}; taking the minimum lower endpoint and the minimum upper endpoint encloses the true minimum. At an uncertain point $x\in D$, a midpoint value plus and minus $L\sup_{x\in D}\|x-\widehat x\|_1$ is also a valid value enclosure. This is the continuation's Lipschitz property, not an actor regularity assertion.

For the actor, let $\mathcal C(D)$ be the finitely many tensor cells whose closures meet a query box. The previous proof implies that every exact winning original index at any point in that box belongs to $\{o_v:v\in\mathcal V(C),C\in\mathcal C(D)\}$. An interval lower bound above the smallest available upper bound rules out a candidate; all other candidates are retained. Their action hull includes the exact actor at every permitted state. Repair, downward action quantization, costs and transitions are then evaluated by interval extensions. This may enlarge the path enclosure, but cannot improve the statistical comparison by hiding a discontinuity. A degenerate ambiguous box is resolved by exact original-cone comparison with the same smallest-index tie rule.

The off-grid work is bounded by the corner count only for a point or one-cell query. A large uncertainty box can intersect many cells and requires their union. We do not use the point-query complexity theorem to claim dimension-independent or uncertainty-independent simulation work. The actual ambiguity and actor-query counts are included in the direct-comparison records.


\subsection{Work to implemented economic accuracy}
For Corollary~\ref{M-cor:work-accuracy48}, the deterministic primitive account gives $G\le C_{T,p}/N+2S_T\xi+2\beta^T\xi_T$. Under its allocations the first term is at most $\varepsilon/4$ and the two numerical terms sum to at most $\varepsilon/4$. The exact-selector acquisition addition is $A2^{-b}+K\Delta\le\varepsilon/2$, so the implemented loss is at most $\varepsilon$. A sufficient integer bit count can be found by increasing $b$ until the rational inequality holds; no unchecked logarithm is required. If observation has a fee, the augmented objective also pays $\lambda dbS_T$. One cannot send $b$ to infinity at fixed positive price and infer a vanishing augmented loss.

There are $(N+1)^2$ state nodes, $N+1$ actions at each node and $N$ innovation midpoints. With a compiled continuation each midpoint requires a fixed number of two-dimensional primitive and corner operations, while transforming each date's labels and storing its feasible actors requires $O(N^2)$ further algebraic work. Thus the full backward construction has $O(TN^4)$ arithmetic operations. The dense alternative multiplies the midpoint workload by a further $(N+1)^2$. The finite-bit actor applies the same compiled owner rule and a constant-dimensional capacity repair. For a dyadic ladder ending at $N$, $\sum_{j\ge0}(N/2^j)^q\le N^q/(1-2^{-q})$ for $q>0$; this charges all previous construction attempts. The storage of the entire ladder of checkpoints also has the same $O(TN^2)$ scalar order.

These are sufficient algebraic upper bounds, not lower bounds for conventional methods. A $p$-bit elementary operation is charged at its actual or separately bounded arithmetic cost, and exact rational owners carry their coefficient sizes. Statistical policy comparison is outside this construction bound and is charged as its own joint service. No physical fsync latency, unknown optimizer iteration count, or unrestricted uncertainty-box enumeration is hidden inside the arithmetic order.

\section{Nonuniform fitted-value comparison and multidimensional primitives}\label{supp:baselines48}
\subsection{A common all-state criterion on unequal coordinate grids}
Let $g_t=\mathcal T_t f_{t+1}$ have graph modulus $L_t$. At nodes $x_i$, numerical feasible-net minimization returns $y_i$ with
$-e_t\le y_i-g_t(x_i)\le D_t^Qk_t+e_t$ and original action $a_i$ with $Q_t(f_{t+1};x_i,a_i)\le y_i+e_t$.
For multilinear interpolation on a rectangle of side lengths $h_j^C$, its nonnegative barycentric weights have sum one and satisfy
\[
 \sum_i w_i(x)\|x-x_i\|_1
 =\sum_j\frac{2(x_j-x_j^-)(x_j^+-x_j)}{h_j^C}
 \le\frac12\sum_jh_j^C\le h_t,
\]
where $h_t$ is half the sum of the largest coordinate gaps. Hence the optimal-residual upper bound is $D_t^Qk_t+e_t+L_th_t$.

Deploy the original action at the nearest coordinatewise node, choosing the smaller original index on a tie, with the same state-feasible repair. On its closed nearest-node cell $D_i$, feasible objective transport gives
\[
 Q_t(f_{t+1};x,\Gamma_i(x))-f_t(x)
 \le e_t+y_i+L_t\|x-x_i\|_1-f_t(x).
\]
Intersect $D_i$ with tensor cells and split at $x_i$ in every coordinate. On each resulting rectangle, the expression after $e_t$ is multi-affine: the absolute value has fixed sign and $f_t$ is multilinear. A multi-affine function attains a maximum at a rectangle vertex, as repeated one-dimensional maximization shows. It therefore suffices to test the left midpoint, node, and right midpoint in each coordinate, at most $3^d$ vertices per node, using outward arithmetic. Denote this computed upper bound by $E_t$. The common one-sided policy criterion becomes
\begin{equation}\label{eq:fvi48}
 G_F=\sum_{t<T}B_t(D_t^Qk_t+2e_t+L_th_t+E_t)
        +B_T(2e_T+2L_gh_T).
\end{equation}
The witness criterion remains~\eqref{M-eq:graphgap46} with exact repair. Both compare against the original optimum over all admissible adapted policies. Neither treats a sampled residual as an all-state error.

The interpolant's $\ell^1$ Lipschitz constant is the largest absolute edge slope: each partial derivative is a convex combination of edge slopes, and the dual norm of $\ell^1$ is $\ell^\infty$. Slopes are computed exactly from rational knots and labels. This modulus, not a neural modulus borrowed from another generator, is passed to that method's previous-date target construction.

\subsection{The prespecified error-driven coordinate grid}
The adaptive conventional comparator is a nonuniform tensor FVI method, not a sparse grid or locally adaptive quadtree. At every date and resolution $N>4$, it computes a fresh $5^d$ pilot array using its own future continuation and the same $N$ feasible action fractions and $N$ innovation midpoints. At terminal date it evaluates the terminal primitive instead. These evaluations are charged even when a later grid repeats a pilot node.

For coordinate $j$, let $D_{j,k}$, $k=1,2,3$, be the maximum, across all other pilot coordinates, of the absolute second difference $|y_{k-1}-2y_k+y_{k+1}|$. Give coarse interval $[k/4,(k+1)/4]$, $k=0,\ldots,3$, weight
\[
 w_{j,k}=1+64\max\{D_{j,\max(1,k)},D_{j,\min(3,k+1)}\}.
\]
Start with the four coarse intervals. Repeatedly bisect the interval maximizing $w_{j,k}$ times its squared length, breaking ties at the leftmost interval, until there are $N$ intervals on that axis. All rankings and knots are rational. This positive curvature heuristic and its factor 64 were fixed in the scientific source before the full catalogue; no reported outcome selects them. The final cover uses the actual largest gaps. In particular, preferentially refining one area can worsen a global covering allowance elsewhere. The common certificate exposes that tradeoff rather than crediting the adaptive method with an assumed uniform radius.

\subsection{Cycle-coupled dimensions two through four}
The dimension stress study retains $x\in[0,1]^d$, $a\in[0,\kappa(x)]$, with
\[
 \kappa(x)=\frac18+\frac1{8d}\sum_i x_i,\qquad
 F_i(x,a,z)=\frac1{16}+\frac{x_i}2+\frac{x_{i+1}}8
 +\frac{x_i(1-x_{i+1})}{16}+b_i a+s_i z.
\]
Indices are cyclic; $b_i$ alternates $1/2,1/4$, and $s_i$ alternates $1,-1$. The shock is the original continuous uniform $z\in[-1/32,1/32]$. Write $\bar x=d^{-1}\sum_i x_i$. The stage state cost is
\[
 \frac2d\sum_i(x_i-5/8)^2+
 \frac1{4d}\sum_i(x_i-x_{i+1})^2+
 2(1/2-2\bar x)_+^2,
\]
with action cost $pa^2+4a^4$. The terminal state cost doubles the first coefficient. At $d=2$ these are exactly the original R47 primitives, including its action set and opposite innovation exposures.

For $d=2,3,4$, valid $\ell^1$ moduli are
\begin{equation}\label{eq:moduli48}
 C_x=\frac{13}{2d},\quad L_g=\frac9d,\quad
 F_x=\frac{11}{16},\quad F_a=\sum_i b_i,\quad
 C_a=\frac p2+\frac14,\quad K^A=\frac1{8d}.
\end{equation}
Indeed, $0\le a\le1/4$ and the displayed transition lies in $[1/32,29/32]$ by direct bounds, so the unit cube is invariant. Each Jacobian column in the state has absolute sum at most $9/16+1/8=11/16$. The action and capacity moduli follow by differentiation.

For the state-cost bound, in the active shortage region the stage derivative with respect to $x_j$ is
\[
 \frac4d(x_j-5/8)+\frac{2x_j-x_{j-1}-x_{j+1}}{2d}
 -\frac8d(1/2-2\bar x).
\]
At $d=2$, the repeated neighbor is counted twice. Its coefficients in the coordinates are nonnegative for $d\le4$, so its minimum on the active region is attained at zero and equals $-13/(2d)$. Its positive part is at most the corresponding derivative without the negative shortage term, whose upper bound is $5/(2d)$. On the inactive region the derivative lies between $-7/(2d)$ and $5/(2d)$. The terminal derivative replaces $4/d$ by $8/d$; the analogous bounds give minimum $-9/d$ and smaller absolute upper bounds elsewhere. Continuity across the shortage boundary yields the global moduli in~\eqref{eq:moduli48}. These constants are proved for the displayed dimensions, not silently extrapolated to a different family.

A feasible action net has points $j\kappa(x)/N$, $j=0,\ldots,N$, and covering radius at most $1/(8N)$. In dimension three a downward binary64 realization adds at most $2^{-50}$ to that radius; its stored actions remain feasible. Dimensions two and four and the study's dyadic knots make these fractions exactly representable. Every path integral uses midpoint queries plus a continuous-law remainder. Since $\|F(x,a,z)-F(x,a,z')\|_1=d|z-z'|$, averaging over $M$ equal shock cells gives
\[
 \left|\E f(F(x,a,Z))-\frac1M\sum_{m=1}^M f(F(x,a,z_m))\right|
 \le\frac{dL_f}{64M}.
\]
The denominator is obtained from the mean absolute displacement within a cell, one quarter of its length. Discounting, interval evaluation and label rounding are additional terms; the deposited query error includes all of them.

\section{Observation charges and direct cost intervals}\label{supp:decisions48}
\subsection{The sensor account and integer choice}
Each observed coordinate cell has width $2^{-b}$ and midpoint error at most $2^{-(b+1)}$, including the boundary cells. For the affine capacity, evaluating the lower endpoint of the observation box produces a feasible lower capacity. Its distance from the true capacity is at most $2K^Ar(b)$. Downward action rounding by $\Delta$ adds at most $\Delta$ to repair displacement. Therefore the acquisition addition at date $t$ is
$2L_tr(b)+D_t^Q(2K_t^Ar(b)+\Delta_t)$.
Summing discounted additions and adding the deterministic observation charge gives~\eqref{M-eq:priced48}. The same summation does not claim that a reference policy-loss bound is a policy-cost point estimate.

The integer precision comparisons follow from the successive difference $C-A2^{-(b+1)}$. If this difference vanishes, both $b$ and $b+1$ minimize locally, and monotonicity makes them global minimizers; the smaller is recorded. Every candidate precision, charge, allowance and augmented bound is stored as a rational number. Checking the two neighbors, or all candidates in the finite range, is an exact optimization certificate for this upper-bound objective.

\subsection{Continuous-law path enclosures}
For a uniform initial law, two independent forty-bit bins cover the initial coordinates. For a fixed initial vector, the coordinate inputs are exact. At each date a further independent bin encloses the common continuous innovation. The bin width is $2^{-40}$ for an initial coordinate and $2^{-44}$ for the shock interval of total length $1/16$.

For each controller separately, its state box determines all possible observation bins; exact midpoints of those bins form an observation enclosure. The exact actor at every possible observation is enclosed using original witness owners or the nearest-node FVI cells. Capacity repair uses the lowest capacity on each possible observation box. Its action hull includes all possible repaired and quantized actions. The natural interval extensions of cost and transition then propagate the enclosure to the next date. The compiler may accelerate these evaluations but never replaces a path box by an unchecked midpoint actor. Dependence between the two policies is arbitrary inside a paired path; their marginals retain the prescribed innovation law.

The running-cost support follows from the separate bounds $\sum_j(x_j-5/8)^2\le25/32$, $(x_1-x_2)^2/4\le1/4$, $2(1/2-x_1-x_2)_+^2\le1/2$, $pa^2\le p/16$, and $4a^4\le1/64$. Terminal cost is at most $37/16$. These inequalities prove~\eqref{M-eq:support48}; separate maxima need not occur at the same state. The wider support is retained as a conservative input, rather than estimated from the simulated sample.

\subsection{Moment and probability accounts}
For a bounded endpoint array $X_1,\ldots,X_n\in[a,b]$, write $M=\max(|a|,|b|)$, $u=2^{-53}$ and $\gamma=(n+4)u/(1-(n+4)u)$. The implementation encloses the computed mean by adding and subtracting $\gamma M$; it enlarges the computed mean square by $\gamma M^2$. These conservative bounds cover multiplication, summation and division in the finite arrays used here, with no overflow. If the mean enclosure is $[m_-,m_+]$ and the second-moment upper bound is $q_+$, an unbiased sample-variance upper bound is
\[
 s_+^2=\max\left\{0,\frac n{n-1}
 \left(q_+-\min_{m\in[m_-,m_+]}m^2\right)\right\}.
\]
The minimum square is zero exactly when the interval contains zero. All moment bounds are converted to rationals before computing the confidence allowance. A square-root upper endpoint is checked by exact squaring and rounded upward until valid. The rational partial sum $\sum_{j=0}^{40}10^j/j!>19200=4K/\alpha$ verifies the logarithmic allowance $\ell=10$.

The empirical-Bernstein inequality of \citet{maurer2009}, applied to the two bounded endpoint samples, gives~\eqref{M-eq:bern48}--\eqref{M-eq:confidence48}. Each of the $2K$ one-sided assertions receives failure probability $\alpha/(2K)$. The union bound requires no independence across policy-pair tasks, but each task's endpoint sample uses the stipulated independent-bin model. Stored PCG64 seeds and stream hashes give deterministic reproducibility only. The path enclosures are deterministic numerical statements; the confidence level is a separate sampling statement under that model.

For an installed FVI controller, changing to witness costs $k+D_{\mu,b}$. For an installed witness controller, changing to FVI costs $k-D_{\mu,b}$. The sign tests and no-recoupment decision in Theorem~\ref{M-thm:decision48} follow. They concern the chosen transaction fee and initial law. They neither establish welfare equivalence under an unmodelled calibration nor prove that both controllers are close to the optimum merely because they are close to each other.

\section{Full work records and preservation}\label{supp:work48}
For each completed service, let $g_r$ be its exact bound at rung $r$ and $w_r$ its accumulated clock through that rung's durable checkpoint. The first-crossing work at tolerance $\varepsilon$ is $w_{\min\{r:g_r\le\varepsilon\}}$, with no attainment when this set is empty. Sorting the union of all exact bound breakpoints yields a complete partition of the positive tolerance axis into left-closed, right-open intervals on which every method's first crossing is fixed. Bounds need not be monotone for this construction: the earliest qualifying rung is always used. The lowest interval with no attainment and the unbounded upper interval are retained.

Every interval in the machine-readable frontier contains each method's selected resolution, median/minimum/maximum prefix wall time, CPU time, component counts and durable checkpoint size. The three repetitions are deterministic constructions with separate process clocks; they are not independent economic tasks or neural initializations. The dimension stress cells have one execution and are labelled accordingly. Peak resident memory includes the interpreter and warm-up. The internal prefix starts at primitives after the fixed non-economic warm-up; the parent process records the startup-inclusive duration separately. CPU affinity and numerical-library threads are fixed, while CPU frequency is not controlled.

The service includes target generation, adaptive pilots, exact compilation, continuous-law integration and its remainder, actor extraction, certification, failed rungs, serialization and a fixed finite acquired-deployment checksum. That checksum is a measured deployment component, not an empirical substitute for the uniform acquisition theorem. An independent direct policy-cost experiment is a separate joint service. Its clock includes both frozen-policy loads and compilation, hash checks, common paths, repair, statistics and durable output. This pair-specific cost is reported separately; it is not assigned silently to one controller or called zero. Historical R46 scalar prefix-time curves are reconstructed from their original records without retiming or adding them to R48 clocks.

The scalar, original constrained and precision findings remain historical evidence. The new compiled construction is a separate prospective service, while the fixed-object compiler diagnostic retains the old labels and original witnesses. This separation prevents faster evaluation of a frozen object from being misreported as better optimizer reliability or a newly superior policy.

\section{Every direct constrained-policy interval}\label{supp:direct-table48}
Intervals concern witness minus FVI expected resource cost. The fee column records the two-direction no-recoupment decision at $k=1/64$. The familywise statement is conditional on the independent-bin model. Displayed endpoints are rounded outward.
\begin{table}[htbp]
\centering
\caption{Direct intervals: horizon 2, price 1}\label{tab:direct-full48-2-1}
\footnotesize
\begin{tabular}{lcrrlr}
\toprule
Initial law & Bits & Lower & Upper & Fee & Joint seconds\\
\midrule
uniform & exact & -0.000821 & 0.001028 & yes & 2.77 \\
uniform & 6 & -0.000738 & 0.001108 & yes & 2.71 \\
uniform & 10 & -0.000806 & 0.001043 & yes & 2.72 \\
1/8 & exact & -0.001247 & 0.000596 & yes & 2.64 \\
1/8 & 6 & -0.001113 & 0.000726 & yes & 2.61 \\
1/8 & 10 & -0.001220 & 0.000622 & yes & 2.62 \\
1/2 & exact & -0.000935 & 0.000901 & yes & 2.62 \\
1/2 & 6 & -0.000918 & 0.000918 & yes & 2.59 \\
1/2 & 10 & -0.000933 & 0.000903 & yes & 2.62 \\
7/8 & exact & 0.000869 & 0.002717 & yes & 2.64 \\
7/8 & 6 & 0.000868 & 0.002716 & yes & 2.62 \\
7/8 & 10 & 0.000869 & 0.002717 & yes & 2.64 \\
\bottomrule
\end{tabular}
\begin{minipage}{0.97\linewidth}\footnotesize\vspace{3pt}
Exact-state policies are unquantized. The finite-bit policies use robust capacity and downward action quantization. Joint seconds charge both policies, source loading, common simulation, interval statistics and durable output; they are not assigned to one method. The total-path numerical enclosure is distinct from sampling uncertainty.
\end{minipage}
\end{table}

\begin{table}[htbp]
\centering
\caption{Direct intervals: horizon 2, price 4}\label{tab:direct-full48-2-4}
\footnotesize
\begin{tabular}{lcrrlr}
\toprule
Initial law & Bits & Lower & Upper & Fee & Joint seconds\\
\midrule
uniform & exact & -0.000694 & 0.001280 & yes & 2.76 \\
uniform & 6 & -0.000687 & 0.001287 & yes & 2.73 \\
uniform & 10 & -0.000694 & 0.001280 & yes & 2.73 \\
1/8 & exact & -0.001142 & 0.000830 & yes & 2.62 \\
1/8 & 6 & -0.001014 & 0.000959 & yes & 2.63 \\
1/8 & 10 & -0.001151 & 0.000821 & yes & 2.62 \\
1/2 & exact & -0.000136 & 0.001840 & yes & 2.63 \\
1/2 & 6 & -0.000374 & 0.001602 & yes & 2.65 \\
1/2 & 10 & -0.000145 & 0.001832 & yes & 2.61 \\
7/8 & exact & -0.000982 & 0.000983 & yes & 2.62 \\
7/8 & 6 & -0.000982 & 0.000982 & yes & 2.62 \\
7/8 & 10 & -0.000982 & 0.000983 & yes & 2.64 \\
\bottomrule
\end{tabular}
\begin{minipage}{0.97\linewidth}\footnotesize\vspace{3pt}
Exact-state policies are unquantized. The finite-bit policies use robust capacity and downward action quantization. Joint seconds charge both policies, source loading, common simulation, interval statistics and durable output; they are not assigned to one method. The total-path numerical enclosure is distinct from sampling uncertainty.
\end{minipage}
\end{table}

\begin{table}[htbp]
\centering
\caption{Direct intervals: horizon 3, price 1}\label{tab:direct-full48-3-1}
\footnotesize
\begin{tabular}{lcrrlr}
\toprule
Initial law & Bits & Lower & Upper & Fee & Joint seconds\\
\midrule
uniform & exact & -0.001028 & 0.001281 & yes & 3.95 \\
uniform & 6 & -0.000912 & 0.001394 & yes & 3.87 \\
uniform & 10 & -0.001007 & 0.001302 & yes & 3.88 \\
1/8 & exact & -0.001501 & 0.000801 & yes & 3.79 \\
1/8 & 6 & -0.001364 & 0.000935 & yes & 3.76 \\
1/8 & 10 & -0.001455 & 0.000846 & yes & 3.73 \\
1/2 & exact & -0.001217 & 0.001079 & yes & 3.77 \\
1/2 & 6 & -0.001147 & 0.001147 & yes & 3.76 \\
1/2 & 10 & -0.001216 & 0.001080 & yes & 3.76 \\
7/8 & exact & 0.000322 & 0.002627 & yes & 3.78 \\
7/8 & 6 & 0.000223 & 0.002528 & yes & 3.78 \\
7/8 & 10 & 0.000334 & 0.002639 & yes & 3.79 \\
\bottomrule
\end{tabular}
\begin{minipage}{0.97\linewidth}\footnotesize\vspace{3pt}
Exact-state policies are unquantized. The finite-bit policies use robust capacity and downward action quantization. Joint seconds charge both policies, source loading, common simulation, interval statistics and durable output; they are not assigned to one method. The total-path numerical enclosure is distinct from sampling uncertainty.
\end{minipage}
\end{table}

\begin{table}[htbp]
\centering
\caption{Direct intervals: horizon 3, price 4}\label{tab:direct-full48-3-4}
\footnotesize
\begin{tabular}{lcrrlr}
\toprule
Initial law & Bits & Lower & Upper & Fee & Joint seconds\\
\midrule
uniform & exact & -0.000688 & 0.001806 & yes & 3.92 \\
uniform & 6 & -0.000709 & 0.001783 & yes & 3.90 \\
uniform & 10 & -0.000685 & 0.001808 & yes & 3.90 \\
1/8 & exact & -0.001356 & 0.001135 & yes & 3.78 \\
1/8 & 6 & -0.001342 & 0.001149 & yes & 3.78 \\
1/8 & 10 & -0.001365 & 0.001126 & yes & 3.77 \\
1/2 & exact & -0.001114 & 0.001373 & yes & 3.80 \\
1/2 & 6 & -0.001048 & 0.001440 & yes & 3.78 \\
1/2 & 10 & -0.001111 & 0.001376 & yes & 3.77 \\
7/8 & exact & -0.000916 & 0.001570 & yes & 3.83 \\
7/8 & 6 & -0.000918 & 0.001568 & yes & 3.80 \\
7/8 & 10 & -0.000916 & 0.001569 & yes & 3.79 \\
\bottomrule
\end{tabular}
\begin{minipage}{0.97\linewidth}\footnotesize\vspace{3pt}
Exact-state policies are unquantized. The finite-bit policies use robust capacity and downward action quantization. Joint seconds charge both policies, source loading, common simulation, interval statistics and durable output; they are not assigned to one method. The total-path numerical enclosure is distinct from sampling uncertainty.
\end{minipage}
\end{table}


\section{Complete first-certificate time partitions}\label{supp:frontier-table48}
Each row is a left-closed, right-open tolerance interval. At the exact rational left endpoint, each method uses its first qualifying checkpoint; a dash denotes nonattainment. The final interval extends to infinity. Displayed endpoints have six significant digits; exact boundaries and complete component counts are in the deposited machine-readable frontier. Times are median prefix seconds for two-state services and singleton prefix seconds for dimension stress cases. These partitions are retrospective descriptions of the frozen complete ladders, not new prospective targets.
\begin{table}[htbp]
\centering
\caption{Full tolerance partition: dimension 2, horizon 2, price 1}\label{tab:frontier-full48-2-2-1}
\footnotesize
\begin{tabular}{rrrrrrrr}
\toprule
$\varepsilon_L$ & $\varepsilon_R$ & $N_C$ & C (s) & $N_U$ & U (s) & $N_A$ & A (s)\\
\midrule
0 & 0.578697 & -- & -- & -- & -- & -- & -- \\
0.578697 & 0.718896 & 64 & 47.996 & -- & -- & -- & -- \\
0.718896 & 1.15739 & 64 & 47.996 & 64 & 36.303 & 64 & 36.684 \\
1.15739 & 1.42452 & 32 & 3.686 & 64 & 36.303 & 64 & 36.684 \\
1.42452 & 2.31479 & 32 & 3.686 & 32 & 2.924 & 32 & 3.052 \\
2.31479 & 2.79602 & 16 & 0.346 & 32 & 2.924 & 32 & 3.052 \\
2.79602 & 4.62958 & 16 & 0.346 & 16 & 0.318 & 16 & 0.363 \\
4.62958 & 5.37034 & 8 & 0.055 & 16 & 0.318 & 16 & 0.363 \\
5.37034 & 9.25915 & 8 & 0.055 & 8 & 0.077 & 8 & 0.092 \\
9.25915 & 10.1253 & 4 & 0.015 & 8 & 0.077 & 8 & 0.092 \\
10.1253 & $\infty$ & 4 & 0.015 & 4 & 0.030 & 4 & 0.030 \\
\bottomrule
\end{tabular}
\begin{minipage}{0.97\linewidth}\footnotesize\vspace{3pt}
The all-state loss certificate is the common stopping criterion. Prefix time includes all earlier failed rungs, own-future construction, integration, compilation or adaptive pilots, repair checks and checkpoint fsync. Independent policy-cost inference is a separate joint service. The adaptive-coordinate comparator is executed only in dimension two.
\end{minipage}
\end{table}

\begin{table}[htbp]
\centering
\caption{Full tolerance partition: dimension 2, horizon 2, price 4}\label{tab:frontier-full48-2-2-4}
\footnotesize
\begin{tabular}{rrrrrrrr}
\toprule
$\varepsilon_L$ & $\varepsilon_R$ & $N_C$ & C (s) & $N_U$ & U (s) & $N_A$ & A (s)\\
\midrule
0 & 0.592281 & -- & -- & -- & -- & -- & -- \\
0.592281 & 0.727139 & 64 & 48.070 & -- & -- & -- & -- \\
0.727139 & 1.18456 & 64 & 48.070 & 64 & 36.386 & 64 & 37.020 \\
1.18456 & 1.44098 & 32 & 3.671 & 64 & 36.386 & 64 & 37.020 \\
1.44098 & 2.36912 & 32 & 3.671 & 32 & 2.936 & 32 & 3.063 \\
2.36912 & 2.82885 & 16 & 0.344 & 32 & 2.936 & 32 & 3.063 \\
2.82885 & 4.73825 & 16 & 0.344 & 16 & 0.318 & 16 & 0.363 \\
4.73825 & 5.43656 & 8 & 0.055 & 16 & 0.318 & 16 & 0.363 \\
5.43656 & 9.4765 & 8 & 0.055 & 8 & 0.078 & 8 & 0.093 \\
9.4765 & 10.2671 & 4 & 0.015 & 8 & 0.078 & 8 & 0.093 \\
10.2671 & $\infty$ & 4 & 0.015 & 4 & 0.030 & 4 & 0.030 \\
\bottomrule
\end{tabular}
\begin{minipage}{0.97\linewidth}\footnotesize\vspace{3pt}
The all-state loss certificate is the common stopping criterion. Prefix time includes all earlier failed rungs, own-future construction, integration, compilation or adaptive pilots, repair checks and checkpoint fsync. Independent policy-cost inference is a separate joint service. The adaptive-coordinate comparator is executed only in dimension two.
\end{minipage}
\end{table}

\begin{table}[htbp]
\centering
\caption{Full tolerance partition: dimension 2, horizon 3, price 1}\label{tab:frontier-full48-2-3-1}
\footnotesize
\begin{tabular}{rrrrrrrr}
\toprule
$\varepsilon_L$ & $\varepsilon_R$ & $N_C$ & C (s) & $N_U$ & U (s) & $N_A$ & A (s)\\
\midrule
0 & 0.830302 & -- & -- & -- & -- & -- & -- \\
0.830302 & 1.00278 & 64 & 71.928 & -- & -- & -- & -- \\
1.00278 & 1.6606 & 64 & 71.928 & 64 & 54.643 & 64 & 55.672 \\
1.6606 & 1.98559 & 32 & 5.534 & 64 & 54.643 & 64 & 55.672 \\
1.98559 & 3.32121 & 32 & 5.534 & 32 & 4.366 & 32 & 4.560 \\
3.32121 & 3.89149 & 16 & 0.503 & 32 & 4.366 & 32 & 4.560 \\
3.89149 & 6.64242 & 16 & 0.503 & 16 & 0.473 & 16 & 0.533 \\
6.64242 & 7.45424 & 8 & 0.078 & 16 & 0.473 & 16 & 0.533 \\
7.45424 & 13.2848 & 8 & 0.078 & 8 & 0.111 & 8 & 0.133 \\
13.2848 & 14.0933 & 4 & 0.021 & 8 & 0.111 & 8 & 0.133 \\
14.0933 & $\infty$ & 4 & 0.021 & 4 & 0.043 & 4 & 0.043 \\
\bottomrule
\end{tabular}
\begin{minipage}{0.97\linewidth}\footnotesize\vspace{3pt}
The all-state loss certificate is the common stopping criterion. Prefix time includes all earlier failed rungs, own-future construction, integration, compilation or adaptive pilots, repair checks and checkpoint fsync. Independent policy-cost inference is a separate joint service. The adaptive-coordinate comparator is executed only in dimension two.
\end{minipage}
\end{table}

\begin{table}[htbp]
\centering
\caption{Full tolerance partition: dimension 2, horizon 3, price 4}\label{tab:frontier-full48-2-3-4}
\footnotesize
\begin{tabular}{rrrrrrrr}
\toprule
$\varepsilon_L$ & $\varepsilon_R$ & $N_C$ & C (s) & $N_U$ & U (s) & $N_A$ & A (s)\\
\midrule
0 & 0.852664 & -- & -- & -- & -- & -- & -- \\
0.852664 & 1.01307 & 64 & 71.958 & -- & -- & -- & -- \\
1.01307 & 1.70533 & 64 & 71.958 & 64 & 54.146 & 64 & 55.021 \\
1.70533 & 2.00604 & 32 & 5.498 & 64 & 54.146 & 64 & 55.021 \\
2.00604 & 3.41066 & 32 & 5.498 & 32 & 4.367 & 32 & 4.550 \\
3.41066 & 3.9321 & 16 & 0.499 & 32 & 4.367 & 32 & 4.550 \\
3.9321 & 6.82132 & 16 & 0.499 & 16 & 0.472 & 16 & 0.530 \\
6.82132 & 7.53684 & 8 & 0.077 & 16 & 0.472 & 16 & 0.530 \\
7.53684 & 13.6426 & 8 & 0.077 & 8 & 0.111 & 8 & 0.133 \\
13.6426 & 14.2814 & 4 & 0.021 & 8 & 0.111 & 8 & 0.133 \\
14.2814 & $\infty$ & 4 & 0.021 & 4 & 0.043 & 4 & 0.043 \\
\bottomrule
\end{tabular}
\begin{minipage}{0.97\linewidth}\footnotesize\vspace{3pt}
The all-state loss certificate is the common stopping criterion. Prefix time includes all earlier failed rungs, own-future construction, integration, compilation or adaptive pilots, repair checks and checkpoint fsync. Independent policy-cost inference is a separate joint service. The adaptive-coordinate comparator is executed only in dimension two.
\end{minipage}
\end{table}

\begin{table}[htbp]
\centering
\caption{Full tolerance partition: dimension 3, horizon 2, price 1}\label{tab:frontier-full48-3-2-1}
\footnotesize
\begin{tabular}{rrrrrrrr}
\toprule
$\varepsilon_L$ & $\varepsilon_R$ & $N_C$ & C (s) & $N_U$ & U (s) & $N_A$ & A (s)\\
\midrule
0 & 2.32986 & -- & -- & -- & -- & -- & -- \\
2.32986 & 2.81383 & 16 & 9.499 & -- & -- & -- & -- \\
2.81383 & 4.65973 & 16 & 9.499 & 16 & 6.671 & -- & -- \\
4.65973 & 5.43228 & 8 & 0.551 & 16 & 6.671 & -- & -- \\
5.43228 & 9.31945 & 8 & 0.551 & 8 & 0.539 & -- & -- \\
9.31945 & 10.3112 & 4 & 0.051 & 8 & 0.539 & -- & -- \\
10.3112 & $\infty$ & 4 & 0.051 & 4 & 0.107 & -- & -- \\
\bottomrule
\end{tabular}
\begin{minipage}{0.97\linewidth}\footnotesize\vspace{3pt}
The all-state loss certificate is the common stopping criterion. Prefix time includes all earlier failed rungs, own-future construction, integration, compilation or adaptive pilots, repair checks and checkpoint fsync. Independent policy-cost inference is a separate joint service. The adaptive-coordinate comparator is executed only in dimension two.
\end{minipage}
\end{table}

\begin{table}[htbp]
\centering
\caption{Full tolerance partition: dimension 3, horizon 3, price 1}\label{tab:frontier-full48-3-3-1}
\footnotesize
\begin{tabular}{rrrrrrrr}
\toprule
$\varepsilon_L$ & $\varepsilon_R$ & $N_C$ & C (s) & $N_U$ & U (s) & $N_A$ & A (s)\\
\midrule
0 & 3.34922 & -- & -- & -- & -- & -- & -- \\
3.34922 & 3.91479 & 16 & 13.900 & -- & -- & -- & -- \\
3.91479 & 6.69844 & 16 & 13.900 & 16 & 9.900 & -- & -- \\
6.69844 & 7.54254 & 8 & 0.801 & 16 & 9.900 & -- & -- \\
7.54254 & 13.3969 & 8 & 0.801 & 8 & 0.799 & -- & -- \\
13.3969 & 14.3472 & 4 & 0.072 & 8 & 0.799 & -- & -- \\
14.3472 & $\infty$ & 4 & 0.072 & 4 & 0.157 & -- & -- \\
\bottomrule
\end{tabular}
\begin{minipage}{0.97\linewidth}\footnotesize\vspace{3pt}
The all-state loss certificate is the common stopping criterion. Prefix time includes all earlier failed rungs, own-future construction, integration, compilation or adaptive pilots, repair checks and checkpoint fsync. Independent policy-cost inference is a separate joint service. The adaptive-coordinate comparator is executed only in dimension two.
\end{minipage}
\end{table}

\begin{table}[htbp]
\centering
\caption{Full tolerance partition: dimension 4, horizon 2, price 1}\label{tab:frontier-full48-4-2-1}
\footnotesize
\begin{tabular}{rrrrrrrr}
\toprule
$\varepsilon_L$ & $\varepsilon_R$ & $N_C$ & C (s) & $N_U$ & U (s) & $N_A$ & A (s)\\
\midrule
0 & 4.62958 & -- & -- & -- & -- & -- & -- \\
4.62958 & 5.45577 & 8 & 7.728 & -- & -- & -- & -- \\
5.45577 & 9.25915 & 8 & 7.728 & 8 & 7.362 & -- & -- \\
9.25915 & 10.3892 & 4 & 0.382 & 8 & 7.362 & -- & -- \\
10.3892 & 18.2484 & 4 & 0.382 & 4 & 0.945 & -- & -- \\
18.2484 & 18.5183 & 4 & 0.382 & 2 & 0.307 & -- & -- \\
18.5183 & $\infty$ & 2 & 0.043 & 2 & 0.307 & -- & -- \\
\bottomrule
\end{tabular}
\begin{minipage}{0.97\linewidth}\footnotesize\vspace{3pt}
The all-state loss certificate is the common stopping criterion. Prefix time includes all earlier failed rungs, own-future construction, integration, compilation or adaptive pilots, repair checks and checkpoint fsync. Independent policy-cost inference is a separate joint service. The adaptive-coordinate comparator is executed only in dimension two.
\end{minipage}
\end{table}

\begin{table}[htbp]
\centering
\caption{Full tolerance partition: dimension 4, horizon 3, price 1}\label{tab:frontier-full48-4-3-1}
\footnotesize
\begin{tabular}{rrrrrrrr}
\toprule
$\varepsilon_L$ & $\varepsilon_R$ & $N_C$ & C (s) & $N_U$ & U (s) & $N_A$ & A (s)\\
\midrule
0 & 6.64242 & -- & -- & -- & -- & -- & -- \\
6.64242 & 7.58853 & 8 & 11.544 & -- & -- & -- & -- \\
7.58853 & 13.2848 & 8 & 11.544 & 8 & 10.976 & -- & -- \\
13.2848 & 14.4743 & 4 & 0.549 & 8 & 10.976 & -- & -- \\
14.4743 & 25.2494 & 4 & 0.549 & 4 & 1.387 & -- & -- \\
25.2494 & 26.5697 & 4 & 0.549 & 2 & 0.454 & -- & -- \\
26.5697 & $\infty$ & 2 & 0.061 & 2 & 0.454 & -- & -- \\
\bottomrule
\end{tabular}
\begin{minipage}{0.97\linewidth}\footnotesize\vspace{3pt}
The all-state loss certificate is the common stopping criterion. Prefix time includes all earlier failed rungs, own-future construction, integration, compilation or adaptive pilots, repair checks and checkpoint fsync. Independent policy-cost inference is a separate joint service. The adaptive-coordinate comparator is executed only in dimension two.
\end{minipage}
\end{table}

\begin{table}[htbp]
\centering
\caption{Historical scalar frontier using its original complete-prefix clocks}\label{tab:historical-clock48}
\footnotesize
\begin{tabular}{rrrrrrrr}
\toprule
$T$ & $p$ & W earlier & S earlier & Tie & W only & S only & Neither\\
\midrule
2 & 1 & 5 & 12 & 0 & 1 & 0 & 1 \\
2 & 4 & 5 & 12 & 0 & 1 & 0 & 1 \\
4 & 1 & 4 & 13 & 0 & 1 & 0 & 1 \\
4 & 4 & 5 & 12 & 0 & 1 & 0 & 1 \\
\bottomrule
\end{tabular}
\begin{minipage}{0.97\linewidth}\footnotesize\vspace{3pt}
Counts are tolerance intervals, not probabilities or weighted economic tasks. W is the R46 witness and S its spline comparator. These are unaltered historical internal prefix clocks, not R48 timings; they are never added to new compiler or direct-comparison clocks. The exact partitions, selected checkpoints and all three repeated times are retained.
\end{minipage}
\end{table}


\clearpage
\section*{Current R52 development}
\section{Primitive domain and exact moduli}\label{supp:primitive49}
Use the family in main equations~\eqref{M-eq:investment49}--\eqref{M-eq:cost49}, with $d=2,3,4$. Capacity is between $1/8$ and $1/4$ and has $\ell^1$ modulus $1/(8d)$. Every component of the dynamics is at least $1/32$. A direct upper bound is $29/32$ for the larger action exposure and $27/32$ for the smaller one. Thus the cube is invariant without an action-dependent exit event. The numerical clipping to $[0,1]$ is inactive on the true primitive image; it does not authorize out-of-domain evaluation.

For $d\ge2$, the sums of the absolute entries in a state column of the transition derivative are bounded by $1/2+1/16+1/8=11/16$. The action derivative has $\ell^1$ norm $\sum_j\gamma_j$. The action derivative of current cost is $2pa+16a^3\le p/2+1/4$. These are bounds between feasible pairs, not a claim that the same action is feasible at both states.

On either shortage region the state gradient is affine. The derivative of the squared-target and imbalance terms in coordinate $j$, after multiplication by $d$, is $4(x_j-5/8)+(2x_j-x_{j-1}-x_{j+1})/2$. In the shortage region add $-8(1/2-2\sum x/d)$. The terminal target derivative replaces 4 by 8. A linear functional on each region reaches its extrema at region vertices. For the declared dimensions, these are cube vertices and vertices with all but at most one coordinate at zero or one and $\sum x=d/4$. Exact evaluation gives absolute derivative bounds $13/(2d)$ and $9/d$. The added regression verifies these vertices rationally. The proof does not assert that these convenient formulas remain valid for untested arbitrary $d$.

The state covering radius of a nonuniform tensor grid is half the sum of its maximum coordinate spacings. Each action net is generated from a downward enclosure of capacity. Its coverage is at most $1/(8K)+2^{-48}$ under the implemented binary64 construction. This last conservative allowance covers capacity enclosure, multiplication, division and the final downward displacement on the bounded action domain. Fixed spacing and valid interval arithmetic are checked; a negative action is clipped to zero. Downward quantization on deployment preserves feasibility because the capacity itself is enclosed from below over the complete acquisition cell.

\section{Numerical labels and the separated account}\label{supp:numerical49}
An objective query has lower and upper endpoints for its midpoint-sum expectation, plus the deterministic continuous-law integration remainder. Labels are rounded on the inherited dyadic lattice, and the largest distance to the objective endpoints supplies the query allowance $e_t$. Selecting the least rounded label and retaining its original action costs one allowance in the label comparison and another in the complete selected-policy account. It does not require a second copy of the action-cover term in main equation~\eqref{M-eq:constructionloss49}.

For a uniform scalar bin of width $w$, the mean absolute distance to its midpoint is $w/4$. Here $w=1/(16M)$ and the transition's shock modulus is $d$. Therefore the numerical account decomposes into $\beta dL_{t+1}/(64M)$ and the remaining arithmetic width. The recorded primitive coefficients in main equation~\eqref{M-eq:coefficients49} reconstruct each witness bound as $a/N+b/K+q/M+\eta$ with a nonnegative numerical remainder. The release checks this decomposition at all witness rungs, as well as the full signed formula for every conventional rung.

The continuous allocation in main Proposition~\ref{M-prop:allocation49} minimizes a leading product, not a full finite-precision program. With positive coefficient $c_i$ and exponent $p_i$, changing variables to error contribution $x_i=c_i/r_i$ gives the objective $\sum_i p_i\log x_i$. Its constrained first-order condition is $x_i=\min(\bar x_i,p_i\lambda)$. The multiplier exists uniquely until every cap binds. A zero coefficient has no benefit from refinement, so its coordinate stays at the resource lower bound. Upward dyadic rounding keeps every error term weakly smaller. Compiler setup, integer counts, coefficient bit lengths, adaptive pilots, memory and durable output are not removed by that relaxation.

\section{Why the compiled actor is the original actor}\label{supp:owner49}
At a tensor node the distance transform minimizes over original sites, not merely neighboring transformed nodes. Each forward relaxation along a coordinate line compares the current value--owner pair with the preceding pair plus the weighted coordinate distance; the backward pass does the same from the opposite direction. The invariant is the lexicographic minimum over all processed original sites in the processed coordinates. Separability proves the invariant after the final coordinate. Keeping only the winning value would lose the actor identity.

For an off-grid point, choose a cell corner coordinatewise between the point and its lexicographically first minimizing site. Distances add exactly. At that corner an owner giving a smaller value, or a tied value and smaller original index, contradicts the minimizer at the original point. This proves the owner identity even for inconsistent labels, zero slope, boundary cells and ties. Nonuniform coordinate spacings enter only the exact distances.

In interval execution, value expressions are enclosed outward. Actor candidates whose lower objective exceeds another candidate's upper objective can be discarded; the remaining possibilities must be retained. An exact point tie is resolved by rational comparison and the original index rule. A genuine uncertain state box is not an exact point: the hull of its surviving actions is propagated. Boxes crossing cells use the union of all relevant corners. These operations may cost more than one point query, and the direct records retain actor and measurement ambiguity counts.

\section{Acquisition for witness and conventional policies}\label{supp:acquisition49}
Let $x\in C$ be the true state and $\widehat x$ the reported midpoint, with radius $r$ in $\ell^1$. Fix the selected original node $x_i$ and node action $a_i$. The ideal repair at the true state is feasible and within $K^A\|x-x_i\|_1$ of $a_i$. Replacing that repair by the cell-wide lower-capacity repair followed by downward spacing $\Delta$ changes the action by at most $2K^Ar+\Delta$. Thus
\[
 Q(x,a^{\rm impl})\le y_i+e+L^{\rm gr}\|x-x_i\|_1+ D(2K^Ar+\Delta).
\]
For a witness selection at $\widehat x$, $y_i+L^{\rm gr}\|\widehat x-x_i\|_1=f(\widehat x)$, up to a separately certified selector error. The triangle inequality and $L^f=L^{\rm gr}$ yield the allowance in main equation~\eqref{M-eq:acquire49}.

For conventional nearest-node selection, subtract and add $f(\widehat x)$. The expression $y_i+L^{\rm gr}\|\widehat x-x_i\|_1-f(\widehat x)$ is bounded by its complete nearest-cell maximum $E_t$. The remaining terms are at most $(L^{\rm gr}+L^f+2DK^A)r+D\Delta$. This proves main equation~\eqref{M-eq:fvisensor49} without a continuity assumption on the actor. It also explains why applying the witness coefficient unchanged to an interpolated conventional critic would be incorrect when the actual critic modulus differs.

For two-state records, $r(b)=2^{-b}$ and the common spacing is $1/4096$. The all-state bit criterion is therefore $G+K_\Delta+A2^{-b}+2\lambda b\sum_{t<T}\beta^t$, where $A$ uses the appropriate actual critic modulus and $K_\Delta=\sum_t\beta^tD_t/4096$. The interval-net criterion uses actual expected costs on the same finite grid instead of replacing $G+K_\Delta+A2^{-b}$ by a claim about true loss.

\section{Continuous path law and simultaneous inference}\label{supp:inference49}
Conditional on all construction records, the studied controllers are fixed. Each simulation row draws independent uniform indices for the initial coordinates and each innovation. The interval simulator encloses the complete continuous cell corresponding to those indices; no bin midpoint is substituted for the true path distribution in the estimand. Uniformly mixing conditional continuous cells reproduces the original uniform initial and innovation laws. Common innovations may couple two policies, but each follows its own endogenous trajectory.

In dimension two, a valid support interval for a raw discounted path-cost difference is $[-H,H]$, where
\[
 H=\frac{99+4p}{64}\sum_{t<T}\beta^t+\frac{37}{16}\beta^T.
\]
Indeed, current cost is nonnegative and at most $(99+4p)/64$ by separately bounding squared-target, imbalance, shortage and action costs. Terminal cost is at most $37/16$. These conservative bounds apply to every feasible repaired path. Clipping interval endpoints to the support preserves inclusion. Lower endpoint expectations lie below the true cost difference; upper endpoint expectations lie above it.

The stored moments consist of outward sample-mean intervals and upper sample-variance bounds for each endpoint sample. In \eqref{M-eq:radius49}, the square root is rounded upward after checking its square against the rational radicand. The linear remainder is rational. With $K\le300$, $\alpha=1/100$ and $\ell=13$, a finite positive Taylor sum verifies $e^{13}>4K/\alpha$. The bounded empirical-Bernstein inequalities of \citet{maurer2009} and a union bound therefore yield the claimed simultaneous event. Dependence between contrasts through shared economic objects is permitted; independence within each bin-index sample is the statistical assumption. The fixed pseudorandom seeds do not prove that assumption.

A first-crossing choice is measurable with respect to construction information. The conditional coverage assertion consequently remains valid after this choice. A bit selector or a fee queried after seeing the simultaneous intervals also needs no additional error allocation: on the one event, each underlying cost difference already belongs to its interval. This does not validate new initial laws, new trained policies, or unrecorded contrasts. Old and new study families have separate error accounts; combining them incurs the corresponding union bound.

\section{Exact frontiers and decision semantics}\label{supp:frontiers49}
For a fixed construction ladder, the first-crossing function can change only at a recorded certificate. A sorted union of all such values, together with zero, partitions all positive tolerances into left-closed, right-open intervals. At an attained endpoint the weak inequality in the stopping rule applies. Later nonmonotone or dominated certificates need not change the function, but remain in the raw record. Adjacent intervals with identical returned rungs may be merged for display; the exact rational partition is deposited.

For a direct interval $[l,u]$, the gain from replacing witness by FVI is in $[l,u]$; the reverse gain is in $[-u,-l]$. The fee $\max(0,u,-l)$ suffices to certify that neither switch is strictly profitable. The smaller $\max(0,l,-u)$ only says that neither switch has an interval lying wholly in its strictly profitable region. When zero lies inside a wide interval, this latter threshold is zero and is not evidence that a free switch cannot pay. The release reports both objects.

For information price $\lambda$, each bit choice supplies lower and upper affine net-cost lines. Enumerating all their nonnegative intersections partitions the price half-line. On an open interval, the active lower and upper lines are fixed. At a boundary all lines are compared exactly with smallest-bit ties. The difference between the active upper and lower lines is the selected implementation's valid regret upper bound. The identity line at 16 bits is exactly zero before information charges. No simulation count or time is assigned to an unperformed identity experiment; its arithmetic identity is recorded explicitly.

\section{Reference evaluation and safe improvement}\label{supp:improve50}
We give details of the reference-policy argument to separate a true improvement theorem from an approximate-greedy heuristic. All comparisons are on the same invariant domain and common conditional-evaluation domain as the main article. Feasible measurable policies include acquisition, capacity repair and numerical action realization.

Let $e_t=J_t^\pi-h_t$. At the terminal date $e_T\in[a_T,b_T]$. If $e_{t+1}\in[a_{t+1},b_{t+1}]$, cash invariance and monotonicity yield
\[
 \beta_ta_{t+1}\le
 \mathcal T_t^\pi J_{t+1}^\pi-\mathcal T_t^\pi h_{t+1}
 \le\beta_tb_{t+1}.
\]
Adding the signed own-policy residual proves the backward enclosure in Proposition~\ref{M-prop:advantage50}. Applying the same statement with either feasible action proves the difference allowance. No independence or differentiability of the fitted critic is required. The domain hypothesis matters: a certainty equivalent need not satisfy the argument for an inadmissible translated continuation.

If $\widetilde h_t=h_t+k_t$, then $\widetilde e_t=e_t-k_t$ and the residual shifts by $\beta_tk_{t+1}-k_t$. The propagated intervals become $[a_t-k_t,b_t-k_t]$, so their widths are unchanged. The two action evaluations of $\widetilde h_{t+1}$ each shift by $\beta_tk_{t+1}$, leaving the contrast unchanged. This is the centered invariance of the improvement test.

For the full expectation identity, define
$A_t^\pi(x,a)=Q_t(J_{t+1}^\pi;x,a)-J_t^\pi(x)$. Along any feasible $\mu$ trajectory,
\begin{align*}
 \E_x^\mu\sum_{t<T}B_tA_t^\pi(X_t,\mu_t(X_t))
 &=\E_x^\mu\sum_{t<T}\{B_tc_t(X_t,\mu_t(X_t))
          +B_{t+1}J_{t+1}^\pi(X_{t+1})-B_tJ_t^\pi(X_t)\}\\
 &=J_0^\mu(x)-J_0^\pi(x).
\end{align*}
Boundedness and finite horizon justify expectation and summation. Rejected actions give exactly zero reference advantage, while accepted ones give at most $-\kappa_t$. This proves Theorem~\ref{M-thm:improve50}. Under a nonlinear conditional recursion, use backward monotonicity instead of this telescoping expectation identity. A nonlinear evaluation does not justify the occupancy formula by changing notation alone.

The state-cell version fixes an action and verifies its inequality uniformly on the whole cell. Testing a cell midpoint alone is insufficient. Every implemented true state belongs to a cell whose gate is valid, and boundary cells use a declared deterministic report rule. In interval simulation, boxes straddling report boundaries retain all eligible actions, independently of whether their point-state selection is discontinuous. The analytic nonincrease may tighten a numerical advantage enclosure by intersection with $(-\infty,0]$, because the support is established independently of the simulated sample. This is not truncation of an unfavorable sample estimate.

\section{Exact integration and action optimization}\label{supp:terminal50}
Write $Z\sim U[-1/32,1/32]$. The deterministic part of each next coordinate is $y_j$, and its random part is $s_jZ$. Hence the expectation of $(y_j+s_jZ-c)^2$ is $(y_j-c)^2+1/3072$. For the cyclic difference use coefficient $s_j-s_{j+1}$. Summing yields the quadratic and variance terms in equation~\eqref{M-eq:exactterminal50}.

For the shortage term, the affine random component is uniform on $[-r,r]$, where $r=|\sum_js_j|/(16d)$. Direct antiderivatives give, for $r>0$,
\[
 \E(b+U)_+^2=\frac{(b+r)_+^3-(b-r)_+^3}{6r},\qquad
 \E(b+U)_+=\frac{(b+r)_+^2-(b-r)_+^2}{4r}.
\]
For $r=0$ take $b_+^2$ and $b_+$. These formulas hold at the two boundaries as well as inside the crossing region. The implementation encloses cancellation outward. Independent regression checks split the innovation interval at the exact shortage root and integrate the resulting quadratic pieces by exact rational Simpson identities; the checks are not used to generate the study's continuation labels.

At fixed $x$, the transition is affine in $a$. The terminal objective is a sum of convex quadratics and squared positive parts of affine functions. Its expectation is convex. The action cost $pa^2+4a^4$ has second derivative $2p+48a^2\ge2p$. Thus final conditional cost has a unique continuous minimizer, possibly on the capacity boundary. Its lattice minimizer lies at the first lattice point with nonnegative derivative or the preceding point, with boundary truncation. The implementation searches at most the fixed 4096 denominator's range, compares the two costs, and uses exact rational comparison when outward point intervals fail to decide a tie. The incumbent need not be one of these two actions: safety is established by the subsequent whole-cell comparison.

For bits $b$, let integer cell indices be $k_1,\ldots,k_d$. The smallest capacity on the cell is $1/8+\sum_jk_j/(8d2^b)$. Its largest lattice index is exactly
\[
 \left\lfloor\frac{4096(d2^b+\sum_jk_j)}{8d2^b}\right\rfloor.
\]
Integer arithmetic avoids an upward feasibility error, including in dimension three. The original action is clipped and rounded downward by this same rule. Feasibility is therefore not inferred from an interval midpoint. The final gate keeps the same earlier policies; its true improvement is relative to that acquired incumbent, not to a different exact-state policy.

\section{Simultaneous actual-cost inference}\label{supp:inference50}
The R50 primary family contains 25 pairs of baseline generators. Each group contains four policies---the two incumbents and their final repairs---and produces four absolute costs and six pairwise differences. Thus there are 250 sampled estimands. Their selection uses frozen construction information only. Within each group the continuous initial law and the first $T-1$ innovations are represented by forty-bit bins with outward propagation. The last innovation is integrated analytically. Conditional expectation over this last innovation preserves each policy's expected total cost. It reduces simulation variance without substituting a new innovation law.

The path-cost support derived in the retained inference section applies in dimensions two through four: the target and imbalance terms are coordinate averages with the same bounds, the shortage term is bounded by $1/2$, and feasible actions do not exceed $1/4$. Consequently $0\le J\text{-score}\le H$, with $H$ as in Section~\ref{supp:inference49}; a paired score lies in $[-H,H]$. An incumbent-minus-own-repair contrast is nonnegative by the independently proved gate and is bounded by $H$. Because the earlier trajectory is common, its numerical enclosure uses the centered final contrast instead of subtracting two complete total-cost intervals.

For endpoint observations $X_i\in[A,B]$, the retained empirical-Bernstein bound uses
\[
 R_n=\sqrt{2\overline{\operatorname{Var}}(X)\ell/n}
       +7(B-A)\ell/[3(n-1)],
\]
with outward moment intervals, an upper sample variance, and a rationally checked upward square root. Lower endpoints and upper endpoints each receive a one-sided bound. Set $n=131072$, $\ell=13$ and $\alpha=1/100$. The family cap is 512; an exact finite Taylor sum proves $e^{13}>4(512)/\alpha$. A union bound therefore covers all 250 actual-cost estimands under the independent-bin model. The cap leaves room but does not certify unrecorded models or initial laws. The pseudorandom stream and its hash establish reproducibility, not independence of ideal observations. The confidence event for the retained R49 family is separate; combining old and new statistical assertions incurs the sum of their failure probabilities.

\section{Mechanism and economic frontiers}\label{supp:mechanism50}
Denote witness and FVI by $W,F$ and their final repairs by $W^+,F^+$. Pure algebra yields
\begin{equation}\label{eq:decomposition50}
 J_W-J_F=(J_W-J_{W^+})+(J_{W^+}-J_{F^+})-(J_F-J_{F^+}).
\end{equation}
All three terms concern actual costs. The first and last are direct measures of removable final-action loss; the middle retains differences from earlier actions, acquired states, continuation approximation and the gate's conservatism. It is not identified as a single causal contribution of neural representation. The experiment reports each term, accepted and blocked cell counts, final action means, final state means and earlier action means by date. This exposes how much of the original ordering can be removed without retraining the continuation. It does not claim that the residual has been exhaustively decomposed into mutually independent mechanisms.

For finite intervals $[L_j,U_j]$, the upper-net-cost choice's true regret is bounded by the selected upper line minus the minimum lower line. To see why this is conservative, consider two overlapping cost intervals: their midpoint ordering may reverse within the same event. The upper-envelope decision remains valid but cannot identify the true minimizer. Pairwise centered intervals can certify individual replacements more sharply than the absolute frontier. For any fixed installation fee, transmission conversion factor and nonnegative resource price, add their declared costs before applying the comparison. A new price does not require new sampling on the same finite event. A new sensing policy does.

\section{Proofs for Finite-Sweep Accuracy}\label{app:finite-sweeps51}
\begin{proof}[Proof of Theorem~\ref{M-thm:finite-sweeps51}]
Fix a pass, date, and state, and suppress their indices. Write $J$ for the
incumbent continuation, $h$ for its approximator, $\pi$ for the incumbent
action, and $v$ for the proposal. By the signed evaluation band, for every
feasible action $a$,
\[
 \beta a_{t+1}\leq Q(J;x,a)-Q(h;x,a)\leq\beta b_{t+1}.
\]
Consequently, the incumbent-relative advantage of the proposal is at most
$\overline d+\beta w_{t+1}$. An accepted proposal therefore does not raise
the incumbent-continuation Bellman cost, while rejection retains exactly that
cost. Backward induction gives $J^{\pi^{k+1}}\leq J^{\pi^k}$ at every
date. This induction uses the new continuation for the implemented new
policy; the comparison itself uses the old continuation.

For an accepted proposal, approximate greediness and the band imply
\[
 Q(J;x,v)\leq (T_th)(x)+\alpha+\beta b_{t+1}
 \leq (T_tJ)(x)+\alpha+\beta w_{t+1}.
\]
For a rejected proposal, $\overline d+\beta w_{t+1}>0$ and
$\overline d\leq d+\zeta$. Thus
\[
 Q(h;x,\pi)<Q(h;x,v)+\zeta+\beta w_{t+1}
 \leq(T_th)(x)+\alpha+\zeta+\beta w_{t+1}.
\]
Transferring this inequality from $h$ to $J$ adds at most one further
$\beta w_{t+1}$. Both cases therefore yield
\begin{equation}\label{eq:approx-greedy-old51}
 T_t^{\pi^{k+1}}J_{t+1}^{\pi^k}
 \leq T_tJ_{t+1}^{\pi^k}+\epsilon_t^k.
\end{equation}
The already proved policy nonincrease, monotonicity of
$T_t^{\pi^{k+1}}$, and \eqref{eq:approx-greedy-old51} give
\[
 J_t^{\pi^{k+1}}\leq T_t^{\pi^{k+1}}J_{t+1}^{\pi^k}
 \leq T_tJ_{t+1}^{\pi^k}+\epsilon_t^k.
\]
The Bellman operator is $\beta_t$-Lipschitz in the supremum norm. Subtracting
$V_t^*=T_tV_{t+1}^*$ and using feasibility proves
\eqref{M-eq:sweep-recursion51}. Repeated substitution proves
\eqref{M-eq:sweep-unroll51}; at $m=T-t$ its terminal term vanishes.
Nonincrease of $E_t^k$ in further passes then gives
\eqref{M-eq:sweep-final51}. The exact case has zero error at every date.
\end{proof}

\begin{proof}[Proof of Proposition~\ref{M-prop:cell-greedy51}]
Fix $x\in C$ and any feasible $a\in A_t(x)$. Choose $m\in M_C$ with
$\|a-m\|\leq\delta_{A,C}$. Then
\begin{align*}
 Q_t(h;x,v_C)&\leq Q_t(h;x_C,v_C)+L_tr_C\\
 &\leq Q_t(h;x_C,m)+\eta_C+L_tr_C\\
 &\leq Q_t(h;x,a)+\eta_C+2L_tr_C+D_t\delta_{A,C}.
\end{align*}
Taking an infimum over $a$ proves the assertion without assuming that the
infimum is attained. Menu feasibility is uniform on $C$ by hypothesis.
\end{proof}

\begin{proof}[Proof of Proposition~\ref{M-prop:sharp-width51}]
Take $T=1$, zero stage cost, $\beta_0=1$, and two deterministic actions
leading to different terminal states $y_0,y_1$. Let $w>0$ and
$0<\delta<w$. Set
\[
 h_1(y_0)=w,\quad h_1(y_1)=\delta,\qquad
 g(y_0)=2w,\quad g(y_1)=\delta.
\]
The exact terminal error lies in $[0,w]$. The incumbent chooses $y_0$ and
the exact $h_1$-greedy proposal chooses $y_1$. Its exact contrast is
$d=\delta-w$, with $\alpha=\zeta=0$. The safe upper comparison is
$d+w=\delta>0$, so the gate retains the incumbent. Its actual optimality
gap is $2w-\delta$. Sending $\delta$ down to zero establishes sharpness.
This example also shows why rejecting an uncertified proposal is not evidence
that the proposal lacks a true economic benefit.
\end{proof}

\paragraph{Computed verification and logical scope.}
The source-bound exact-rational tests use finite transition matrices, signed
residual propagation, and independently evaluated policy and optimal costs.
They verify policy nonincrease, both gate cases, the one-pass recurrence,
constant-shift invariance, finite-horizon exact termination, and the two-action
sharpness construction. They are regression tests for the implementation of
the displayed formulas; the preceding arguments, rather than enumeration,
prove the theorems for the stated model class. When bands are statistical,
all passes and proposed comparisons require a common simultaneous coverage
event or a valid sequential error allocation. Reusing a confidence interval
outside its stated family does not satisfy the hypotheses.

\section{Proofs for policy-relevant continuation errors}\label{supp:contrast52}
\begin{proof}[Proof of Theorem~\ref{M-thm:contrast52}]
Fix a pass and suppress its index. For any feasible $a,b$, subtraction of the two Bellman expressions gives
\[
 Q_t(J_{t+1}^\pi;x,a)-Q_t(J_{t+1}^\pi;x,b)
 =Q_t(h_{t+1};x,a)-Q_t(h_{t+1};x,b)
       +\beta_t(P_t^a-P_t^b)e_{t+1}(x).
\]
Thus the accepted proposal has nonpositive true advantage relative to the incumbent. At a rejected state the implemented action is exactly the incumbent, whose advantage is zero. Backward induction gives $J^{\pi^+}\le J^\pi$. This proves safety even when only the proposal--incumbent contrast is certified.

For accuracy, suppose first that the proposal is accepted. For every feasible $a$, its approximate greediness and the uniform pair certificate imply
\[
 Q_t(J_{t+1}^\pi;x,v)-Q_t(J_{t+1}^\pi;x,a)
 \le\alpha_t+\beta_t\chi_t.
\]
Taking the infimum over $a$ requires neither existence of a minimizer nor a new selection assumption. If the proposal is rejected, then
$\overline d_t+\beta_tR_t(x;v,\pi_t)>0$ and
$\overline d_t\le d_t+\zeta_t$. Hence
\[
 Q_t(h_{t+1};x,\pi_t)
 <Q_t(h_{t+1};x,v)+\zeta_t+\beta_tR_t(x;v,\pi_t).
\]
Comparison with each feasible $a$ and another use of the action certificate yield an excess over $T_tJ_{t+1}^\pi$ of at most
$\alpha_t+\zeta_t+2\beta_t\chi_t$.
This common upper bound covers acceptance as well. Since future costs weakly decrease,
\[
 J_t^{\pi^+}=\mathcal T_t^{\pi^+}J_{t+1}^{\pi^+}
 \le\mathcal T_t^{\pi^+}J_{t+1}^{\pi}
 \le T_tJ_{t+1}^{\pi}+\tilde\epsilon_t.
\]
Monotonicity and the $\beta_t$ supremum-norm Lipschitz property of $T_t$ give the asserted recursion. The terminal gap is zero, and repeated substitution along successive passes gives the stated finite-horizon unrolling. Every retained policy is feasible for the original optimum, so a zero upper bound implies equality with that optimum. A fixed information contract may preclude zero greedy-search error; the theorem does not assume it away.

For sharpness take one period, two actions leading to distinct deterministic terminal states, zero current cost, and $\beta=1$. For $w>0$ and $0<\delta<w$, let the terminal critic be $(w,\delta)$ and the terminal cost be $(2w,\delta)$. The incumbent selects the first state and the exact greedy proposal the second. The exact action-error contrast is $w$, so $\chi=w$. The gate rejects because the critic contrast plus its allowance is $\delta>0$. The retained policy's gap is $2w-\delta$, whereas the terminal optimality gap and both search and arithmetic errors are zero. Letting $\delta$ decrease to zero excludes every universal coefficient below two.
\end{proof}

\begin{proof}[Proof of Proposition~\ref{M-prop:coupled52}]
The terminal assertion follows from $e_T=r_T$. At an earlier date, the policy-evaluation identity is
$e_t(x)=r_t(x)+\beta_tP_t^{\pi_t(x)}e_{t+1}(x)$.
The coupling's two marginals allow subtraction at states $x,y$ inside one integral. The triangle inequality and the inductive pair bound give
\[
 |e_t(x)-e_t(y)|\le H_t(x,y)+\beta_t
       \int D_{t+1}(u,v)\,d\Lambda_t^\pi(x,y)(u,v).
\]
The signed range supplies the separate bound $w_t$, so the minimum is valid. Coupling two action kernels at the same state gives the identical argument without the residual term and establishes the action certificate. Replacing an integral by an upper enclosure can only enlarge the bound before its intersection with the separately valid scalar width. This argument uses the marginals, not independence or optimality of the coupling.
\end{proof}

\begin{proof}[Proof of Corollary~\ref{M-cor:null52}]
The action difference annihilates $n$. Each of the two conditional means of $u$ lies in its verified interval, hence their difference in absolute value is at most its width. Further,
$J^\pi-(h+n)=u$, so the existing signed-residual recursion applied to the corrected critic yields the required interval without using an unknown true value. Adding $n$ to the continuation changes every feasible action objective by the same state-dependent amount; all action contrasts and approximate-greedy inequalities are therefore unchanged. The conclusion applies separately at each date and pass for which the membership and residual certificates are supplied.
\end{proof}

\begin{proof}[Proof of Proposition~\ref{M-prop:investment-null52}]
In the original dynamics the coefficients of $a$ in the first two coordinates are $1/2$ and $1/4$. For each common innovation $z$, the coefficient in $F_1(x,a,z)-2F_2(x,a,z)$ is consequently zero. Every function of that difference has identical realized values for two feasible actions with the same $z$, and thus identical conditional means under the original uniform innovation law. On the cube $y_1-2y_2+1$ ranges from $-1$ to $2$; its positive part ranges from zero to two. The global terminal error width is exactly $2M$. Its action contrast is nevertheless zero. The corrected critic is the original terminal cost $g$, for which the evaluation error vanishes. Proposal contrasts and their symbolic, outward-enclosed implementation are unchanged. Keeping the same proposals, tie conventions and cellwise gate therefore gives exactly the same terminal policy for every amplitude, not merely a similar observed cost.
\end{proof}

\subsection{Finite-state reference and limitations}
The exact reference implementation constructs a common-uniform coupling by exhausting ordered probability masses. It verifies both marginals, propagates pairwise residual bounds backward, and forms every same-state action contrast. Its regression suite independently evaluates the true finite-state policy values to test the certificates; those true values are not used by the certificate constructor. Tests cover one through five full improvement passes, a nonconstant action-null error with amplitudes zero, one, sixteen and 4096, non-null errors, constant shifts, exact feasibility, acceptance and rejection, and the sharp two-action example. These finite-state programs validate executable identities and counterexamples; they do not discretize the investment benchmark or establish full-horizon empirical convergence in that continuous economy.

\section{Complete construction records}\label{supp:records49}
The next table reports all distinct method, economy and resource rungs. Each row has three isolated repetitions with identical policy and certificate checkpoint hashes. Times are the median complete prefix, followed by its minimum and maximum; they do not omit unsuccessful earlier rungs. Counts and storage at full precision are in the machine-readable records. The table's decimal rounding is descriptive: exact rational bounds, not printed rounding, determine attainment.
{\small
\setlength{\tabcolsep}{3.5pt}
\begin{longtable}{@{}llrrrr@{}}
\caption{All distinct construction rungs and complete-prefix times}\label{tab:rungs49}\\
\toprule
$d;T;p$ & Method & $N,K,M$ & Bound & Median sec. & Range sec. \\
\midrule
\endfirsthead
\toprule
$d;T;p$ & Method & $N,K,M$ & Bound & Median sec. & Range sec. \\
\midrule
\endhead
\bottomrule
\endfoot
2;2;1 & W & 4,2,1 & 10.0042 & 0.014 & [0.014,0.017] \\
2;2;1 & W & 8,2,1 & 5.591673 & 0.034 & [0.034,0.036] \\
2;2;1 & W & 16,4,2 & 2.795836 & 0.092 & [0.091,0.093] \\
2;2;1 & W & 32,4,2 & 1.692705 & 0.290 & [0.290,0.291] \\
2;2;1 & W & 64,8,4 & 0.8463524 & 1.458 & [1.449,1.459] \\
2;2;1 & W & 128,16,8 & 0.4231762 & 10.658 & [10.612,10.696] \\
2;2;1 & A & 4,2,1 & 10.68755 & 0.042 & [0.042,0.043] \\
2;2;1 & A & 8,2,1 & 6.132349 & 0.090 & [0.090,0.091] \\
2;2;1 & A & 16,4,2 & 3.216146 & 0.176 & [0.176,0.178] \\
2;2;1 & A & 32,4,2 & 1.906506 & 0.399 & [0.398,0.402] \\
2;2;1 & A & 64,8,4 & 0.9622566 & 1.514 & [1.509,1.521] \\
2;2;1 & A & 128,16,8 & 0.4833928 & 9.817 & [9.748,9.900] \\
2;2;1 & F & 4,2,1 & 10.68755 & 0.029 & [0.029,0.029] \\
2;2;1 & F & 8,2,1 & 6.132349 & 0.061 & [0.061,0.063] \\
2;2;1 & F & 16,4,2 & 3.216146 & 0.119 & [0.118,0.122] \\
2;2;1 & F & 32,4,2 & 1.906506 & 0.269 & [0.268,0.274] \\
2;2;1 & F & 64,8,4 & 0.9622566 & 1.134 & [1.130,1.143] \\
2;2;1 & F & 128,16,8 & 0.4833928 & 8.458 & [8.455,8.491] \\
2;2;4 & W & 4,2,1 & 10.31855 & 0.014 & [0.014,0.014] \\
2;2;4 & W & 8,2,1 & 5.844473 & 0.034 & [0.034,0.035] \\
2;2;4 & W & 16,4,2 & 2.922236 & 0.092 & [0.092,0.096] \\
2;2;4 & W & 32,4,2 & 1.803718 & 0.290 & [0.287,0.303] \\
2;2;4 & W & 64,8,4 & 0.9018591 & 1.452 & [1.452,1.510] \\
2;2;4 & W & 128,16,8 & 0.4509295 & 10.689 & [10.584,10.884] \\
2;2;4 & A & 4,2,1 & 10.92323 & 0.043 & [0.042,0.043] \\
2;2;4 & A & 8,2,1 & 6.330362 & 0.091 & [0.090,0.092] \\
2;2;4 & A & 16,4,2 & 3.31506 & 0.177 & [0.176,0.179] \\
2;2;4 & A & 32,4,2 & 2.000228 & 0.397 & [0.397,0.403] \\
2;2;4 & A & 64,8,4 & 1.009136 & 1.509 & [1.509,1.524] \\
2;2;4 & A & 128,16,8 & 0.5068371 & 9.797 & [9.763,9.822] \\
2;2;4 & F & 4,2,1 & 10.92323 & 0.029 & [0.028,0.033] \\
2;2;4 & F & 8,2,1 & 6.330362 & 0.061 & [0.060,0.065] \\
2;2;4 & F & 16,4,2 & 3.31506 & 0.118 & [0.118,0.123] \\
2;2;4 & F & 32,4,2 & 2.000228 & 0.268 & [0.268,0.273] \\
2;2;4 & F & 64,8,4 & 1.009136 & 1.134 & [1.131,1.145] \\
2;2;4 & F & 128,16,8 & 0.5068371 & 8.453 & [8.373,8.520] \\
2;3;1 & W & 4,2,1 & 14.5143 & 0.019 & [0.019,0.019] \\
2;3;1 & W & 8,2,1 & 8.228071 & 0.046 & [0.046,0.047] \\
2;3;1 & W & 16,4,2 & 4.114036 & 0.126 & [0.126,0.129] \\
2;3;1 & W & 32,4,2 & 2.542479 & 0.401 & [0.398,0.408] \\
2;3;1 & W & 64,8,4 & 1.271239 & 2.070 & [2.062,2.110] \\
2;3;1 & W & 128,16,8 & 0.6356197 & 15.635 & [15.607,15.924] \\
2;3;1 & A & 4,2,1 & 14.97615 & 0.061 & [0.060,0.062] \\
2;3;1 & A & 8,2,1 & 8.615403 & 0.130 & [0.129,0.131] \\
2;3;1 & A & 16,4,2 & 4.533221 & 0.252 & [0.251,0.255] \\
2;3;1 & A & 32,4,2 & 2.723572 & 0.568 & [0.565,0.574] \\
2;3;1 & A & 64,8,4 & 1.374878 & 2.148 & [2.145,2.184] \\
2;3;1 & A & 128,16,8 & 0.6908084 & 14.176 & [14.130,14.444] \\
2;3;1 & F & 4,2,1 & 14.97615 & 0.041 & [0.041,0.042] \\
2;3;1 & F & 8,2,1 & 8.615403 & 0.088 & [0.087,0.088] \\
2;3;1 & F & 16,4,2 & 4.533221 & 0.171 & [0.170,0.192] \\
2;3;1 & F & 32,4,2 & 2.723572 & 0.388 & [0.388,0.412] \\
2;3;1 & F & 64,8,4 & 1.374878 & 1.651 & [1.647,1.677] \\
2;3;1 & F & 128,16,8 & 0.6908084 & 12.420 & [12.409,12.456] \\
2;3;4 & W & 4,2,1 & 15.02035 & 0.019 & [0.019,0.019] \\
2;3;4 & W & 8,2,1 & 8.625736 & 0.046 & [0.046,0.048] \\
2;3;4 & W & 16,4,2 & 4.312868 & 0.127 & [0.125,0.129] \\
2;3;4 & W & 32,4,2 & 2.714215 & 0.402 & [0.397,0.408] \\
2;3;4 & W & 64,8,4 & 1.357108 & 2.085 & [2.053,2.103] \\
2;3;4 & W & 128,16,8 & 0.6785539 & 15.815 & [15.506,15.955] \\
2;3;4 & A & 4,2,1 & 15.30603 & 0.061 & [0.061,0.061] \\
2;3;4 & A & 8,2,1 & 8.884068 & 0.131 & [0.130,0.131] \\
2;3;4 & A & 16,4,2 & 4.667275 & 0.253 & [0.252,0.254] \\
2;3;4 & A & 32,4,2 & 2.853489 & 0.569 & [0.569,0.569] \\
2;3;4 & A & 64,8,4 & 1.440006 & 2.169 & [2.165,2.170] \\
2;3;4 & A & 128,16,8 & 0.723386 & 14.257 & [14.181,14.260] \\
2;3;4 & F & 4,2,1 & 15.30603 & 0.042 & [0.041,0.042] \\
2;3;4 & F & 8,2,1 & 8.884068 & 0.088 & [0.087,0.088] \\
2;3;4 & F & 16,4,2 & 4.667275 & 0.170 & [0.170,0.171] \\
2;3;4 & F & 32,4,2 & 2.853489 & 0.390 & [0.387,0.391] \\
2;3;4 & F & 64,8,4 & 1.440006 & 1.643 & [1.637,1.647] \\
2;3;4 & F & 128,16,8 & 0.723386 & 12.435 & [12.393,12.446] \\
3;2;1 & W & 4,2,1 & 10.09191 & 0.033 & [0.033,0.033] \\
3;2;1 & W & 8,4,2 & 5.045957 & 0.225 & [0.224,0.225] \\
3;2;1 & W & 16,8,4 & 2.522979 & 2.424 & [2.422,2.453] \\
3;2;1 & W & 32,8,4 & 1.415656 & 16.615 & [16.590,16.644] \\
3;2;1 & F & 4,2,1 & 10.91291 & 0.096 & [0.095,0.096] \\
3;2;1 & F & 8,4,2 & 5.756259 & 0.318 & [0.318,0.322] \\
3;2;1 & F & 16,8,4 & 2.983083 & 2.074 & [2.059,2.097] \\
3;2;1 & F & 32,8,4 & 1.654898 & 13.687 & [13.539,13.776] \\
3;3;1 & W & 4,2,1 & 14.67402 & 0.045 & [0.044,0.046] \\
3;3;1 & W & 8,4,2 & 7.337012 & 0.317 & [0.313,0.322] \\
3;3;1 & W & 16,8,4 & 3.668506 & 3.508 & [3.476,3.527] \\
3;3;1 & W & 32,8,4 & 2.088657 & 23.987 & [23.952,24.055] \\
3;3;1 & F & 4,2,1 & 15.28673 & 0.140 & [0.140,0.141] \\
3;3;1 & F & 8,4,2 & 8.038961 & 0.467 & [0.463,0.468] \\
3;3;1 & F & 16,8,4 & 4.173215 & 3.032 & [3.027,3.059] \\
3;3;1 & F & 32,8,4 & 2.334703 & 20.040 & [19.984,20.139] \\
4;2;1 & W & 4,2,1 & 10.0042 & 0.167 & [0.164,0.168] \\
4;2;1 & W & 8,4,2 & 5.002099 & 2.588 & [2.561,2.635] \\
4;2;1 & W & 16,8,4 & 2.50105 & 58.569 & [58.203,59.507] \\
4;2;1 & F & 4,2,1 & 10.98489 & 0.523 & [0.520,0.526] \\
4;2;1 & F & 8,4,2 & 5.772966 & 4.388 & [4.386,4.449] \\
4;2;1 & F & 16,8,4 & 2.981482 & 60.809 & [60.637,61.583] \\
4;3;1 & W & 4,2,1 & 14.5143 & 0.230 & [0.230,0.231] \\
4;3;1 & W & 8,4,2 & 7.257149 & 3.663 & [3.649,3.688] \\
4;3;1 & W & 16,8,4 & 3.628575 & 82.156 & [81.795,82.676] \\
4;3;1 & F & 4,2,1 & 15.40485 & 0.782 & [0.776,0.784] \\
4;3;1 & F & 8,4,2 & 8.075833 & 6.457 & [6.444,6.547] \\
4;3;1 & F & 16,8,4 & 4.176837 & 90.099 & [89.833,90.155] \\
\end{longtable}
}


\section{All direct returned-policy comparisons}\label{supp:direct49}
The labels U, L, C and H denote the uniform initial law and point laws at $(1/8,1/8)$, $(1/2,1/2)$ and $(7/8,7/8)$. Sensor E denotes full information; sensor 8 uses the eight-bit cell contract and downward action spacing. The historical favorable pair is marked R48; other pairs are the actual R49 first crossings. Every interval concerns expected witness cost minus expected FVI cost. All sampled intervals in this table and the next belong to the same new 99 percent family under the independent-bin model.
{\small
\setlength{\tabcolsep}{3.5pt}
\begin{longtable}{@{}llrrrrr@{}}
\caption{Every actual first-crossing policy-cost interval}\label{tab:all-direct49}\\
\toprule
Pair & $T;p$ & $\varepsilon$ & Law & Bits & Lower & Upper \\
\midrule
\endfirsthead
\toprule
Pair & $T;p$ & $\varepsilon$ & Law & Bits & Lower & Upper \\
\midrule
\endhead
\bottomrule
\endfoot
R49 & 2;1 & 4 & U & E & 0.00543396 & 0.00798729 \\
R49 & 2;1 & 4 & U & 8 & 0.00541635 & 0.00796936 \\
R49 & 2;1 & 4 & L & E & -0.00156041 & 0.000833084 \\
R49 & 2;1 & 4 & L & 8 & -0.00156513 & 0.000828396 \\
R49 & 2;1 & 4 & C & E & 0.00372573 & 0.0061508 \\
R49 & 2;1 & 4 & C & 8 & 0.00375513 & 0.00618077 \\
R49 & 2;1 & 4 & H & E & 0.00609378 & 0.00853627 \\
R49 & 2;1 & 4 & H & 8 & 0.00589939 & 0.00833643 \\
R49 & 2;1 & 2 & U & E & 0.0143531 & 0.0170811 \\
R49 & 2;1 & 2 & U & 8 & 0.0143519 & 0.0170783 \\
R49 & 2;1 & 2 & L & E & 0.0547607 & 0.0574082 \\
R49 & 2;1 & 2 & L & 8 & 0.0562672 & 0.0589226 \\
R49 & 2;1 & 2 & C & E & 0.00754534 & 0.0100938 \\
R49 & 2;1 & 2 & C & 8 & 0.00764868 & 0.0101989 \\
R49 & 2;1 & 2 & H & E & 0.00545776 & 0.00790326 \\
R49 & 2;1 & 2 & H & 8 & 0.00540822 & 0.00785334 \\
R49 & 2;1 & 1 & U & E & 0.00608323 & 0.00862331 \\
R49 & 2;1 & 1 & U & 8 & 0.00608821 & 0.00862733 \\
R49 & 2;1 & 1 & L & E & 0.0241465 & 0.0266462 \\
R49 & 2;1 & 1 & L & 8 & 0.0243754 & 0.0268759 \\
R49 & 2;1 & 1 & C & E & 0.0101329 & 0.0126139 \\
R49 & 2;1 & 1 & C & 8 & 0.0101188 & 0.012601 \\
R49 & 2;1 & 1 & H & E & 0.00787923 & 0.0103208 \\
R49 & 2;1 & 1 & H & 8 & 0.00773077 & 0.0101721 \\
R49 & 2;4 & 4 & U & E & 0.0048609 & 0.00758426 \\
R49 & 2;4 & 4 & U & 8 & 0.00488276 & 0.00760722 \\
R49 & 2;4 & 4 & L & E & 0.0185747 & 0.0211388 \\
R49 & 2;4 & 4 & L & 8 & 0.0186476 & 0.0212113 \\
R49 & 2;4 & 4 & C & E & 0.0347841 & 0.0374079 \\
R49 & 2;4 & 4 & C & 8 & 0.0346533 & 0.0372768 \\
R49 & 2;4 & 4 & H & E & -8.33958e-05 & 0.00253601 \\
R49 & 2;4 & 4 & H & 8 & 0.000139912 & 0.00276397 \\
R49 & 2;4 & 2 & U & E & 0.0126061 & 0.0154244 \\
R49 & 2;4 & 2 & U & 8 & 0.0124777 & 0.0152934 \\
R49 & 2;4 & 2 & L & E & 0.0213281 & 0.0239062 \\
R49 & 2;4 & 2 & L & 8 & 0.0216307 & 0.0242254 \\
R49 & 2;4 & 2 & C & E & 0.0363241 & 0.0389543 \\
R49 & 2;4 & 2 & C & 8 & 0.0365806 & 0.0392119 \\
R49 & 2;4 & 2 & H & E & 0.0408571 & 0.0435165 \\
R49 & 2;4 & 2 & H & 8 & 0.0405602 & 0.0432247 \\
R49 & 2;4 & 1 & U & E & 0.00493422 & 0.0076087 \\
R49 & 2;4 & 1 & U & 8 & 0.00498356 & 0.00765916 \\
R49 & 2;4 & 1 & L & E & 0.00791584 & 0.0104887 \\
R49 & 2;4 & 1 & L & 8 & 0.00799778 & 0.010571 \\
R49 & 2;4 & 1 & C & E & 0.016087 & 0.018698 \\
R49 & 2;4 & 1 & C & 8 & 0.0157071 & 0.0183196 \\
R49 & 2;4 & 1 & H & E & 0.00782793 & 0.0104847 \\
R49 & 2;4 & 1 & H & 8 & 0.00787666 & 0.0105306 \\
R49 & 3;1 & 4 & U & E & 0.0253656 & 0.0288147 \\
R49 & 3;1 & 4 & U & 8 & 0.0252662 & 0.0287151 \\
R49 & 3;1 & 4 & L & E & 0.0788767 & 0.0824397 \\
R49 & 3;1 & 4 & L & 8 & 0.0778317 & 0.0813824 \\
R49 & 3;1 & 4 & C & E & 0.100968 & 0.104201 \\
R49 & 3;1 & 4 & C & 8 & 0.101486 & 0.104716 \\
R49 & 3;1 & 4 & H & E & 0.0197496 & 0.0229893 \\
R49 & 3;1 & 4 & H & 8 & 0.019754 & 0.0229886 \\
R49 & 3;1 & 2 & U & E & 0.0109993 & 0.0141923 \\
R49 & 3;1 & 2 & U & 8 & 0.0110585 & 0.0142512 \\
R49 & 3;1 & 2 & L & E & 0.0282946 & 0.0314077 \\
R49 & 3;1 & 2 & L & 8 & 0.0279109 & 0.0310247 \\
R49 & 3;1 & 2 & C & E & 0.0176355 & 0.0207641 \\
R49 & 3;1 & 2 & C & 8 & 0.0171906 & 0.0203209 \\
R49 & 3;1 & 2 & H & E & 0.0126606 & 0.0156975 \\
R49 & 3;1 & 2 & H & 8 & 0.0126999 & 0.0157435 \\
R49 & 3;1 & 1 & U & E & 0.00423544 & 0.00731756 \\
R49 & 3;1 & 1 & U & 8 & 0.00451493 & 0.0075994 \\
R49 & 3;1 & 1 & L & E & 0.00514129 & 0.008227 \\
R49 & 3;1 & 1 & L & 8 & 0.00545945 & 0.00854446 \\
R49 & 3;1 & 1 & C & E & 0.00460682 & 0.00763251 \\
R49 & 3;1 & 1 & C & 8 & 0.00492024 & 0.00794768 \\
R49 & 3;1 & 1 & H & E & 0.00035797 & 0.00337421 \\
R49 & 3;1 & 1 & H & 8 & 0.00036318 & 0.00337929 \\
R49 & 3;4 & 4 & U & E & 0.0186753 & 0.022252 \\
R49 & 3;4 & 4 & U & 8 & 0.0186419 & 0.0222187 \\
R49 & 3;4 & 4 & L & E & 0.0276014 & 0.0309773 \\
R49 & 3;4 & 4 & L & 8 & 0.0279277 & 0.0313061 \\
R49 & 3;4 & 4 & C & E & 0.0549145 & 0.058394 \\
R49 & 3;4 & 4 & C & 8 & 0.0535423 & 0.0569789 \\
R49 & 3;4 & 4 & H & E & 0.0500905 & 0.0536733 \\
R49 & 3;4 & 4 & H & 8 & 0.0501068 & 0.0536738 \\
R49 & 3;4 & 2 & U & E & 0.00722022 & 0.0105987 \\
R49 & 3;4 & 2 & U & 8 & 0.00727216 & 0.0106525 \\
R49 & 3;4 & 2 & L & E & 0.0111161 & 0.0144007 \\
R49 & 3;4 & 2 & L & 8 & 0.0111398 & 0.0144261 \\
R49 & 3;4 & 2 & C & E & 0.0119992 & 0.0152703 \\
R49 & 3;4 & 2 & C & 8 & 0.0120479 & 0.0153198 \\
R49 & 3;4 & 2 & H & E & 0.00896933 & 0.0122714 \\
R49 & 3;4 & 2 & H & 8 & 0.00896107 & 0.0122625 \\
R49 & 3;4 & 1 & U & E & 0.00231571 & 0.00560954 \\
R49 & 3;4 & 1 & U & 8 & 0.00249689 & 0.00579256 \\
R49 & 3;4 & 1 & L & E & 0.00405598 & 0.00731121 \\
R49 & 3;4 & 1 & L & 8 & 0.000239139 & 0.00350007 \\
R49 & 3;4 & 1 & C & E & 0.00285917 & 0.00611809 \\
R49 & 3;4 & 1 & C & 8 & 0.00270615 & 0.00596212 \\
R49 & 3;4 & 1 & H & E & 0.000774402 & 0.00403881 \\
R49 & 3;4 & 1 & H & 8 & 0.000787911 & 0.0040506 \\
R48 & 3;4 & 2 & U & E & -0.00139342 & 0.00183546 \\
R48 & 3;4 & 2 & U & 8 & -0.00138819 & 0.00184069 \\
R48 & 3;4 & 2 & L & E & -0.00148104 & 0.00174625 \\
R48 & 3;4 & 2 & L & 8 & -0.00148102 & 0.00174626 \\
R48 & 3;4 & 2 & C & E & -0.00135341 & 0.0018749 \\
R48 & 3;4 & 2 & C & 8 & -0.00134267 & 0.00188569 \\
R48 & 3;4 & 2 & H & E & -0.00139997 & 0.00182783 \\
R48 & 3;4 & 2 & H & 8 & -0.00141355 & 0.00181414 \\
\end{longtable}
}
\noindent{\footnotesize Endpoints are displayed to six significant digits; exact rational endpoints and policy hashes govern every sign and fee conclusion. R48 rows evaluate the old favorable N=32 witness versus N=64 FVI pair under fresh independent interval paths.\par}


\section{All sensor-policy contrasts}\label{supp:sensors49}
Each row compares the actual $b$-bit implementation with the same returned controller at 16 bits, under the uniform initial law. The 16-bit self-rows are identities. The corresponding price selectors and exact continuous-price envelopes use these rows without additional sampling. Upper-net-cost regret is relative to the best bit in this catalogue, not an unobserved optimum over continuous precision or arbitrary sensors.
{\small
\setlength{\tabcolsep}{3.5pt}
\begin{longtable}{@{}llrrr@{}}
\caption{Every actual sensor-cost difference relative to 16 bits}\label{tab:all-sensor49}\\
\toprule
$T;p$ & Method & Bits & Lower & Upper \\
\midrule
\endfirsthead
\toprule
$T;p$ & Method & Bits & Lower & Upper \\
\midrule
\endhead
\bottomrule
\endfoot
2;1 & W & 4 & 0.000261375 & 0.00285711 \\
2;1 & W & 5 & -0.000427494 & 0.00215671 \\
2;1 & W & 6 & -0.00204196 & 0.000521085 \\
2;1 & W & 7 & -0.00088098 & 0.00160987 \\
2;1 & W & 8 & -0.00122806 & 0.00124059 \\
2;1 & W & 9 & -0.00121323 & 0.00123098 \\
2;1 & W & 10 & -0.00120642 & 0.00122091 \\
2;1 & W & 11 & -0.00119893 & 0.00121541 \\
2;1 & W & 12 & -0.00120501 & 0.00120022 \\
2;1 & W & 13 & -0.0011989 & 0.00119995 \\
2;1 & W & 14 & -0.00119712 & 0.00119717 \\
2;1 & W & 15 & -0.00119517 & 0.00119587 \\
2;1 & W & 16 & 0 & 0 \\
2;1 & F & 4 & 0.000199127 & 0.0026252 \\
2;1 & F & 5 & -0.000584413 & 0.00181963 \\
2;1 & F & 6 & -0.00085593 & 0.00153971 \\
2;1 & F & 7 & -0.00111845 & 0.00126841 \\
2;1 & F & 8 & -0.00117522 & 0.00120999 \\
2;1 & F & 9 & -0.00118398 & 0.00120093 \\
2;1 & F & 10 & -0.00118806 & 0.00119669 \\
2;1 & F & 11 & -0.00119023 & 0.00119442 \\
2;1 & F & 12 & -0.00119128 & 0.00119331 \\
2;1 & F & 13 & -0.00119182 & 0.00119274 \\
2;1 & F & 14 & -0.00119207 & 0.00119247 \\
2;1 & F & 15 & -0.00119219 & 0.00119235 \\
2;1 & F & 16 & 0 & 0 \\
2;4 & W & 4 & -4.77011e-05 & 0.00267196 \\
2;4 & W & 5 & -0.000905067 & 0.00180125 \\
2;4 & W & 6 & -0.00220586 & 0.000485462 \\
2;4 & W & 7 & -0.00108854 & 0.00153925 \\
2;4 & W & 8 & -0.00124312 & 0.00137376 \\
2;4 & W & 9 & -0.00131086 & 0.00128822 \\
2;4 & W & 10 & -0.00128939 & 0.00129611 \\
2;4 & W & 11 & -0.00128416 & 0.00129219 \\
2;4 & W & 12 & -0.00128539 & 0.00128378 \\
2;4 & W & 13 & -0.00128173 & 0.0012827 \\
2;4 & W & 14 & -0.00128018 & 0.00128022 \\
2;4 & W & 15 & -0.00127809 & 0.00128043 \\
2;4 & W & 16 & 0 & 0 \\
2;4 & F & 4 & -0.00112704 & 0.00143668 \\
2;4 & F & 5 & -0.00121844 & 0.00133951 \\
2;4 & F & 6 & -0.00125571 & 0.00129926 \\
2;4 & F & 7 & -0.00126208 & 0.00129244 \\
2;4 & F & 8 & -0.00127554 & 0.00127722 \\
2;4 & F & 9 & -0.00127593 & 0.00127681 \\
2;4 & F & 10 & -0.00127615 & 0.00127658 \\
2;4 & F & 11 & -0.00127626 & 0.00127646 \\
2;4 & F & 12 & -0.0012763 & 0.00127642 \\
2;4 & F & 13 & -0.00127634 & 0.00127638 \\
2;4 & F & 14 & -0.00127635 & 0.00127637 \\
2;4 & F & 15 & -0.00127636 & 0.00127636 \\
2;4 & F & 16 & 0 & 0 \\
3;1 & W & 4 & -0.000101036 & 0.00300605 \\
3;1 & W & 5 & -0.00041666 & 0.00269816 \\
3;1 & W & 6 & -0.000850866 & 0.00225551 \\
3;1 & W & 7 & -0.0023459 & 0.000745795 \\
3;1 & W & 8 & -0.00124523 & 0.00180352 \\
3;1 & W & 9 & -0.00147051 & 0.00156613 \\
3;1 & W & 10 & -0.00150886 & 0.00151075 \\
3;1 & W & 11 & -0.00150001 & 0.00150822 \\
3;1 & W & 12 & -0.00150165 & 0.00149869 \\
3;1 & W & 13 & -0.00149858 & 0.00149581 \\
3;1 & W & 14 & -0.00149663 & 0.00149349 \\
3;1 & W & 15 & -0.00149339 & 0.00149403 \\
3;1 & W & 16 & 0 & 0 \\
3;1 & F & 4 & 0.000691523 & 0.00372718 \\
3;1 & F & 5 & -0.000505628 & 0.00250065 \\
3;1 & F & 6 & -0.0010836 & 0.00190806 \\
3;1 & F & 7 & -0.00122457 & 0.00176312 \\
3;1 & F & 8 & -0.001473 & 0.00150813 \\
3;1 & F & 9 & -0.0014817 & 0.00149918 \\
3;1 & F & 10 & -0.0014861 & 0.00149463 \\
3;1 & F & 11 & -0.00148819 & 0.00149247 \\
3;1 & F & 12 & -0.00148929 & 0.00149132 \\
3;1 & F & 13 & -0.00148981 & 0.00149079 \\
3;1 & F & 14 & -0.00149008 & 0.0014905 \\
3;1 & F & 15 & -0.00149022 & 0.00149036 \\
3;1 & F & 16 & 0 & 0 \\
3;4 & W & 4 & -0.00115463 & 0.00216779 \\
3;4 & W & 5 & -0.00128854 & 0.00202944 \\
3;4 & W & 6 & -0.0012672 & 0.00204406 \\
3;4 & W & 7 & -0.00203778 & 0.00126435 \\
3;4 & W & 8 & -0.00147323 & 0.00179482 \\
3;4 & W & 9 & -0.00162558 & 0.00163607 \\
3;4 & W & 10 & -0.00162825 & 0.00162325 \\
3;4 & W & 11 & -0.00162181 & 0.00162209 \\
3;4 & W & 12 & -0.00161759 & 0.00162064 \\
3;4 & W & 13 & -0.00161723 & 0.00161692 \\
3;4 & W & 14 & -0.00161637 & 0.00161497 \\
3;4 & W & 15 & -0.00161356 & 0.00161547 \\
3;4 & W & 16 & 0 & 0 \\
3;4 & F & 4 & -0.00143754 & 0.00179918 \\
3;4 & F & 5 & -0.0015455 & 0.00168516 \\
3;4 & F & 6 & -0.0015881 & 0.00163939 \\
3;4 & F & 7 & -0.00159646 & 0.0016305 \\
3;4 & F & 8 & -0.00161162 & 0.00161347 \\
3;4 & F & 9 & -0.00161206 & 0.00161302 \\
3;4 & F & 10 & -0.00161229 & 0.00161278 \\
3;4 & F & 11 & -0.00161241 & 0.00161265 \\
3;4 & F & 12 & -0.00161247 & 0.00161258 \\
3;4 & F & 13 & -0.0016125 & 0.00161255 \\
3;4 & F & 14 & -0.00161251 & 0.00161254 \\
3;4 & F & 15 & -0.00161252 & 0.00161253 \\
3;4 & F & 16 & 0 & 0 \\
\end{longtable}
}
\noindent{\footnotesize All rows use the uniform initial law and the actual first-target-one controller. Self-comparisons are exact identities. Both policies use the declared robust repair and action quantum.\par}


\section{Reconstruction and preservation}\label{supp:audit49}
The release verifies frozen science-source identities, every recorded checkpoint, all signed policy formulas, every first crossing, complete-prefix counts, repeated policy identities and the confidence intervals reconstructed from endpoint moments. Additional exact tests check the primitive regions, resource allocation, original-owner identity, fee semantics and net-choice bound. The complete development article and supplement retain all baseline labeled results; the original paths and adverse results are unchanged. A clean Git-archive rebuild removes generated PDFs and derived tables, then rebuilds from ordinary committed sources and frozen records without network or new scientific execution. Compilation, finite tests and source identities do not constitute independent mathematical peer approval.
\section{Every R50 actual-cost interval and mechanism record}\label{supp:records50}
The following tables use the corrected outward-subtraction replay. C denotes the separate common-accuracy block, P the primary target-five plan, Q the primary target-two plan, O the original R49 target-two policies, A residual-driven FVI, and I the isotropic block. U is the uniform initial law; L, M and H denote point laws at 1/8,1/2 and 7/8 in every coordinate. Each group has four absolute costs and six contrasts; the table lists all 280 estimands. Printed endpoints are rounded for display. Exact rational endpoints and source-bound moments govern inference.
{\small
\setlength{\tabcolsep}{3pt}
\begin{longtable}{@{}llrr@{}}
\caption{Every new absolute expected policy-cost interval}\label{tab:all-costs50}\\
\toprule
Case & Policy & Lower & Upper \\
\midrule
\endfirsthead
\toprule
Case & Policy & Lower & Upper \\
\midrule
\endhead
\bottomrule
\endfoot
C2;2;1;U & left & 0.6242 & 0.6378 \\
C2;2;1;U & left-repaired & 0.61117 & 0.62482 \\
C2;2;1;U & right & 0.59821 & 0.61177 \\
C2;2;1;U & right-repaired & 0.59628 & 0.60988 \\
C3;2;1;U & left & 0.58992 & 0.60068 \\
C3;2;1;U & left-repaired & 0.58062 & 0.59146 \\
C3;2;1;U & right & 0.56904 & 0.57982 \\
C3;2;1;U & right-repaired & 0.56768 & 0.57849 \\
C4;2;1;U & left & 0.62188 & 0.63199 \\
C4;2;1;U & left-repaired & 0.60399 & 0.61391 \\
C4;2;1;U & right & 0.5855 & 0.59525 \\
C4;2;1;U & right-repaired & 0.5854 & 0.59515 \\
I2;2;1;U & left & 0.63959 & 0.65386 \\
I2;2;1;U & left-repaired & 0.61456 & 0.6285 \\
I2;2;1;U & right & 0.59407 & 0.60781 \\
I2;2;1;U & right-repaired & 0.59363 & 0.60737 \\
O2;2;1;M & left & 0.30128 & 0.30382 \\
O2;2;1;M & left-repaired & 0.29236 & 0.29488 \\
O2;2;1;M & right & 0.29238 & 0.29489 \\
O2;2;1;M & right-repaired & 0.29236 & 0.29488 \\
O2;2;1;L & left & 1.5737 & 1.5766 \\
O2;2;1;L & left-repaired & 1.5626 & 1.565 \\
O2;2;1;L & right & 1.5163 & 1.5188 \\
O2;2;1;L & right-repaired & 1.5162 & 1.5186 \\
O2;2;1;H & left & 0.19157 & 0.19409 \\
O2;2;1;H & left-repaired & 0.18888 & 0.19135 \\
O2;2;1;H & right & 0.18497 & 0.18741 \\
O2;2;1;H & right-repaired & 0.18423 & 0.18668 \\
O2;2;1;U & left & 0.61089 & 0.62464 \\
O2;2;1;U & left-repaired & 0.60392 & 0.61765 \\
O2;2;1;U & right & 0.59513 & 0.60882 \\
O2;2;1;U & right-repaired & 0.59491 & 0.6086 \\
O2;2;4;M & left & 0.42746 & 0.43017 \\
O2;2;4;M & left-repaired & 0.42249 & 0.42507 \\
O2;2;4;M & right & 0.38962 & 0.39224 \\
O2;2;4;M & right-repaired & 0.38952 & 0.39214 \\
O2;2;4;L & left & 1.6588 & 1.6614 \\
O2;2;4;L & left-repaired & 1.6564 & 1.659 \\
O2;2;4;L & right & 1.6358 & 1.6385 \\
O2;2;4;L & right-repaired & 1.6358 & 1.6385 \\
O2;2;4;H & left & 0.2474 & 0.2501 \\
O2;2;4;H & left-repaired & 0.24264 & 0.24526 \\
O2;2;4;H & right & 0.20559 & 0.20818 \\
O2;2;4;H & right-repaired & 0.20554 & 0.20813 \\
O2;2;4;U & left & 0.69261 & 0.7071 \\
O2;2;4;U & left-repaired & 0.68704 & 0.70155 \\
O2;2;4;U & right & 0.67867 & 0.69319 \\
O2;2;4;U & right-repaired & 0.67849 & 0.69301 \\
O2;3;1;M & left & 0.42649 & 0.42982 \\
O2;3;1;M & left-repaired & 0.42326 & 0.42657 \\
O2;3;1;M & right & 0.40778 & 0.41096 \\
O2;3;1;M & right-repaired & 0.40776 & 0.41095 \\
O2;3;1;L & left & 1.6968 & 1.7 \\
O2;3;1;L & left-repaired & 1.6932 & 1.6964 \\
O2;3;1;L & right & 1.6673 & 1.6705 \\
O2;3;1;L & right-repaired & 1.6672 & 1.6704 \\
O2;3;1;H & left & 0.29136 & 0.29449 \\
O2;3;1;H & left-repaired & 0.28954 & 0.29265 \\
O2;3;1;H & right & 0.27712 & 0.28025 \\
O2;3;1;H & right-repaired & 0.27696 & 0.28009 \\
O2;3;1;U & left & 0.71942 & 0.73411 \\
O2;3;1;U & left-repaired & 0.71638 & 0.73106 \\
O2;3;1;U & right & 0.70681 & 0.72143 \\
O2;3;1;U & right-repaired & 0.70676 & 0.72139 \\
O2;3;4;M & left & 0.58099 & 0.58436 \\
O2;3;4;M & left-repaired & 0.57925 & 0.58262 \\
O2;3;4;M & right & 0.56732 & 0.57065 \\
O2;3;4;M & right-repaired & 0.56709 & 0.57042 \\
O2;3;4;L & left & 1.8582 & 1.8616 \\
O2;3;4;L & left-repaired & 1.856 & 1.8594 \\
O2;3;4;L & right & 1.8454 & 1.8487 \\
O2;3;4;L & right-repaired & 1.8454 & 1.8487 \\
O2;3;4;H & left & 0.35451 & 0.35788 \\
O2;3;4;H & left-repaired & 0.35176 & 0.35509 \\
O2;3;4;H & right & 0.34392 & 0.34723 \\
O2;3;4;H & right-repaired & 0.34391 & 0.34721 \\
O2;3;4;U & left & 0.86235 & 0.87799 \\
O2;3;4;U & left-repaired & 0.86048 & 0.87613 \\
O2;3;4;U & right & 0.85339 & 0.86906 \\
O2;3;4;U & right-repaired & 0.85323 & 0.86891 \\
Q2;2;1;U & left & 0.60972 & 0.62332 \\
Q2;2;1;U & left-repaired & 0.60276 & 0.61638 \\
Q2;2;1;U & right & 0.59518 & 0.60882 \\
Q2;2;1;U & right-repaired & 0.59444 & 0.6081 \\
P2;2;1;U & left & 0.6222 & 0.63577 \\
P2;2;1;U & left-repaired & 0.60912 & 0.62274 \\
P2;2;1;U & right & 0.59639 & 0.60993 \\
P2;2;1;U & right-repaired & 0.59445 & 0.60803 \\
Q2;2;4;U & left & 0.69243 & 0.70702 \\
Q2;2;4;U & left-repaired & 0.68754 & 0.70213 \\
Q2;2;4;U & right & 0.68106 & 0.69562 \\
Q2;2;4;U & right-repaired & 0.68086 & 0.69542 \\
Q2;3;1;U & left & 0.72237 & 0.73696 \\
Q2;3;1;U & left-repaired & 0.71906 & 0.73365 \\
Q2;3;1;U & right & 0.70782 & 0.7224 \\
Q2;3;1;U & right-repaired & 0.70769 & 0.72228 \\
Q2;3;4;U & left & 0.86908 & 0.88468 \\
Q2;3;4;U & left-repaired & 0.86539 & 0.881 \\
Q2;3;4;U & right & 0.85511 & 0.87074 \\
Q2;3;4;U & right-repaired & 0.85451 & 0.87015 \\
P3;2;1;U & left & 0.58981 & 0.60055 \\
P3;2;1;U & left-repaired & 0.58046 & 0.59128 \\
P3;2;1;U & right & 0.56902 & 0.57978 \\
P3;2;1;U & right-repaired & 0.56764 & 0.57844 \\
P4;2;1;U & left & 0.61951 & 0.62957 \\
P4;2;1;U & left-repaired & 0.60168 & 0.61155 \\
P4;2;1;U & right & 0.58357 & 0.59328 \\
P4;2;1;U & right-repaired & 0.58332 & 0.59303 \\
A2;2;1;U & left & 0.59772 & 0.61129 \\
A2;2;1;U & left-repaired & 0.59585 & 0.60945 \\
A2;2;1;U & right & 0.59777 & 0.61134 \\
A2;2;1;U & right-repaired & 0.59585 & 0.60946 \\
\end{longtable}
}

{\small
\setlength{\tabcolsep}{3pt}
\begin{longtable}{@{}llrr@{}}
\caption{Every new direct expected policy-cost contrast}\label{tab:all-contrasts50}\\
\toprule
Case & Contrast & Lower & Upper \\
\midrule
\endfirsthead
\toprule
Case & Contrast & Lower & Upper \\
\midrule
\endhead
\bottomrule
\endfoot
C2;2;1;U & L minus L+ & 0.011448 & 0.014567 \\
C2;2;1;U & L minus R & 0.023033 & 0.028988 \\
C2;2;1;U & L minus R+ & 0.024948 & 0.030894 \\
C2;2;1;U & L+ minus R & 0.010133 & 0.015872 \\
C2;2;1;U & L+ minus R+ & 0.012051 & 0.017775 \\
C2;2;1;U & R minus R+ & 0.00067113 & 0.0031495 \\
C3;2;1;U & L minus L+ & 0.0077535 & 0.010757 \\
C3;2;1;U & L minus R & 0.01796 & 0.023776 \\
C3;2;1;U & L minus R+ & 0.019309 & 0.025125 \\
C3;2;1;U & L+ minus R & 0.0087758 & 0.014449 \\
C3;2;1;U & L+ minus R+ & 0.010129 & 0.015794 \\
C3;2;1;U & R minus R+ & 0.0001224 & 0.0025755 \\
C4;2;1;U & L minus L+ & 0.016432 & 0.019542 \\
C4;2;1;U & L minus R & 0.033636 & 0.039473 \\
C4;2;1;U & L minus R+ & 0.033741 & 0.039578 \\
C4;2;1;U & L+ minus R & 0.015807 & 0.021327 \\
C4;2;1;U & L+ minus R+ & 0.015912 & 0.021433 \\
C4;2;1;U & R minus R+ & 0 & 0.0012997 \\
I2;2;1;U & L minus L+ & 0.023541 & 0.026849 \\
I2;2;1;U & L minus R & 0.042713 & 0.048868 \\
I2;2;1;U & L minus R+ & 0.043151 & 0.049305 \\
I2;2;1;U & L+ minus R & 0.017772 & 0.023418 \\
I2;2;1;U & L+ minus R+ & 0.018209 & 0.023856 \\
I2;2;1;U & R minus R+ & 0 & 0.00164 \\
O2;2;1;M & L minus L+ & 0.007622 & 0.010239 \\
O2;2;1;M & L minus R & 0.0064146 & 0.011415 \\
O2;2;1;M & L minus R+ & 0.00643 & 0.011431 \\
O2;2;1;M & L+ minus R & -0.0023997 & 0.002369 \\
O2;2;1;M & L+ minus R+ & -0.0023843 & 0.0023843 \\
O2;2;1;M & R minus R+ & 0 & 0.0012077 \\
O2;2;1;L & L minus L+ & 0.0099411 & 0.012722 \\
O2;2;1;L & L minus R & 0.054972 & 0.060123 \\
O2;2;1;L & L minus R+ & 0.055155 & 0.060307 \\
O2;2;1;L & L+ minus R & 0.043822 & 0.048609 \\
O2;2;1;L & L+ minus R+ & 0.044007 & 0.048792 \\
O2;2;1;L & R minus R+ & 0 & 0.001381 \\
O2;2;1;H & L minus L+ & 0.0014934 & 0.003946 \\
O2;2;1;H & L minus R & 0.0042173 & 0.0090667 \\
O2;2;1;H & L minus R+ & 0.0049502 & 0.0097948 \\
O2;2;1;H & L+ minus R & 0.001523 & 0.0063216 \\
O2;2;1;H & L+ minus R+ & 0.0022567 & 0.007049 \\
O2;2;1;H & R minus R+ & 0 & 0.0019272 \\
O2;2;1;U & L minus L+ & 0.0056474 & 0.0083165 \\
O2;2;1;U & L minus R & 0.013161 & 0.018416 \\
O2;2;1;U & L minus R+ & 0.013385 & 0.01864 \\
O2;2;1;U & L+ minus R & 0.0062318 & 0.011381 \\
O2;2;1;U & L+ minus R+ & 0.0064562 & 0.011605 \\
O2;2;1;U & R minus R+ & 0 & 0.0014206 \\
O2;2;4;M & L minus L+ & 0.0036933 & 0.0063836 \\
O2;2;4;M & L minus R & 0.035287 & 0.040496 \\
O2;2;4;M & L minus R+ & 0.035387 & 0.040597 \\
O2;2;4;M & L+ minus R & 0.03028 & 0.035426 \\
O2;2;4;M & L+ minus R+ & 0.030381 & 0.035526 \\
O2;2;4;M & R minus R+ & 0 & 0.0013773 \\
O2;2;4;L & L minus L+ & 0.0010828 & 0.0036974 \\
O2;2;4;L & L minus R & 0.020345 & 0.025506 \\
O2;2;4;L & L minus R+ & 0.020345 & 0.025506 \\
O2;2;4;L & L+ minus R & 0.017976 & 0.023095 \\
O2;2;4;L & L+ minus R+ & 0.017976 & 0.023095 \\
O2;2;4;L & R minus R+ & 0 & 0.0012764 \\
O2;2;4;H & L minus L+ & 0.0034416 & 0.006155 \\
O2;2;4;H & L minus R & 0.039236 & 0.044497 \\
O2;2;4;H & L minus R+ & 0.039288 & 0.044549 \\
O2;2;4;H & L+ minus R & 0.034497 & 0.039639 \\
O2;2;4;H & L+ minus R+ & 0.034549 & 0.039692 \\
O2;2;4;H & R minus R+ & 0 & 0.0013293 \\
O2;2;4;U & L minus L+ & 0.004188 & 0.0069254 \\
O2;2;4;U & L minus R & 0.011181 & 0.016658 \\
O2;2;4;U & L minus R+ & 0.011361 & 0.016838 \\
O2;2;4;U & L+ minus R & 0.0056495 & 0.011076 \\
O2;2;4;U & L+ minus R+ & 0.0058295 & 0.011256 \\
O2;2;4;U & R minus R+ & 0 & 0.0014585 \\
O2;3;1;M & L minus L+ & 0.0016956 & 0.0047777 \\
O2;3;1;M & L minus R & 0.015697 & 0.021868 \\
O2;3;1;M & L minus R+ & 0.015714 & 0.021886 \\
O2;3;1;M & L+ minus R & 0.012467 & 0.018624 \\
O2;3;1;M & L+ minus R+ & 0.012485 & 0.018641 \\
O2;3;1;M & R minus R+ & 0 & 0.0015083 \\
O2;3;1;L & L minus L+ & 0.0020031 & 0.005116 \\
O2;3;1;L & L minus R & 0.026367 & 0.032516 \\
O2;3;1;L & L minus R+ & 0.026477 & 0.032626 \\
O2;3;1;L & L+ minus R & 0.022819 & 0.028945 \\
O2;3;1;L & L+ minus R+ & 0.022928 & 0.029055 \\
O2;3;1;L & R minus R+ & 0 & 0.0016012 \\
O2;3;1;H & L minus L+ & 0.00030976 & 0.0033545 \\
O2;3;1;H & L minus R & 0.011221 & 0.017268 \\
O2;3;1;H & L minus R+ & 0.011378 & 0.017424 \\
O2;3;1;H & L+ minus R & 0.0093942 & 0.01543 \\
O2;3;1;H & L+ minus R+ & 0.0095515 & 0.015586 \\
O2;3;1;H & R minus R+ & 0 & 0.0016481 \\
O2;3;1;U & L minus L+ & 0.0014967 & 0.0045926 \\
O2;3;1;U & L minus R & 0.0095127 & 0.015773 \\
O2;3;1;U & L minus R+ & 0.0095603 & 0.01582 \\
O2;3;1;U & L+ minus R & 0.0064826 & 0.012713 \\
O2;3;1;U & L+ minus R+ & 0.0065302 & 0.012761 \\
O2;3;1;U & R minus R+ & 0 & 0.0015385 \\
O2;3;4;M & L minus L+ & 9.4671e-05 & 0.0033779 \\
O2;3;4;M & L minus R & 0.010433 & 0.016948 \\
O2;3;4;M & L minus R+ & 0.01066 & 0.017175 \\
O2;3;4;M & L+ minus R & 0.0087064 & 0.015202 \\
O2;3;4;M & L+ minus R+ & 0.0089337 & 0.015429 \\
O2;3;4;M & R minus R+ & 0 & 0.00184 \\
O2;3;4;L & L minus L+ & 0.0005048 & 0.0037993 \\
O2;3;4;L & L minus R & 0.0095446 & 0.01608 \\
O2;3;4;L & L minus R+ & 0.0095504 & 0.016086 \\
O2;3;4;L & L+ minus R & 0.0074107 & 0.01391 \\
O2;3;4;L & L+ minus R+ & 0.0074164 & 0.013916 \\
O2;3;4;L & R minus R+ & 0 & 0.0016183 \\
O2;3;4;H & L minus L+ & 0.0011154 & 0.0044219 \\
O2;3;4;H & L minus R & 0.0073415 & 0.013898 \\
O2;3;4;H & L minus R+ & 0.0073538 & 0.013911 \\
O2;3;4;H & L+ minus R & 0.004596 & 0.011107 \\
O2;3;4;H & L+ minus R+ & 0.0046083 & 0.011119 \\
O2;3;4;H & R minus R+ & 0 & 0.0016249 \\
O2;3;4;U & L minus L+ & 0.00021625 & 0.0035063 \\
O2;3;4;U & L minus R & 0.0056055 & 0.012273 \\
O2;3;4;U & L minus R+ & 0.0057651 & 0.012433 \\
O2;3;4;U & L+ minus R & 0.0037509 & 0.010406 \\
O2;3;4;U & L+ minus R+ & 0.0039105 & 0.010565 \\
O2;3;4;U & R minus R+ & 0 & 0.001774 \\
Q2;2;1;U & L minus L+ & 0.0055893 & 0.0083122 \\
Q2;2;1;U & L minus R & 0.01186 & 0.017187 \\
Q2;2;1;U & L minus R+ & 0.01259 & 0.017917 \\
Q2;2;1;U & L+ minus R & 0.0049684 & 0.010177 \\
Q2;2;1;U & L+ minus R+ & 0.0056983 & 0.010907 \\
Q2;2;1;U & R minus R+ & 0 & 0.0019382 \\
P2;2;1;U & L minus L+ & 0.011497 & 0.014615 \\
P2;2;1;U & L minus R & 0.022852 & 0.028801 \\
P2;2;1;U & L minus R+ & 0.024778 & 0.030719 \\
P2;2;1;U & L+ minus R & 0.0099052 & 0.015635 \\
P2;2;1;U & L+ minus R+ & 0.011835 & 0.01755 \\
P2;2;1;U & R minus R+ & 0.00068331 & 0.0031614 \\
Q2;2;4;U & L minus L+ & 0.003535 & 0.0062382 \\
Q2;2;4;U & L minus R & 0.0086864 & 0.014074 \\
Q2;2;4;U & L minus R+ & 0.0088869 & 0.014274 \\
Q2;2;4;U & L+ minus R & 0.0038216 & 0.0091655 \\
Q2;2;4;U & L+ minus R+ & 0.004022 & 0.0093656 \\
Q2;2;4;U & R minus R+ & 0 & 0.0014794 \\
Q2;3;1;U & L minus L+ & 0.0017403 & 0.0048791 \\
Q2;3;1;U & L minus R & 0.01136 & 0.017753 \\
Q2;3;1;U & L minus R+ & 0.011488 & 0.01788 \\
Q2;3;1;U & L+ minus R & 0.0080678 & 0.014426 \\
Q2;3;1;U & L+ minus R+ & 0.0081953 & 0.014553 \\
Q2;3;1;U & R minus R+ & 0 & 0.0016207 \\
Q2;3;4;U & L minus L+ & 0.0020091 & 0.0053568 \\
Q2;3;4;U & L minus R & 0.010559 & 0.017349 \\
Q2;3;4;U & L minus R+ & 0.011154 & 0.017943 \\
Q2;3;4;U & L+ minus R & 0.006888 & 0.013654 \\
Q2;3;4;U & L+ minus R+ & 0.0074835 & 0.014248 \\
Q2;3;4;U & R minus R+ & 0 & 0.0022148 \\
P3;2;1;U & L minus L+ & 0.0078066 & 0.010814 \\
P3;2;1;U & L minus R & 0.017874 & 0.023687 \\
P3;2;1;U & L minus R+ & 0.019232 & 0.025044 \\
P3;2;1;U & L+ minus R & 0.008637 & 0.014304 \\
P3;2;1;U & L+ minus R+ & 0.0099988 & 0.015657 \\
P3;2;1;U & R minus R+ & 0.00013073 & 0.0025845 \\
P4;2;1;U & L minus L+ & 0.016374 & 0.019482 \\
P4;2;1;U & L minus R & 0.033196 & 0.039029 \\
P4;2;1;U & L minus R+ & 0.033446 & 0.039279 \\
P4;2;1;U & L+ minus R & 0.015426 & 0.020943 \\
P4;2;1;U & L+ minus R+ & 0.015676 & 0.021193 \\
P4;2;1;U & R minus R+ & 0 & 0.0014477 \\
A2;2;1;U & L minus L+ & 0.00062241 & 0.0030972 \\
A2;2;1;U & L minus R & -0.0024517 & 0.0023493 \\
A2;2;1;U & L minus R+ & -0.00057956 & 0.0042851 \\
A2;2;1;U & L+ minus R & -0.0043435 & 0.00052149 \\
A2;2;1;U & L+ minus R+ & -0.0024125 & 0.0023985 \\
A2;2;1;U & R minus R+ & 0.00066516 & 0.0031428 \\
\end{longtable}
}
\noindent{\footnotesize Both tables are reconstructed from stored outward means, variance bounds and deterministic supports. Numerical replay uses the same original samples. Statistical interpretation requires the independent-bin model.\par}

{\small
\setlength{\tabcolsep}{3pt}
\begin{longtable}{@{}llrrrrr@{}}
\caption{Final-action mechanism diagnostics}\label{tab:mechanism50}\\
\toprule
Case & Policy & Accepted & Blocked & Old action & New action & Derivative calls \\
\midrule
\endfirsthead
\toprule
Case & Policy & Accepted & Blocked & Old action & New action & Derivative calls \\
\midrule
\endhead
\bottomrule
\endfoot
C2;2;1;U & L & 81251 & 1686 & 0.13805 & 0.16199 & 1216761 \\
C2;2;1;U & R & 92005 & 2261 & 0.17813 & 0.16288 & 1223163 \\
C3;2;1;U & L & 71835 & 1430 & 0.15264 & 0.16793 & 1212973 \\
C3;2;1;U & R & 86739 & 1858 & 0.18142 & 0.16825 & 1220030 \\
C4;2;1;U & L & 105827 & 1124 & 0.11931 & 0.16893 & 1209765 \\
C4;2;1;U & R & 79985 & 9020 & 0.16899 & 0.16935 & 1219031 \\
I2;2;1;U & L & 115800 & 1824 & 0.10471 & 0.16094 & 1216727 \\
I2;2;1;U & R & 87457 & 6803 & 0.16155 & 0.16239 & 1222172 \\
O2;2;1;M & L & 123097 & 0 & 0.13434 & 0.18236 & 1179648 \\
O2;2;1;M & R & 66240 & 0 & 0.18188 & 0.18236 & 1179648 \\
O2;2;1;L & L & 51516 & 0 & 0.13357 & 0.14812 & 1179648 \\
O2;2;1;L & R & 25162 & 0 & 0.14947 & 0.14977 & 1179648 \\
O2;2;1;H & L & 131072 & 0 & 0.11338 & 0.11033 & 1301723 \\
O2;2;1;H & R & 131072 & 0 & 0.10626 & 0.12107 & 1286030 \\
O2;2;1;U & L & 96199 & 2180 & 0.13976 & 0.16238 & 1217297 \\
O2;2;1;U & R & 86179 & 5606 & 0.16295 & 0.16247 & 1223021 \\
O2;2;4;M & L & 131072 & 0 & 0.06499 & 0.09143 & 1253562 \\
O2;2;4;M & R & 131072 & 0 & 0.08886 & 0.08423 & 1245289 \\
O2;2;4;L & L & 131072 & 0 & 0.12539 & 0.13479 & 1227082 \\
O2;2;4;L & R & 0 & 131072 & 0.13106 & 0.13106 & 1232486 \\
O2;2;4;H & L & 131072 & 0 & 0.04179 & 0.04061 & 1294458 \\
O2;2;4;H & R & 131072 & 0 & 0.05060 & 0.04725 & 1278365 \\
O2;2;4;U & L & 127091 & 2453 & 0.08431 & 0.08510 & 1255705 \\
O2;2;4;U & R & 116259 & 12245 & 0.08476 & 0.08500 & 1255578 \\
O2;3;1;M & L & 74044 & 0 & 0.15964 & 0.17640 & 1179648 \\
O2;3;1;M & R & 50391 & 0 & 0.17885 & 0.17907 & 1179648 \\
O2;3;1;L & L & 47744 & 0 & 0.14872 & 0.15570 & 1179648 \\
O2;3;1;L & R & 68952 & 0 & 0.15623 & 0.15647 & 1179648 \\
O2;3;1;H & L & 127923 & 3149 & 0.15631 & 0.16935 & 1267556 \\
O2;3;1;H & R & 130156 & 880 & 0.16622 & 0.15981 & 1272977 \\
O2;3;1;U & L & 93095 & 1636 & 0.15635 & 0.17217 & 1200960 \\
O2;3;1;U & R & 68286 & 6835 & 0.17297 & 0.17285 & 1207373 \\
O2;3;4;M & L & 131009 & 63 & 0.08959 & 0.08991 & 1253921 \\
O2;3;4;M & R & 131072 & 0 & 0.08589 & 0.09307 & 1249060 \\
O2;3;4;L & L & 129450 & 1622 & 0.11835 & 0.11804 & 1234478 \\
O2;3;4;L & R & 75986 & 44091 & 0.11664 & 0.11708 & 1237417 \\
O2;3;4;H & L & 130535 & 487 & 0.06486 & 0.07197 & 1266116 \\
O2;3;4;H & R & 81893 & 46520 & 0.06993 & 0.07110 & 1270467 \\
O2;3;4;U & L & 127548 & 3104 & 0.09220 & 0.09283 & 1251214 \\
O2;3;4;U & R & 119926 & 9604 & 0.09243 & 0.09264 & 1251454 \\
Q2;2;1;U & L & 84626 & 1862 & 0.14266 & 0.16266 & 1217576 \\
Q2;2;1;U & R & 89341 & 2780 & 0.16443 & 0.16248 & 1222725 \\
P2;2;1;U & L & 81414 & 1734 & 0.13768 & 0.16190 & 1217185 \\
P2;2;1;U & R & 92131 & 2223 & 0.17824 & 0.16278 & 1223408 \\
Q2;2;4;U & L & 127314 & 3044 & 0.08473 & 0.08526 & 1255543 \\
Q2;2;4;U & R & 116513 & 12198 & 0.08462 & 0.08509 & 1255620 \\
Q2;3;1;U & L & 78872 & 1217 & 0.15697 & 0.17211 & 1200463 \\
Q2;3;1;U & R & 72138 & 3059 & 0.17400 & 0.17287 & 1207783 \\
Q2;3;4;U & L & 127944 & 2619 & 0.09086 & 0.09284 & 1251053 \\
Q2;3;4;U & R & 124704 & 5413 & 0.09262 & 0.09261 & 1251454 \\
P3;2;1;U & L & 71967 & 1425 & 0.15260 & 0.16790 & 1213011 \\
P3;2;1;U & R & 86489 & 1881 & 0.18143 & 0.16824 & 1220010 \\
P4;2;1;U & L & 106376 & 1130 & 0.11948 & 0.16893 & 1210149 \\
P4;2;1;U & R & 86994 & 7077 & 0.16745 & 0.16930 & 1218786 \\
A2;2;1;U & L & 90517 & 2162 & 0.17906 & 0.16290 & 1223395 \\
A2;2;1;U & R & 92064 & 2277 & 0.17821 & 0.16286 & 1223276 \\
\end{longtable}
}
\noindent{\footnotesize Means are descriptive over the recorded paths, not confidence intervals for separate mechanism parameters. Complete date-wise earlier actions, terminal state vectors, capacity and quantum repair counts, and exact fallbacks are deposited.\par}

{\small
\setlength{\tabcolsep}{3pt}
\begin{longtable}{@{}llrrrrr@{}}
\caption{New-construction operation and storage account}\label{tab:work50}\\
\toprule
Case & Method & Target eval. & Actor scalars & KiB & Coeff. bits & MiB \\
\midrule
\endfirsthead
\toprule
Case & Method & Target eval. & Actor scalars & KiB & Coeff. bits & MiB \\
\midrule
\endhead
\bottomrule
\endfoot
C2;2;1;U & W & 1,156 & 578 & 41.8 & 36 & 41.4 \\
C2;2;1;U & F & 1,156 & 578 & 26.1 & 42 & 41.7 \\
C3;2;1;U & W & 19,652 & 9,826 & 851.1 & 42 & 44.3 \\
C3;2;1;U & F & 19,652 & 9,826 & 434.4 & 42 & 42.3 \\
C4;2;1;U & W & 236,196 & 13,122 & 1024.5 & 42 & 50.5 \\
C4;2;1;U & F & 3,242,952 & 167,042 & 7464.3 & 42 & 169.8 \\
I2;2;1;U & W & 11,664 & 162 & 14.2 & 41 & 42.4 \\
I2;2;1;U & F & 11,664 & 162 & 9.3 & 41 & 42.0 \\
Q2;2;1;U & W & 13,068 & 2,178 & 171.9 & 41 & 41.6 \\
Q2;2;1;U & F & 13,068 & 2,178 & 98.6 & 42 & 42.1 \\
P2;2;1;U & W & 1,156 & 578 & 41.8 & 36 & 42.0 \\
P2;2;1;U & F & 1,156 & 578 & 26.1 & 42 & 42.0 \\
Q2;2;4;U & W & 19,602 & 2,178 & 178.4 & 42 & 41.6 \\
Q2;2;4;U & F & 19,602 & 2,178 & 99.4 & 42 & 42.0 \\
Q2;3;1;U & W & 126,750 & 12,675 & 1005.0 & 42 & 45.8 \\
Q2;3;1;U & F & 126,750 & 12,675 & 550.0 & 42 & 41.9 \\
Q2;3;4;U & W & 126,750 & 12,675 & 983.8 & 42 & 45.8 \\
Q2;3;4;U & F & 126,750 & 12,675 & 539.4 & 42 & 41.8 \\
P3;2;1;U & W & 19,652 & 9,826 & 851.1 & 42 & 44.5 \\
P3;2;1;U & F & 19,652 & 9,826 & 434.4 & 42 & 42.5 \\
P4;2;1;U & W & 236,196 & 13,122 & 1024.5 & 42 & 51.3 \\
P4;2;1;U & F & 236,196 & 13,122 & 575.2 & 42 & 48.3 \\
A2;2;1;U & A & 4,472 & 578 & 39.5 & 42 & 41.9 \\
A2;2;1;U & F & 1,156 & 578 & 26.1 & 42 & 42.0 \\
\end{longtable}
}
\noindent{\footnotesize Target evaluations include every failed common-accuracy attempt. Actor scalars and checkpoint KiB describe the returned model. Coefficient bits are the true maximum, not the additive descriptor in the raw prefix map; MiB is process peak memory. Full operation categories, including terminal gates and exact fallbacks, remain in JSON.\par}


\section{Numerical correction and sequential common-accuracy design}\label{supp:corrections50}
The primary shared allocation is sufficient for the witness bound, not automatically for conventional FVI's actual interpolated modulus. Its failed certificates are retained. The common-accuracy amendment was specified after that outcome and frozen before its own execution; it charges every failed construction before doubling the state quota. It is not a retrospective recoding of the primary study.

During publication review, nearest-rounded subtraction of different policies' already outward path endpoints was replaced by directed outward subtraction. The same bin streams, policies and sample counts were replayed. Original sources and results remain unchanged; a separate correction manifest binds the new evaluator, and the audit checks identical streams and unchanged sign classifications. Replaying the same observations adds no statistical sample size. The additive raw prefix field for coefficient-bit descriptors is not interpreted as a maximum; maximum bit length and live storage are reconstructed separately. The correction and development logs are part of the release.


\section{Exhaustive robustness-service accounting}\label{supp:work52}
\begin{table}[htbp]
\caption{Complete robustness-service work and memory}
\label{tab:work52}
\centering\small
\begin{tabular}{lrrrr}
\toprule
Incumbent & $T$ & $p$ & Seconds & Peak MiB \\
\midrule
Witness & 2 & 1 & 3.73 & 50.6 \\
Witness & 2 & 4 & 3.72 & 50.6 \\
Witness & 3 & 1 & 4.37 & 56.5 \\
Witness & 3 & 4 & 4.32 & 56.5 \\
FVI & 2 & 1 & 3.63 & 50.7 \\
FVI & 2 & 4 & 3.76 & 51.0 \\
FVI & 3 & 1 & 3.91 & 52.5 \\
FVI & 3 & 4 & 3.86 & 52.3 \\
\bottomrule
\end{tabular}
\parbox{0.96\textwidth}{\footnotesize Each isolated service loads a frozen incumbent, computes all proposals and centered intervals, independently replays the inherited terminal gate, evaluates every amplitude gate, and durably serializes the cell records. Process startup is excluded; inherited construction time is not added. These are verification-service clocks, not a work-to-policy-cost superiority claim.}
\end{table}

\bibliographystyle{ecta-fullname}
\bibliography{references}
\end{document}
